\documentclass[JEP,XML,SOM,Unicode,published]{cedram}
\datereceived{2021-03-11}
\dateaccepted{2021-10-25}
\dateepreuves{2021-11-12}

\usepackage[scr=rsfs,cal=euler]{mathalfa}
\multlinegap0pt
\newenvironment{enumeratei}
{\bgroup\def\theenumi{\roman{enumi}}\def\theenumii{\arabic{enumii}}\begin{enumerate}}
{\end{enumerate}\egroup}
\newcommand{\RedefinitSymbole}[1]{%
\expandafter\let\csname old\string#1\endcsname=#1
\let#1=\relax
\newcommand{#1}{\csname old\string#1\endcsname\,}%
}
\RedefinitSymbole{\forall} \RedefinitSymbole{\exists}
\newcommand\mto{\mathchoice{\longmapsto}{\mapsto}{\mapsto}{\mapsto}}
\def\To#1{\mathchoice{\xrightarrow{\textstyle\kern4pt#1\kern3pt}}{\stackrel{#1}{\longrightarrow}}{}{}}
\newcommand{\psfrac}[2]{\sfrac{(#1)}{#2}}
\newcommand{\spfrac}[2]{\sfrac{#1}{(#2)}}
\newcommand{\Psfrac}[2]{(\sfrac{#1}{#2})}

\DeclareMathOperator{\Ker}{Ker}
\newcommand{\rmc}{\mathrm{c}}
\newcommand{\GH}{\mathrm{GH}}
\newcommand{\sGH}{{\scriptscriptstyle\mathrm{GH}\,}}
\let\sep:
\let\backslash\setminus

\newcommand{\Changel}{\mathcode`l="7160}
\newcommand{\Changelback}{\mathcode`l="716C}
\newcommand{\NoChangel}[1]{%
\expandafter\let\csname old\string#1\endcsname=#1
\let#1=\relax
\newcommand{#1}{\mathcode`l="716C\csname old\string#1\endcsname\mathcode`l="7160 }%
}

\NoChangel{\log}
\NoChangel{\ln}
\NoChangel{\lim}
\NoChangel{\limsup}
\NoChangel{\liminf}
\NoChangel{\varlimsup}
\NoChangel{\varliminf}
\NoChangel{\varinjlim}
\NoChangel{\varprojlim}

\usepackage{microtype}
\usepackage{upgreek}
\usepackage{mathtools}
\usepackage{mathrsfs}
\usepackage{aliascnt}
\usepackage{esint}

\mathtoolsset{showonlyrefs}

\theoremstyle{plain}
\newtheorem{theorem}{Theorem}[section]
\newtheorem{proposition}[theorem]{Proposition}
\newtheorem{lemma}[theorem]{Lemma}
\newtheorem{claim}[theorem]{Claim}
\newtheorem{corollary}[theorem]{Corollary}

\theoremstyle{definition}
\newtheorem{definition}[theorem]{Definition}
\newtheorem{remark}[theorem]{Remark}

\newcommand{\setC}{\mathbb{C}}
\newcommand{\setN}{\mathbb{N}}
\newcommand{\setQ}{\mathbb{Q}}
\newcommand{\setR}{\mathbb{R}}
\newcommand{\setX}{\mathbb{X}}
\newcommand{\bP}{\mathbb{P}}
\newcommand{\setB}{\mathbb{B}}
\newcommand{\cE}{\mathcal{E}}
\newcommand{\cV}{\mathcal{V}}
\newcommand{\ucV}{\underline{\mathcal{V}}}
\newcommand{\cW}{\underline{\mathcal{W}}}
\newcommand{\cD}{\mathcal{D}}

\let\eps\upepsilon
\let\phi\varphi

\DeclareMathOperator{\Id}{Id}
\DeclareMathOperator{\tr}{tr}
\let\div\undefined
\DeclareMathOperator{\div}{div}
\DeclareMathOperator{\Rad}{Rad}
\DeclareMathOperator{\sign}{sign}
\newcommand{\di}{\mathop{}\!\mathrm{d}}
\newcommand{\bs}{\mathrm{bs}}
\newcommand{\loc}{\mathrm{loc}}\NoChangel{\loc}
\newcommand{\ev}{\mathrm{ev}}
\DeclareMathOperator{\supp}{supp}
\newcommand{\rmD}{\mathrm{D}}
\newcommand{\scal}[2]{\ensuremath{\langle #1, #2 \rangle}}

\DeclareMathOperator{\Lip}{Lip}
\newcommand{\leb}{\mathscr{L}}
\newcommand{\measrestr}{%
\,\raisebox{-.127ex}{\reflectbox{\rotatebox[origin=br]{-90}{$\lnot$}}}\,%
}

\newcommand{\dist}{\mathsf{d}}
\DeclareMathOperator{\Span}{Span}
\DeclareMathOperator{\vol}{vol}\NoChangel{\vol}
\DeclareMathOperator{\Ric}{Ric}

\newcommand{\uX}{\underline{X}}
\newcommand{\uB}{\underline{B}}
\newcommand{\ub}{\underline{b}}
\newcommand{\usigma}{\underline{\sigma}}
\DeclareMathOperator{\udist}{\underline{\dist}}
\newcommand{\uo}{\underline{o}}
\newcommand{\ux}{\underline{x}}
\newcommand{\uz}{\underline{z}}
\newcommand{\uc}{\underline{c}}
\newcommand{\uL}{\underline{L}}
\newcommand{\uphi}{\underline{\phi}}
\newcommand{\uD}{\underline{D}}
\newcommand{\uH}{\underline{H}}
\newcommand{\uS}{\underline{S}}
\newcommand{\umu}{\underline{\mu}}
\newcommand{\unu}{\underline{\nu}}
\newcommand{\ubeta}{\underline{\beta}}
\newcommand{\ugamma}{\underline{\gamma}}
\newcommand{\uy}{\underline{y}}
\newcommand{\upp}{\underline{p}}
\newcommand{\uh}{\underline{h}}

\DeclareMathOperator*{\essssup}{ess\,sup}
\DeclareMathOperator*{\esssinf}{ess\,inf}

\datepublished{2021-11-17}
\begin{document}
\frontmatter

\title{A rigidity result for metric measure spaces with Euclidean heat kernel}

\author[\initial{G.} \lastname{Carron}]{\firstname{Gilles} \lastname{Carron}}
\address{Université de Nantes, Département de Mathématiques\\
2 rue de la Houssinière, BP 92208, 44322 Nantes cedex 03, France}
\email{Gilles.Carron@univ-nantes.fr}
\urladdr{https://www.math.sciences.univ-nantes.fr/~carron/}

\thanks{G.C. was partially supported by the ANR grants ANR-17-CE40-0034 (CCEM) and ANR-18-CE40-0012(RAGE).}

\author[\initial{D.} \lastname{Tewodrose}]{\firstname{David} \lastname{Tewodrose}}
\address{CY Cergy Paris University, AGM\\
2 av. Adolphe Chauvin, 95302 Cergy-Pontoise cedex, France}
\email{David.Tewodrose@univ-nantes.fr}
\urladdr{http://www.davidtewodrose.com/}

\begin{abstract}
We prove that a metric measure space equipped with a Dirichlet form admitting an Euclidean heat kernel is necessarily isometric to the Euclidean space. This helps us providing an alternative proof of Colding's celebrated almost rigidity volume theorem via a quantitative version of our main result. We also discuss the case of a metric measure space equipped with a Dirichlet form admitting a spherical heat kernel.
\end{abstract}

\subjclass{35K08, 31C25, 53C23, 53C21}

\keywords{Heat kernel, harmonic functions, asymptotic cone}

\altkeywords{Noyau de la chaleur, fonctions harmoniques, cône asymptotique}

\alttitle{Un résultat de rigidité pour les espaces métriques mesurés à noyau de la chaleur euclidien}

\begin{altabstract}
Nous prouvons qu'un espace métrique mesuré équipé d'une forme de Dirichlet admettant un noyau de la chaleur euclidien est nécessairement isométrique à l'espace euclidien. Nous en déduisons une preuve alternative du célèbre théorème de presque rigidité du volume de Colding grâce à une version quantitative de notre résultat principal. Nous traitons aussi le cas d'un espace métrique mesuré équipé d'une forme de Dirichlet admettant un noyau de la chaleur sphérique.
\end{altabstract}

\maketitle
\tableofcontents
\mainmatter

\section{Introduction}

In $\setR^n$, the classical Dirichlet energy is the functional defined on $H^1$ by
\[
E(u):= \int_{\setR^n} |\nabla u|^2
\]
for any $u \in H^1$. As well-known, it is related to the Laplace operator $\Delta:=\sum_{k=1}^{n}\partial_{kk}$ by the integration by parts formula, namely
\[
E(u,v)=-\int_{\setR^n} (\Delta u) v
\]
for any $u,v \in H^1$ such that $\nabla u\in H^1$, where $E(u,v):=\int_{\setR^n} \langle \nabla u, \nabla v\rangle$. Standard tools from spectral theory show that $\Delta$ generates a semi-group of operators $(e^{t\Delta})_{t>0}$ sending any $u_0 \in L^2$ to the family $(u_t)_{t>0}\subset H^1$ satisfying the heat equation $\partial_t u_t = \Delta u_t$ with $u_0$ as an initial condition. The semi-group $(e^{t\Delta})_{t>0}$ admits a smooth kernel $p$, so that for any $f \in L^2$, $x \in \setR^n$ and $t>0$,
\[
e^{t\Delta} f(x) = \int_{\setR^n} p(x,y,t)f(y) \di y.
\]
The explicit expression of this heat kernel is well-known: for any $x,y \in \setR^n$ and $t>0$,
\[
p(x,y,t) = \frac{1}{(4 \pi t)^{n/2}}\, e^{-\sfrac{|x-y|^2}{4t}}.
\]

In the more general context of a measured space $(X,\mu)$, the Dirichlet energy possesses abstract analogues called Dirichlet forms. Associated with any such a form~$\cE$ is a self-adjoint operator $L$ whose properties are similar to the Laplace operator; in particular, the spectral theorem applies to it and provides a semi-group $(P_t)_{t>0}$ delivering the solution of the equation $\partial_t u_t = L u_t$ starting from any square integrable initial condition. Under suitable assumptions, this semi-group admits a kernel. When the space $X$ is equipped with a metric $\dist$ generating the $\sigma$-algebra on which $\mu$ is defined, this kernel is often compared with the exact Gaussian kernel
\[
\frac{1}{(4 \pi t)^{n/2}}\, e^{-\sfrac{\dist^2(x,y)}{4t}}
\]
through upper and lower estimates: see \cite{Sturm2}, for instance. From this perspective, a natural question arises: what happens when the kernel of $\cE$ coincides with this Gaussian term? In this article, we answer this question by showing that the unique metric measure space admitting such a kernel is the Euclidean space. The precise statement of our main result is the following:

\begin{theorem}\label{th:main}
Let $(X,\dist)$ be a complete metric space equipped with a non-negative regular Borel measure $\mu$. Assume that there exists a symmetric Dirichlet form $\cE$ on $(X,\mu)$ admitting a heat kernel $p$ such that for some $\alpha>0$,
\begin{equation}\label{eq:heatkernel}
p(x,y,t) = \frac{1}{(4 \pi t)^{\alpha/2}}\, e^{-\sfrac{\dist^2(x,y)}{4t}}
\end{equation}
holds for any $x,y \in X$ and any $t>0$. Then $\alpha$ is an integer, $(X,\dist)$ is isometric to $(\setR^\alpha,\dist_e)$, where $\dist_e$ stands for the classical Euclidean distance, and $\mu$ is the $\alpha$-dimensional Hausdorff measure.
\end{theorem}

Then we show that this rigidity result can be turned quantitative via a suitable contradiction argument. Denoting by $\dist_{\sGH}$ the Gromov-Hausdorff distance and by $\setB^n_r$ any Euclidean ball in $\setR^n$ with radius $r>0$, we obtain the following:

\begin{theorem}\label{th:almostrigidity} Let $n$ be a positive integer. For any $\upepsilon>0$, there exists $\updelta>0$ depending only on $\upepsilon$ and $n$ such that if $(X,\dist,\mu)$ is a complete metric measure space endowed with a symmetric Dirichlet form $\cE$
admitting a heat kernel $p$ satisfying
\begin{equation}\label{eq:almostheatkernel}
(1-\updelta)\, \frac{1}{(4 \pi t)^{n/2}}\, e^{-\sfrac{\dist^2(x,y)}{4(1-\updelta)t}}\le p(x,y,t) \le (1+\updelta)\, \frac{1}{(4 \pi t)^{n/2}}\, e^{-\sfrac{\dist^2(x,y)}{4(1+\updelta)t}}
\end{equation}
for any $x,y \in X$ and $t \in (0,T]$, for some given $T>0$, then for any $x \in X$ and $r \in (0,\sqrt{T})$,
\begin{equation}%\label{eq:Reifenberg}
\dist_{\sGH}\left(B_r(x), \setB^n_r\right)< \upepsilon r.
\end{equation}
\end{theorem}

The intrinsic Reifenberg theorem of Cheeger and Colding \cite[Th.\,A.1.1.]{CheegerColding} provides the following immediate topological consequence, where $\Psi(\cdot|n)$ is a function depending only on $n$ with $\Psi(r|n) \to 0$ when $r \to 0^+$.

\begin{corollary}\label{cor:topo} Let $n$ be a positive integer. There exists $\updelta_n>0$ depending only on~$n$ such that if $(X,\dist,\mu)$ is a complete metric measure space endowed with a symmetric Dirichlet form $\cE$ admitting a heat kernel $p$ such that for some numbers $\delta \in (0,\delta_n)$ and $T>0$,
\[
(1-\updelta)\,\frac{1}{(4 \pi t)^{n/2}}\, e^{-\sfrac{\dist^2(x,y)}{4(1-\updelta) t}}\le p(x,y,t) \le (1+\updelta)\,\frac{1}{(4 \pi t)^{n/2}}\, e^{-\sfrac{\dist^2(x,y)}{4(1+\updelta) t}}
\]holds for all $x, y \in X$ and $t \in (0,T)$, then for any $x \in X$, there exists a topological embedding of $\mathbb{B}_{\sqrt{T}}^{n}$ into $B_{\sqrt{T}}(x)$ whose image contains $B_{(1-\Psi(\updelta|n))\sqrt{T}}(x)$.\end{corollary}

We point out the two previous results are also true in case $T=+\infty$. Moreover, Theorem \ref{th:almostrigidity} can be used to give an alternative proof of a celebrated result established by Colding \cite[Th.\,0.8]{Colding}, namely the almost rigidity of the volume for Riemannian manifolds with non-negative Ricci curvature. Let us recall this statement:

\begin{theorem}[Colding]\label{th:Colding} Let $n$ be a positive integer. For any $\upepsilon>0$, there exists $\updelta>0$ depending only on $\upepsilon$ and $n$ such that if
$(M^n,g)$ is a complete Riemannian manifold with non-negative Ricci curvature such that for any $x \in M$ and $r>0$,
\begin{equation}\label{eq:vol}\vol B_r(x) \ge \left(1-\updelta\right) \omega_n\, r^n,\end{equation}
then for
any $x\in M$ and $r>0$,
\[
\dist_{\sGH}\left(B_r(x), \setB^n_r\right)\le \upepsilon r.
\]
\end{theorem}

This theorem is a direct consequence of our almost rigidity theorem coupled with an intermediary result, Theorem \ref{th:estimate}, which states, roughly speaking, that a complete Riemannian manifold satisfying the volume estimate \eqref{eq:vol} necessarily has an almost Euclidean heat kernel. Our proof of this result is based on previous works by Cheeger and Yau \cite{CheegerYau}, Li and Yau \cite{LiYau} and especially Li, Tam and Wang \cite{LiTamWang}.

Finally, in the last section of this paper, we investigate the case of a metric measure space equipped with a spherical heat kernel. To be precise, the sphere $\mathbb{S}^n$ has a heat kernel which can be written as
\[
K_t^{(n)}(\dist_{\mathbb{S}^n}(x,y)),
\]
where $K_t^{(n)}$ is an explicit function and $\dist_{\mathbb{S}^n}$ is the classical round Riemannian distance. We show that if a metric measure space $(X,\dist,\mu)$ is equipped with a Dirichlet form admitting a heat kernel $p$ such that
\[
p(x,y,t)=K_t^{(n)}(\dist(x,y))
\]
for all $x, y \in X$ and $t>0$, then $(X,\dist)$ is isometric to $(\mathbb{S}^n,\dist_{\mathbb{S}^n})$.

Let us spend some words to describe our proof of Theorem \ref{th:main}. A key point is the celebrated result of Colding and Minicozzi asserting that on any complete Riemannian manifold satisfying the doubling and Poincaré properties, the space of harmonic maps with linear growth is finite-dimensional \cite{ColdingMinicozzi}. As already observed in non-smooth contexts \cite{Hua, HuaKellXia}, the proof of this result can be carried out on any complete metric measure spaces satisfying the doubling and Poincaré properties. It turns out that admitting a Dirichlet form with an Euclidean heat kernel forces the metric measure space to satisfy these two properties, see Proposition \ref{prop:important}.

Then we consider the functions
\[
B(x,\cdot):=\frac{1}{2}(\dist^2(o,x)+\dist^2(o,\cdot)-\dist^2(x,\cdot)), \qquad x \in X,
\]
which are easily shown to have linear growth. When $(X,\dist,\mu)$ is equipped with a Dirichlet form $\cE$ satisfying the assumptions of Theorem \ref{th:main}, these functions are locally $L$-harmonic: this follows from establishing
\[
L\textbf{1}=0 \qquad \text{and} \qquad L\dist^2(x,\cdot) = 2 \alpha.
\]
Therefore, the vector space $\cV$ generated by the functions $B(x,\cdot)$ has a finite dimension~$n$. Choosing a suitable basis $(h_1,\ldots,h_n)$ of this space, we embed $X$ into $\setR^n$ by setting
\[
H(x)=(h_1(x),\ldots,h_n(x))
\]
for any $x \in X$. More precisely, there exists $x_1,\ldots,x_n \in X$ such that $(\delta_{x_1},\ldots,\delta_{x_n})$ is a basis of $\cV^{*}$, where $\delta_{x}(h):=h(x)$ for any $x \in X$ and any $h \in \cV$, and $(h_1,\ldots,h_n)$ is chosen as the dual of this basis. Setting $Q(\xi):= \sum_{i,j} B(x_i,x_j)\xi_i \xi_j$ for any $\xi=(\xi_1,\dots,\xi_n) \in \setR^n$, we easily get
\begin{equation}\label{eq:isometry1}
Q(H(x)-H(y))=\dist^2(x,y)
\end{equation}
for any $x, y \in X$, thus $H$ is an embedding.

To conclude, we establish $\alpha = n$ and show that $Q$ is non-degenerate, so that $\dist_Q(\xi,\xi')=\sqrt{Q(\xi-\xi')}$ defines a distance on $\setR^n$ that is isometric to the Euclidean distance: then \eqref{eq:isometry1} shows that $H$ is an isometric embedding onto its image which, by a final argument, is shown to be $\setR^n$. We prove these two concluding assertions by the study of asymptotic cones at infinity of $(X,\dist,\mu)$.

Concerning the rigidity result for the spherical case, we instead embed $X$ into $E_1:=\Ker(L-\lambda_1 I)$, where $\lambda_1$ is the first non-zero eigenvalue of $L$, and show that $H(X)$ is isometric to $\Sigma:=\{Q=1\}$ for some suitable quadratic form $Q$.

The paper is organized as follows. Our proof of Theorem \ref{th:main} relies on several notions and results from different areas that we collect in the preliminary Section \ref{sec:2}. Then in Section \ref{sec:3} we establish simple rigidity results for metric measure spaces with an Euclidean heat kernel. We use these results in Section \ref{sec:4} which is dedicated to the proof of Theorem \ref{th:main}. Section \ref{sec:5} is devoted to the almost rigidity result, namely Theorem~\ref{th:almostrigidity}, and Section \ref{sec:6} explains our new proof of Colding's volume almost rigidity theorem. Finally Section \ref{sec:7} contains our study of the case of metric measure spaces equipped with a spherical heat kernel.

\subsubsection*{Acknowledgements}
The first author thanks the Centre Henri Lebesgue ANR-11-LABX-0020-01 for creating an attractive mathematical environment. The second author thanks S.~Honda for interesting remarks and questions at a late stage of this work and for the good working conditions in Tohoku University that grandly helped completing this article. Both authors thank the anonymous referees for their precious suggestions, notably to simplify the proof of Lemma 3.3.

\section{Preliminaries}\label{sec:2}

Throughout the article, we shall call metric measure space any triple $(X,\dist,\mu)$ where $(X,\dist)$ is a $\sigma$-compact metric space and $\mu$ is a non-negative $\sigma$-finite Radon measure on $(X,\dist)$ such that $\supp \mu = X$. Here $\supp \mu$ denotes the support of $\mu$. We shall keep fixed a number $\alpha>0$ and denote by $\omega_\alpha$ the quantity
\begin{equation}\label{eq:omega_n}
\omega_\alpha = \frac{\pi^{\alpha/2}}{\Gamma(\alpha/2+1)},
\end{equation}
where $\Gamma$ denotes the usual Gamma function $\{\mathrm{Re}>0\} \ni z \mto \int_0^{+\infty} t^{z-1}e^{-t} \di t$. Note that $\omega_n$ is the volume of the unit sphere in $\setR^n$ when $\alpha=n$ is an integer.

We shall use classical notations for the functional spaces defined on $(X,\dist, \mu)$, like $C(X)$ (\resp $C_\rmc(X)$) for the space of continuous (\resp compactly supported continuous) functions, $\Lip(X)$ (\resp $\Lip_\rmc(X)$)) for the space of Lipschitz (\resp compactly supported Lipschitz) functions, $L^p(X,\mu)$, where $p \in [1,+\infty)$, for the space of (equivalent classes of) $\mu$-measurable functions whose $p$-th power is $\mu$-integrable, $L^\infty(X,\mu)$ for the space of $\mu$-essentially bounded functions, and so on.
We shall write $\supp f$ for the support of a function $f$ and $1_A$ for the characteristic function of a set $A \subset X$.

A generic open ball in $(X,\dist)$ will be denoted by $B$, and we will write $\lambda B$ for the ball with same center as $B$ but radius multiplied by $\lambda>0$.

We will extensively make use of the following definition.

\begin{definition}
We say that a metric measure space $(X,\dist,\mu)$ has an $\alpha$-dimensional volume whenever $\mu(B)=\omega_\alpha r^\alpha$ for any metric ball $B\subset X$ with radius $r>0$.
\end{definition}

\subsection*{Dirichlet forms}

Let us recall some basic facts about Dirichlet forms, referring to \eg \cite{FukushimaOshidaTakeda, Sturm1, KoskelaZhou} for more details. Let $(X,\tau)$ be a topological space equipped with a $\sigma$-finite Borel measure $\mu$. A Dirichlet form $\cE$ on $(X,\mu)$ is a non-negative definite bilinear map $\cE: \cD(\cE) \times \cD(\cE) \to \setR$, with $\cD(\cE)$ being a dense subset of $L^2(X,\mu)$, satisfying closedness, meaning that the space $\cD(\cE)$ is a Hilbert space once equipped with the scalar product
\begin{equation}\label{eq:scalE}
\langle f, g \rangle_{\cE}:= \int_{X} f g \di \mu + \cE(f,g)\qquad \forall f, g \in \cD(\cE),
\end{equation}
and the Markov property: for any $f \in \cD(\cE)$, the function $f_{0}^{1} = \min (\max(f,0),1)$ belongs to $\cD(\cE)$ and $\cE(f_{0}^{1},f_{0}^{1}) \le \cE(f,f)$. We denote by $|\cdot|_\cE$ the norm associated with $\langle \cdot, \cdot \rangle_\cE$.

We focus only on symmetric Dirichlet forms, \ie those $\cE$ for which $\cE(f,g)=\cE(g,f)$ holds for all $f,g \in \cD(\cE)$. Therefore, in the rest of the article, by Dirichlet form we will always tacitly mean \textit{symmetric} Dirichlet form.

Finally, let us recall that any Dirichlet form is associated with a non-positive self-adjoint operator $L$ with dense domain $\cD(L) \subset L^2(X,\mu)$ characterized by the following:
\[
\cD(L):=\biggl\{f \in \cD(\cE) \sep \exists h=:Lf \in L^2(X,\mu)\, \, \text{s.t.}\, \, \cE(f,g)= -\int_X h g \di \mu \, \, \, \forall g \in \cD(\cE)\biggr\}.
\]

We now additionally assume that $(X,\tau)$ is locally compact and separable and that $\mu$ is a Radon measure such that $\supp \mu = X$. A Dirichlet form $\cE$ on $(X,\mu)$ is called \textit{strongly local} if $\cE(f,g)=0$ for any $f, g \in \cD(\cE)$ such that $f$ is constant on a neighborhood of $\supp g$, and \textit{regular} if $C_\rmc(X) \cap \cD(\cE)$ contains a subset (called a \textit{core}) which is both dense in $C_\rmc(X)$ for $\|\cdot\|_{\infty}$ and in $\cD(\cE)$ for $|\cdot|_\cE$. A celebrated result by Beurling and Deny \cite{BeurlingDeny} implies that any strongly local regular Dirichlet form $\cE$ on $(X,\mu)$ admits a non-negative definite symmetric bilinear map $\Gamma: \cD(\cE) \times \cD(\cE) \to \Rad$, where $\Rad$ denotes the set of signed Radon measures on $(X,\tau)$, such that
\[
\cE(f,g) = \int_X \di \Gamma(f,g) \qquad \forall f,g \in \cD(\cE),
\]
where $\int_X \di \Gamma(f,g)$ denotes the total mass of the measure $\Gamma(f,g)$. From now until the end of this paragraph, we assume that $\cE$ is strongly local and regular.

Let us mention that the map $\Gamma$ is concretely given as follows: for any $f \in \cD(\cE) \cap L^\infty(X,\mu)$, the measure $\Gamma(f):=\Gamma(f,f)$ is defined by its action on test functions:
\begin{equation}\label{eq:2000}
\int_X \phi \di \Gamma(f):= \cE(f,f\phi) - \frac{1}{2}\,\cE(f^2, \phi) \qquad \forall \phi \in \cD(\cE) \cap C_\rmc(X).
\end{equation}
Regularity of $\cE$ allows to extend \eqref{eq:2000} to the set of functions $\phi \in C_\rmc(X)$, providing a well-posed definition of $\Gamma(f)$ by duality between $C_\rmc(X)$ and $\Rad$. In case $f \in \cD(\cE)$ is not essentially bounded, $\Gamma(f)$ is obtained as the limit of the increasing sequence of measures $(\Gamma(f_{-n}^{n}))_{n\in\setN}$ where $f_{-n}^n:=\min(\max(f,-n),n)$ for any $n \in \setN$. The general expression of $\Gamma(f,g)$ for any $f,g \in \cD(\cE)$ is then obtained by polarization:
\[
\Gamma(f,g):= \frac{1}{4}\bigl(\Gamma(f+g,f+g) - \Gamma(f-g,f-g)\bigr).
\]

Strong locality of $\cE$ implies locality of $\Gamma$, that is
\[
\int_A \di \Gamma(u,w) = \int_A \di \Gamma(v,w)
\]
for any open set $A\subset X$ and any functions $u,v,w \in \cD(\cE)$ such that $u=v$ on $A$. This property allows to extend $\Gamma$ to the set $\cD_{\loc}(\cE)$ made of those $\mu$-measurable functions~$f$ for which for any compact set $K\subset X$ there exists $g \in \cD(\cE)$ such that $f=g$ $\mu$-a.e.\ on~$K$. Then $\Gamma$ satisfies the Leibniz rule:
\begin{equation}\label{eq:Leibniz}
\Gamma(fg,h)=f \Gamma(g,h) + g \Gamma(f,h) \qquad \forall u,v \in \cD_{\loc}(\cE) \cap L^{\infty}_{\loc}(X,\mu), \, \, \forall h \in \cD_{\loc}(\cE),
\end{equation}
and the chain rule:
\begin{equation}
\Gamma(\eta \circ f,g)=(\eta'\circ f) \Gamma(f,g)\quad
\begin{aligned}\label{eq:chain}
& \forall \eta \in C^1_\mathrm{b,bd}(\setR), \, \, \forall f \in \cD_{\loc}(\cE), \\
\text{or} \quad & \forall \eta \in C^1(\setR), \, \, \forall f \in \cD_{\loc}(\cE)\cap L^{\infty}(X,\mu),
\end{aligned}
\end{equation}
where $C^1_\mathrm{b,bd}(\setR)$ stands for the set of bounded $C^1$ functions on $\setR$ with bounded derivative.

For our purposes, we also need to define $\cD_{\loc}(\Omega,\cE)$ as the set of functions $f\in L^2_{\loc}(\Omega)$ for which for any compact set $K\subset \Omega$ there exists $g \in \cD(\cE)$ such that $f=g$ $\mu$-a.e.\ on~$K$; here, $\Omega$ is an open subset of $X$.

The so-called \textit{intrinsic} extended pseudo-distance $\dist_\cE$ associated with $\cE$ is defined by:
\begin{equation}\label{eq:defdist}
\dist_\cE(x,y):=\sup \{|f(x)-f(y)| \sep f \in C(X) \cap \cD_{\loc}(\cE) \, \, \, \text{s.t.} \, \, \Gamma(f) \le \mu\}\quad \forall x,y \in X.
\end{equation}
Here $\Gamma(f) \le \mu$ means that $\Gamma(f)$ is absolutely continuous with respect to $\mu$ with density lower than $1$ $\mu$-a.e.\ on $X$, and ``extended'' refers to the fact that $\dist_\cE(x,y)$ may be infinite. When the topology $\tau$ is generated by a distance $\dist$ on $X$, we call assumption~\eqref{eq:A} the following statement:
\begin{equation}\label{eq:A}\tag{A}
\text{$\dist_\cE$ is a distance inducing the same topology as $\dist$}.
\end{equation}

A final consequence of strong locality and regularity is that the operator $L$ canonically associated to $\cE$ satisfies the classical chain rule:
\begin{equation}\label{eq:chainrule}
L(\phi \circ f) = (\phi' \circ f) Lf + (\phi''\circ f) \Gamma(f) \qquad \forall f \in \mathbb{G}, \,\, \, \forall \phi \in C^\infty([0,+\infty),\setR),
\end{equation}
where $\mathbb{G}$ is the set of functions $f \in \cD(L)$ such that $\Gamma(f)$ is absolutely continuous with respect to $\mu$ with density also denoted by $\Gamma(f)$. In particular:
\begin{equation}\label{eq:chainrulesquare}
Lf^2 = 2fLf + 2 \Gamma(f) \qquad \forall f \in \mathbb{G}.
\end{equation}

\subsection*{Heat kernel associated to a Dirichlet form}

Let $(X,\tau)$ be a topological space equipped with a $\sigma$-finite Borel measure $\mu$. Let $\cE$ be a Dirichlet form on $(X,\tau)$. The spectral theorem (see \eg \cite[Th.\,VIII.5]{ReedSimon}) implies that the operator $L$ associated to $\cE$ defines an analytic sub-Markovian semi-group $(P_t)_{t>0}$ acting on $L^2(X,\mu)$, where for any $f \in L^2(X,\mu)$, the map $t \mto P_tf$ is characterized as the unique $C^1$ map $(0,+\infty)\to L^2(X,\mu)$, with values in $\cD(L)$, such that\vspace*{-3pt}
\[
\begin{cases}
\dfrac{\di}{\di t} P_tf = L(P_t f) \qquad \forall t>0,\\
\lim_{t \to 0} \|P_t f - f\|_{L^2(X,\mu)}=0.
\end{cases}
\]
One can then recover $\cD(L)$ and $L$ from $(P_t)_{t>0}$ in the following manner:\vspace*{-3pt}
\begin{align*}
\cD(L)&=\biggl\{f \in L^2(X,\mu) \sep \Bigl(\frac{P_t f - f}{t}\Bigr)_{t>0} \, \text{converges in the $L^2$-norm when $t\downarrow 0$} \biggr\},\\
Lf &= \lim_{t \downarrow 0} \frac{P_t f-f}{t} \qquad \forall f \in \cD(L).
\end{align*}

We say that $\cE$ admits a heat kernel if there exists a family of $(\mu \otimes \mu)$-measurable functions $(p(\cdot, \cdot, t))_{t>0}$ on $X\times X$ such that for all $t>0$ and $f\in L^2(X,\mu)$, one has
\[
P_t f (x) = \int_X p(x,y,t) f(y) \di \mu(y) \qquad \text{for $\mu$-a.e.\ $x \in X$};
\]
the function $p=p(\cdot,\cdot,\cdot)$ is then called the heat kernel of $\cE$.
In this case, the semi-group property (namely $P_{s+t}=P_s \circ P_t$ for any $s,t>0$) implies that $p$ satisfies the so-called Chapman-Kolmogorov property:
\begin{equation}\label{eq:ChapmanKolmogorov}
\int_X p(x,z,t)p(z,y,s) \di \mu(z) = p(x,y,t+s), \qquad \forall x,y \in X, \, \, \forall s,t>0.
\end{equation}
Moreover, for any $t>0$, $p(\cdot,\cdot,t)$ is symmetric and uniquely determined up to a $(\mu\otimes \mu)$-negligible subset of $X\times X$.

When $\cE$ admits a heat kernel, the space $(X,\tau,\mu,\cE)$ is called stochastically complete whenever\vspace*{-3pt}
\[
\int_X p(x,y,t) \di \mu(y) = 1 \qquad \forall x \in X,\, \, \forall t >0.
\]
Under stochastic completeness, one can show that\vspace*{-4pt}
\begin{multline}\label{eq:Grigoryan}
\cD(\cE)= \biggl\{f \in L^2(X,\mu) \,: \\[-5pt]
t \mto \frac{1}{2t} \iint_{X\times X} (f(x)-f(y))^2 p(x,y,t) \di \mu(x) \di \mu(y) \, \, \text{is bounded} \biggr\}
\end{multline}
and\vspace*{-3pt}
\[
\cE(f,g) = \lim_{t \downarrow 0} \frac{1}{2t} \iint_{X\times X} (f(x)-f(y))(g(x)-g(y)) p(x,y,t) \di \mu(x) \di \mu(y)
\]
for any $f, g \in \cD(\cE)$: see \cite[2.2]{Grigor'yan}, for instance. In particular, if $(X,\dist,\mu,\cE)$ is strongly local and regular, then
\begin{equation}\label{eq:Grigor'yan2}
\int_X \phi \di \Gamma(f) = \lim_{t \downarrow 0} \frac{1}{2t} \iint_{X\times X} \phi(x)(f(x)-f(y))^2 p(x,y,t) \di \mu(x) \di \mu(y)
\end{equation}
for any $f \in \cD_{\loc}(\cE)$ and $\phi \in C_\rmc(X)$.

As well-known, the classical Dirichlet energy on $\setR^n$ admits the Gaussian heat kernel
\[
p(x,y,t) = \frac{1}{(4 \pi t)^{n/2}}\, e^{-\sfrac{\dist_e^2(x,y)}{4t}} \qquad \forall x,y \in X, \, \forall t >0,
\]
where $\dist_e$ is the usual Euclidean distance. This motivates the next definition.

\begin{definition}\label{def:Euclideanheatkernel}
Let $(X,\dist,\mu)$ be a metric measure space and $\cE$ a Dirichlet form on $(X,\mu)$. We say that $(X,\dist,\mu,\cE)$ has an $\alpha$-dimensional Euclidean heat kernel if $\cE$ admits a heat kernel $p$ such that:
\[
p(x,y,t) = \frac{1}{(4 \pi t)^{\alpha/2}}\, e^{-\sfrac{\dist^2(x,y)}{4t}}\qquad \forall x,y \in X, \, \forall t>0.
\]
\end{definition}

\subsection*{Harnack inequalities}

Let $(X,\dist,\mu)$ be a metric measure space equipped with a Dirichlet form $\cE$ with associated operator $L$. Let $\leb^1$ be the Lebesgue measure on~$\setR$. In order to properly state what a Harnack inequality means for $(X,\dist,\mu,\cE)$, let us introduce some notions. We refer \eg to \cite{Sturm2} and the references therein for more details. Note first that any element $f \in L^2(X,\mu)$ uniquely defines a continuous linear form on $\cD(\cE)$, namely $g \mto \int_X fg \di \mu$. Thus $L^2(X,\mu)$ embeds into $\cD(\cE)^*$ whose norm we denote $|\cdot|_{\cE,*}$.

For any open interval $I \subset \setR$, we consider the following functional spaces:
\begin{itemize}
\item $L^2(I,\cD(\cE))$ is the space of $\leb^1$-measurable functions $u:I\to \cD(\cE)$, $u_t:=u(t)$, equipped with the Hilbert norm $\|u\|_{L^2(I,\cD(\cE))}:=(\int_I |u_t|_{\cE}^2 \di t)^{1/2}$;

\item $H^1(I,\cD(\cE)^*)$ is the space of $\leb^1$-measurable functions $u:I\to \cD(\cE)^*$ admitting a distributional derivative $\partial_t u \in L^2(I,\cD(\cE)^*)$ on $I$ equipped with the Hilbert norm $\|u\|_{H^1(I,\cD(\cE)^*)}:=(\int_I |u_t|_{\cE,*}^2 \di t + \int_I |(\partial_t u)_t |_{\cE,*}^2 \di t)^{1/2}$, where $(\partial_t u)_t:= \partial_t u(t)$;

\item $\cD_{\mathrm{par},I}(\cE):=L^2(I,\cD(\cE))\cap H^1(I,\cD(\cE)^*)$ equipped with the Hilbert norm $\|u\|_{\mathrm{par},I}:=(\int_I |u_t|_{\cE}^2 \di t + \int_I |(\partial_t u)_t |_{\cE,*}^2 \di t)^{1/2}$.
\end{itemize}
We can define a Dirichlet form $\cE_I$ on $\cD_{\mathrm{par},I}(\cE)$ by setting
\[
\cE_I(u,v):=\int_I\cE(u_t,v_t) \di t - \int_I (\partial_t u)_t \cdot v_t \di t \qquad \forall u, v \in \cD_{\mathrm{par},I}(\cE).
\]
Let $\Omega \subset X$ be an open set. Denote by $Q$ the parabolic cylinder $I \times \Omega$. Let $\cD_{Q}(\cE)$ be the set of $(\leb^1 \otimes \mu)$-measurable functions defined on $Q$ such that for every relatively compact open set $\Omega' \Subset \Omega$ and every open interval $I' \Subset I$ there exists a function $u' \in \cD_{\mathrm{par},I}(\cE)$ such that $u=u'$ on $I' \times \Omega'$. We also define $\cD_{Q,c}(\cE)$ as the set of functions $u \in \cD_{Q}(\cE)$ such that for any $t \in I$, the function $u_t$ has compact support in~$\Omega$.

\begin{definition}
We call local solution on $Q$ of the equation $(\partial_t+L)u = 0$ any function $u \in \cD_Q(\cE)$ such that $\cE_I(u,\phi)=0$ holds for any $\phi \in \cD_{Q,c}(\cE)$.
\end{definition}

When $\cE$ admits a heat kernel $p$, one can show that for any $x \in X$ and $t>0$ the function $p(x,\cdot,t)$ is a local solution of the equation $(\partial_t+L)u=0$.

The next important proposition is a combination of several famous results \cite{Grigor'yan92,Saloff-Coste,Sturm3}.

\begin{proposition}\label{prop:important}
Let $(X,\dist,\mu)$ be a metric measure space equipped with a strongly local and regular Dirichlet form $\cE$ satisfying assumption \eqref{eq:A}. Let $L$ be the operator canonically associated to $\cE$. Then the following statements are equivalent:
\begin{enumerate}
\item the combination of

\begin{enumerate}
\item
the doubling property: there exists a constant $C_D>0$ such that for any ball $B \subset X$,
\begin{equation}\label{eq:doubling}
\mu(2B)\le C_D \mu(B),
\end{equation}

\item
the local Poincaré inequality: there exists a constant $C_P>0$ such that for any $f \in \cD(\cE)$ and any ball $B \subset X$ with radius $r>0$, setting $f_B:= \mu(B)^{-1}\int_B f \di \mu$,
\begin{equation}\label{eq:Poincare}
\int_B |f-f_B|^2 \di \mu \le C_P r^2 \cE(f),
\end{equation}
\end{enumerate}

\item the existence of a heat kernel $p$ for $\cE$ satisfying double-sided Gaussian estimates: there exists $A>0$ such that for any $x,y \in X$ and any $t>0$,
\begin{equation}\label{eq:LiYau}
\frac{1}{A \mu(B_{\sqrt{t}}(x))}\, e^{-A\sfrac{\dist^2(x,y)}{t}} \le p(x,y,t) \le \frac{A}{\mu(B_{\sqrt{t}}(x))}\, e^{-\sfrac{\dist^2(x,y)}{At}},
\end{equation}

\item the parabolic Harnack inequality: there exists a constant $C_H>0$ such that for any $s\in \setR$, any ball $B$ with radius $r>0$ and any non-negative local solution $u$ on $(s-r^2,s)\times B$ of the parabolic equation $(\partial_t + L)u=0$, we have
\begin{equation}\label{eq:parabolicHarnack}
\essssup_{Q_-} (u) \le C_H \esssinf_{Q_+} (u),
\end{equation}
where $Q_-:=(s-(3/4)r^2,s-(1/2)r^2) \times (1/2)B$ and $Q_+:=(s-(1/4)r^2,s)\times (1/2)B$.
\end{enumerate}
\end{proposition}

Note that the parabolic Harnack inequality \eqref{eq:parabolicHarnack} implies the elliptic one introduced below in Lemma \ref{lem:ellipticHarnack}.

\subsection*{Locally $L$-harmonic functions}

Let $(X,\dist,\mu)$ be a metric measure space equipped with a strongly local and regular Dirichlet form $\cE$ with associated operator $L$. We set
\[
\cD_\rmc(\cE):=\{ \phi \in \cD(\cE) \, \text{with compact support}\}.
\]

\begin{definition}\label{def:localsol} Let $\Omega \subset X$ be an open set.

\begin{enumerate}
\item
We call local solution on $\Omega$ of the Laplace equation $Lu=0$ any function $u \in \cD_{\loc}(\Omega,\cE)$ such that $\cE(u,\phi)=0$ holds true for any $\phi \in \cD_\rmc(\cE)$ with $\supp \phi \subset \Omega$; here $\supp \phi$ denotes the support of the measure $A \mto \int_A \phi \di \mu$.

\item
We call locally $L$-harmonic function any function $u \in \cD(\cE)$ such that $\cE(u,\phi)=0$ holds true for any $\phi \in \cD_\rmc(\cE)$.

\item
For any $f \in L^1_{\loc}(X,\mu)$, we call local solution on $\Omega$ of the Poisson equation $Lu=f$ any function $u \in \cD_{\loc}(\Omega, \cE)$ such that $\cE(u,\phi)=-\int_X f \phi \di \mu$ holds true for any $\phi \in \cD_\rmc(\cE)$ with $\supp \phi \subset \Omega$.
\end{enumerate}
\end{definition}

We shall often simply write ``$Lu=f$ on $\Omega$'' to mean that $u \in \cD_{\loc}(\Omega, \cE)$ is a local solution on $\Omega$ of the equation $Lu=f$, and ``$Lv=0$'' to express that $v \in \cD_{\loc}(\cE)$ is locally $L$-harmonic. Lastly, we point out that strong locality directly implies that constant functions are locally $L$-harmonic, \ie
\[
L\mathbf{1}=0.
\]

Let us state a classical lemma (Liouville theorem under elliptic Harnack inequality) whose proof is omitted here (see \eg \cite[Lem.\,7.3]{AldanaCarronTapie}).

\begin{lemma}\label{lem:ellipticHarnack}
Let $(X,\dist,\mu)$ be a metric measure space equipped with a Dirichlet form~$\cE$ whose associated operator $L$ satisfies an elliptic Harnack inequality, meaning that there exists a constant $C_E>0$ such that for any ball $B \subset X$ and any non-negative local solution $h$ of $Lu=0$ on $B$, we have
\begin{equation}\label{eq:ellipticHarnack}
\essssup_{(1/2)B} \, h \le C_E \esssinf_{(1/2)B} \, h.
\end{equation}
Then any non-negative locally $L$-harmonic function is constant.
\end{lemma}

\subsection*{Strongly harmonic functions}

Let $(X,\dist,\mu)$ be a metric measure space. Following the terminology adopted in \cite{GaczkowskiGorka, AGG19}, for any open set $\Omega \subset X$ we call strongly harmonic function on $\Omega$ any function $h:\Omega\to\setR$ satisfying the mean value property:
\[
h(x) = \fint_{B_r(x)} h \di \mu \qquad \forall x \in \Omega, \, \, \forall \, r\in(0, \dist(x,^\rmc\!\Omega)).
\]
\begin{remark}\label{rem:defalternative}
It can easily be checked that a function $h:\Omega\to\setR$ is strongly harmonic if and only if for any $x\in \Omega$ and any $u\in C_\rmc^1\left(\left[0, \dist(x,^\rmc\!\Omega)\right]\right)$ with
$\int_X u(\dist(x,y))\di \mu(y) = 1$ one has
\[
h(x) = \int_{X}u\left(\dist(x,y)\right) h(y) \di \mu(y).
\]
\end{remark}
Under mild assumptions on $(X,\dist,\mu)$, an elliptic Harnack inequality holds true for strongly harmonic functions, provided the doubling condition \eqref{eq:doubling} is satisfied: see \cite[Lem.\,4.1]{AGG19}. The next lemma is an easy consequence of this fact. We recall that a metric space is called proper if any closed ball is compact, and that proper metric spaces are complete and locally compact.

\begin{lemma}\label{lem:Harnackstrongly}
Let $(X,\dist)$ be a proper metric space equipped with a regular Borel measure $\mu$ such that $0<\mu(B)<+\infty$ for any metric ball $B \subset X$. Assume that $(X,\dist,\mu)$ satisfies the doubling condition \eqref{eq:doubling}. Then any non-negative strongly harmonic function on $X$ is constant.
\end{lemma}

When $(X,\dist,\mu)$ has an $\alpha$-dimensional volume, strongly harmonic functions satisfy the following two properties:

\begin{lemma}\label{lem:0109}
Let $(X,\dist,\mu)$ be with an $\alpha$-dimensional volume and $h:X\to\setR$ be strongly harmonic. Then:
\begin{enumeratei}
\item\label{lem:0109i}
if $h$ has linear growth -- meaning that there exists $C>0$ such that $|h|\le C(1+\dist(o,\cdot))$ for some $o \in X$ -- then $h$ is Lipschitz;
\item\label{lem:0109ii}
if $h$ is continuous and such that $\sup_{\partial B_{r_i}(o)} |h| =o(r_i)$ for some point $o \in X$ and some sequence $\{r_i\}_i \subset (0,+\infty)$ such that $r_i \to +\infty$, then $h$ is constant.
\end{enumeratei}
\end{lemma}

\begin{proof}
Let us first prove \eqref{lem:0109i}. Assuming $h$ to have linear growth, we know that there exists $o \in X$, $r_o>0$ and $M>0$ such that $|h(z)| \le M\dist(o,z)$ for all $z \in X \backslash B_{r_o}(o)$. Since $h$ is strongly harmonic, we have
\[
\mu(B_{r+d}(x)) h(x) - \mu(B_r(y)) h(y) = \int_{B_{r+d}(x) \backslash B_r(y)} h \di \mu
\]
for all $r>0$ and any given $x,y \in X$, where we have set $d:=\dist(x,y)$. Since $\mu(B_{r+d}(x)) = \omega_\alpha (r+d)^\alpha$ and $\mu(B_r(y)) = \omega_\alpha r^\alpha$, we obtain
\begin{equation}\label{eq:D4.1}
\begin{aligned}
|\omega_\alpha (r+d)^\alpha h(x) - \omega_\alpha r^\alpha h(y)| & \le \omega_\alpha ((r+d)^\alpha - r^\alpha)\, \mathrm{sup}_{B_{r+d}(x)\backslash B_r(y)} \, |h| \\
& \le \omega_\alpha ((r+d)^\alpha - r^\alpha) \, \mathrm{sup}_{B_{r+d+\dist(o,x)}(o)\backslash B_r(y)} \, |h|
\end{aligned}
\end{equation}
since $B_{r+d}(x) \subset B_{r+d+\dist(o,x)}(o)$. Choosing $r>r_o+\dist(o,y)$ in order to ensure that $B_r(y)$ contains $B_{r_o}(o)$, we get $ \mathrm{sup}_{B_{r+d+\dist(o,x)}(o)\backslash B_r(y)} \, |h| \le M(r+d+\dist(o,x))$,
hence
\[
|(1+d/r)^\alpha h(x) - h(y)| \le ((1+d/r)^\alpha-1)M(r+d+\dist(o,x)).
\]
Letting $r \to +\infty$ and applying $(1+d/r)^\alpha-1 =\alpha d/r +o(1/r)$ yields to $|h(x) - h(y)| \le \alpha dM$.

To prove \eqref{lem:0109ii}, apply \eqref{eq:D4.1} with $r=R_i:=r_i-d-\dist(o,x)$ to get
\[
\left| h(x) (1+d/R_i)^\alpha - h(y)\right| \le ((1+d/R_i)^\alpha-1) \,\mathrm{sup}_{B_{r_i}(o)} |h|.
\]
By the weak maximum principle \cite[Cor.\,4.3]{AGG19}, we have
\[
\sup_{B_{r_i}(o)} |h| = \sup_{\partial B_{r_i}(o)} |h|.
\]
Since $(1+d/R_i)^\alpha-1= \alpha d/ R_i + o(1/R_i) = O(1/r_i)$ when $i\to +\infty$, then there exists $i_o$ and $C>0$ such that
\[
\left| h(x) (1+d/R_i)^\alpha - h(y)\right| \le C r_i^{-1}\mathrm{sup}_{\partial B_{r_i}(o)} |h|
\]
for all $i \ge i_o$. This implies $h(x)=h(y)$ by letting $i$ tend to $+\infty$.
\end{proof}

\subsection*{Tangent cones at infinity}

We refer to \cite{Gromov} for a definition of the Gromov-Hausdorff distance $\dist_{\sGH}$ between compact metric spaces and only mention here that a sequence of compact metric spaces $\{(X_i,\dist_i)\}$ converges to another compact metric space $(X,\dist)$ with respect to the Gromov-Hausdorff distance (what we denote by $\dist_{\sGH}(X_i,X) \to 0$) if and only if there exists an infinitesimal sequen\-ce $\{\eps_i\}_i \subset (0,+\infty)$ and functions $\phi_i: X_i \to X$ called $\eps_i$-isometries such that \hbox{$|\dist(\phi_i(x),\phi_i(x'))-\dist_i(x,x')|\le \eps_i$} for any $x,x' \in X_i$ and any~$i$. If $x_i \in X_i$ for any~$i$ and $x \in X$ are such that $\dist(\phi_i(x_i),x) \to 0$, we write $x_i \To{\sGH} x$.

When dealing with non-compact spaces, we say that a sequence of pointed metric spaces $\{(X_i, \dist_i, x_i)\}_i$ converges in the pointed Gromov-Hausdorff topology to $(X, \dist, x)$ if there exist sequences of positive numbers $\eps_i \downarrow 0$, $R_i \uparrow \infty$, and of Borel maps $\phi_i:B_{R_i}(x_i) \to X$, also called $\eps_i$-isometries, such that such that for any $i$ the ball $B_{R_i}(x)$ is included in the $\eps_i$-neighborhood of $\phi_i(B_{R_i}(x_i))$, namely $ \bigcup_{y \in \phi_i(B_{R_i}(x_i))} B_{\eps_i}(y)$, $|\dist_i(y, z)-\dist(\phi_i(y), \phi_i(z))|<\eps_i$ for any $y,\, z \in B_{R_i}(x_i)$, and $\dist(\phi_i(x_i),x)\to 0$ (which we also abbreviate to $x_i \To{\sGH} x$).

Pointed measured Gromov-Hausdorff convergence of a sequence of pointed metric measure spaces $\{(X_i, \dist_i, \mu_i, x_i)\}$ to $(X, \dist,\mu,x)$ is set as pointed Gromov-Hausdorff convergence of $\{(X_i, \dist_i, x_i)\}$ to $(X, \dist, x)$ with the additional requirement \hbox{$(\phi_i)_{\sharp}\mu_i \!\stackrel{C_\bs(X)}{\rightharpoonup}\! \mu$}, where $C_\bs(X)$ is the space of continuous functions with bounded support and $ f_\sharp $ is
the push forward operator between measures induced by a Borel map $f$.

A metric space $(X,\dist)$ is called metric doubling if there exists a positive integer $N$ such that any ball in $(X,\dist)$ can be covered by at most $N$ balls with half its radius. Whenever $(X,\dist)$ is a doubling space, for any $o \in X$, the family of pointed spaces $\{(X,r^{-1}\dist,o)\}_{r>1}$ satisfies the assumptions of Gromov's precompactness theorem \cite[Prop.\,5.2]{Gromov}, henceforth it admits limit points in the pointed Gromov-Haudorff topology as $r\uparrow +\infty$. These pointed metric spaces are called tangent cones at infinity of $(X,\dist)$ in $o$.

It is well-known that when $(X,\dist,\mu)$ is satisfies the doubling property \eqref{eq:doubling}, then the metric space $(X,\dist)$ is metric doubling: see \eg \cite[\S 2.5]{AmbrosioColomboDiMarino}.

When a metric measure space $(X,\dist,\mu)$ has an $\alpha$-dimensional volume, a simple computation shows that it is measure doubling, with $C_D=2^\alpha$. Moreover, one can equip any of its tangent cones at infinity $(\uX,\udist,\uo)$ with a limit measure $\umu$ in the following way. Let $\{r_i\}_i$ be a sequence of positive real numbers diverging to $+\infty$ such that $(\uX,\udist,\uo)$ is the pointed Gromov-Hausdorff limit of $\{(X,r_i^{-1}\dist,o)\}_i$. Set $\mu_i:=r_i^{-\alpha}\mu$ for any $i$, and note that
\[
\mu_i(B_r^{\dist_i}(x)) = \mu_i(B_{r r_i}(x)) = \omega_\alpha r^\alpha \qquad \forall x \in X,\; r>0.
\]
Set $\underline{V}(x,r):=\omega_\alpha r^\alpha$ for any $x \in X$ and $r>0$. Then for any $\delta>0$ and any Borel set~$A$ of $(\uX,\udist)$, setting
\[\textstyle
\umu_\delta(A):= \inf\bigl\{ \sum_{i} \underline{V}(z_i,r) \sep \{B_{r_i}(z_i)\}_i\, \, \text{s.t. $A \subset \bigcup_i B_{r_i}(z_i)$ and $r_i \le \delta$} \bigr\}
\]
and then $\umu(A) = \lim_{\delta \to 0} \umu_\delta(A)$ defines a metric outer measure $\umu$ on $(\uX,\udist)$ whose canonically associated measure, still denoted by $\umu$, is a Radon measure satisfying $\umu(B_r(\ux)) = \omega_\alpha r^\alpha$ for any $\ux \in \uX$ and $r>0$. This shows that $(\uX,\udist,\umu)$ has an $\alpha$\nobreakdash-dimen\-sional volume. Moreover, we obviously have $\umu(B_r(\ux))=\lim_{i \to +\infty} \mu_i(B_r^{\dist_i}(x_i))$ for any $r>0$ and any sequence $x_i \To{\sGH} x$; by density in $C_\bs(\uX)$ of the space spanned by the collection of characteristic functions of balls, this implies the pointed measured Gromov-Hausdorff convergence $(X,r_i^{-1}\dist,\mu_i,o) \to (\uX,\udist,\umu,\uo)$.

\subsection*{Ascoli-Arzelà type theorems}

Let $\{(X_i,\dist_i,x_i)\}_i, (X,\dist,x)$ be pointed proper metric spaces such that
\[
(X_i,\dist_i,x_i) \to (X,\dist,x)
\]
in the pointed Gromov-Hausdorff topology and $\phi_i:B_{R_i}(x_i) \to X$ be $\eps_i$-isometries, where $\{\eps_i\}_i, \{R_i\}_i \subset (0,+\infty)$ are such that $\eps_i \downarrow 0$ and $R_i \uparrow +\infty$. For any $i$, let $K_i$ be a compact subset of $X_i$, and assume that there exists $K \subset X$ compact such that $\dist_{\sGH}(K_i,K) \to 0$. We say that functions $f_i:X_i\to\setR$ converge to $f:X\to\setR$ uniformly over $K_i \to K$ if $\sup_{K_i} |f_i - f \circ \phi_i| \to 0$. Note that this definition depends on the choice of the $\eps_i$-isometries $\phi_i$ that we keep fixed for the rest of this paragraph.

\begin{remark}
In the rest of the article, whenever we consider a convergent sequen\-ce of pointed metric spaces $(X_i,\dist_i,x_i) \to (X,\dist,x)$, we always implicitly assume that sequences $\{\eps_i\}_i, \{R_i\}_i \subset (0,+\infty)$ with $\eps_i\downarrow 0, R_i \uparrow +\infty$ and $\eps_i$-isometries $\phi_i:B_{R_i}(x_i)\to X$ have been chosen a priori and that the statements ``$x_i \To{\sGH} x$'' and ``$f_i \to f$ uniformly on compact sets'' are meant with these $\eps_i$-isometries.
\end{remark}

In this context, we have the following Ascoli-Arzelà theorem:

\begin{proposition}\label{prop:AA1}
Let $\{(X_i,\dist_i,x_i)\}_i, (X,\dist,x)$ be as above, and $r>0$. For any $i$, let $f_i \in C(X_i)$ be such that:
\begin{itemize}
\item $\sup_i \|f_i\|_{L^{\infty}(\overline{B}_r(x_i))} < +\infty$,
\item the sequence $\{f_i\}_i$ is asymptotically uniformly continuous on $\overline{B}_r(x)$ (see \cite[Def.\,3.2]{Honda}).
\end{itemize}
Then $\{f_i\}_i$ admits a subsequence $(f_{i(j)})_j$ which converges to $f$ uniformly over $\overline{B}_r(x_i) \to \overline{B}_r(x)$.
\end{proposition}

\begin{proof}
From \cite[Prop.\,3.3]{Honda}, we know that for $\{f_i\}_i$ satisfying the above assumptions, there exists $f \in C(B_r(x))$ and a subsequence $(f_{i(j)})_j$ such that $f_{i(j)}(x_j) \to f(x)$ whenever $x_j \To{\sGH} x \in B_r(x)$. With no loss of generality, we can assume that the subsequence is the whole sequence itself. By contradiction, assume that the uniform convergence $f_i\to f$ over $\overline{B}_r(x_i)\to \overline{B}_r(x)$ is not satisfied. Then there is some $\eps>0$ and a subsequence $(f_{i(\ell)})_\ell$ such that $\inf_\ell \{\sup_{\overline{B}_r(x_{i(\ell)})} |f_{i(\ell)}-f \circ \phi_{i(\ell)}|\} \ge \eps$. Again, we can assume that the subsequence is the whole sequence itself. For any $i$, choose $y_i \in \overline{B}_r(x_i)$ such that $|f_i(y_i)-f \circ \phi_i(y_i)| \ge \eps/2$ and set $z_i:=\phi_i(y_i) \in \overline{B}_{r+\eps_i}(x)$. Properness of $X$ implies that the sequence $\{z_i\}_i$ converges to some $z \in \overline{B}_r(x)$, up to extraction. In particular, $y_i \To{\sGH} z$. Then in
\[
\eps/2 \le |f_i(y_i)-f(z)| + |f(z)-f\circ \phi_i(y_i)|,
\]
the first term in the right-hand side goes to $0$ when $i$ tend to $+\infty$. Since $f$ is continuous, we also have $|f(z)-f\circ \phi_i(y_i)|\to 0$ when $i$ tend to $+\infty$, hence a contradiction.
\end{proof}

Let $\{(X_i,\dist_i,x_i)\}_i, (X,\dist,x),\{\phi_i\}_i$ be as above. Let $(Y,\dist_Y)$ be another metric space. We say that $f_i: Y\to X_i$ converge to $f: Y\to X$ uniformly on compact subsets of $Y$ if $\sup_K \dist_i(\phi_i \circ f_i,f) \to 0$ for any compact set $K \subset Y$.

An Ascoli-Arzelà theorem is also available in this context. We state it with an equi-Lipschitz assumption which is enough for our purposes. The proof is omitted for brevity.

\begin{proposition}\label{prop:AA2}
Let $\{(X_i,\dist_i,x_i)\}_i, (X,\dist,x)$ be as above. Let $(Y,\dist_Y)$ be a metric space and $f_i: Y\to X_i$ be Lipschitz functions such that:
\begin{itemize}
\item $L:=\sup_i \Lip(f_i) < +\infty$,
\item there exists $y\in Y$ and $r>0$ such that $\dist_i(f_i(y),x_i)\le r$ for any $i$.
\end{itemize} Then $\{f_i\}_i$ admits a subsequence converging uniformly on compact sets of $Y$ to some Lipschitz function $f:Y\to X$, and $\Lip(f) \le L$.
\end{proposition}

Let us conclude this paragraph with a stability result for strongly harmonic functions.

\begin{proposition}\label{prop:stabstrongharmonic}
Let $\{(X_i,\dist_i,\mu_i,x_i)\}_i, (X,\dist,\mu,x)$ be proper pointed metric measured spaces such that $(X_i,\dist_i,\mu_i,x_i) \to (X,\dist,\mu,x)$ in the pointed measured Gromov-Hausdorff topology. Let $f_i \in C(X_i)$ be converging to $f \in C(X)$ uniformly over $\overline{B}_r(x_i) \to \overline{B}_r(x)$ for any $r>0$. Assume that $f_i$ is strongly harmonic for any $i$. Then $f$ is strongly harmonic.
\end{proposition}

\begin{proof}
By the characterization of strongly harmonic functions stated in Remark~\ref{rem:defalternative}, it is enough to establish
\begin{equation}\label{eq:fstrongharm}
f(y) = \fint_{B_r(y)}u(\dist(y,z))f(z)\di \mu(z)
\end{equation}
for any given $r>0$, $y \in X$ and $u \in C^1_\rmc([0,+\infty))$ such that $\int_X u(\dist(y,z)) \di \mu(z)=1$.
Let $y_i \in X_i$ for any $i$ be such that $y_i \To{\sGH} y$. For any $i$, set
\[
u_i:=\frac{u}{\int_{X_i} u(\dist_i(y_i,z)) \di \mu_i(z)}
\]
and note that
\begin{equation}
\int_{X_i} u_i(\dist_i(y_i,z)) \di \mu_i(z) = 1
\end{equation}
so that $f_i$ being strongly harmonic implies
\begin{equation}\label{eq:fistronharm}
f_i(y_i) = \fint_{B_r(y_i)}u_i(\dist_i(y_i,z))f_i(z)\di \mu_i(z).
\end{equation}
But
\begin{multline*}
\int_{X_i} u(\dist_i(y_i,z)) \di \mu_i(z) = - \int_0^{+\infty} u'(r) \mu_i(B_r^{\dist_i}(y_i)) \di r\\
\to - \int_0^{+\infty} u'(r) \mu(B_r^{\dist}(y)) \di r = \int_{X} u(\dist(y,z)) \di \mu(z)=1,
\end{multline*}
so $u_i \to u$ uniformly on $(0,+\infty)$: this implies that the functions $u_i(\dist_i(y_i,\cdot))f_i \in C(X_i)$ converge uniformly over all compact sets to $u(\dist(y,\cdot))f \in C(X)$. Therefore, letting $i$ tend to $+\infty$ in \eqref{eq:fistronharm} provides \eqref{eq:fstrongharm}.
\end{proof}

\subsection*{Length structures}

Let $(X,\dist)$ be a metric space. A path in $X$ is a continuous map $c:[0,1]\to X$. It is called rectifiable if its length
\[\textstyle
L_\dist(c):=\sup \bigl\{ \sum_{i=1}^n \dist(c(t_i),c(t_{i-1})) \, \sep \, 0 = t_0 < \ldots < t_n = 1, \, \, n \in \setN\backslash\{0\} \bigr\}
\]
is finite. $(X,\dist)$ is called length metric space if for any $x, y \in X$,
\[
\dist(x,y) = \inf\{L_\dist (c) \sep c \in \Omega_{xy}\},
\]
where $\Omega_{xy}$ is the set of rectifiable paths in $X$ such that $c(0)=x$ and $c(1)=y$. A~geodesic space is a trivial example of length space. Equivalently, $(X,\dist)$ is length if~$\dist$ coincides with its associated length distance $\overline{\dist}$ defined by:
\[
\overline{\dist}(x,y):= \inf\{L_\dist (c) \sep c \in \Omega_{xy}\}\qquad \forall x,y \in X,
\]
in which case we say that $\dist$ is a length distance. Note that we always have $\dist \le \overline{\dist}$ and $L_{\dist}(c)=L_{\overline{\dist}}(c)$ whenever $c$ is a rectifiable path in $X$. Moreover,
\[
L_\dist(c) = \lim_{\alpha \to 0^+} L_{\dist,\alpha}(c),
\]
where
\begin{multline*}
\textstyle
L_{\dist,\alpha}(c) = \sup \bigl\{ \sum_{i=1}^n \dist(c(t_i),c(t_{i-1})) \sep\\ 0 = t_0 < \ldots < t_n = 1, |t_i - t_{i+1}|<\alpha \, \, \forall i,\, \, n \in \setN\backslash\{0\} \bigr\}.
\end{multline*}
In this context, we have the following lemma.

\begin{lemma}\label{lem:length}
Let $(X,\delta)$ be a length metric space. Assume that $\dist$ defined as $\dist:=2\sin(\delta/2)$ is a distance. Then its associated length distance $\overline{\dist}$ coincides with $\delta$.
\end{lemma}

\begin{proof}
First note that $(2/\pi) \delta \le \dist \le \delta$ because $(2/\pi)x\le 2\sin(x/2)\le x$ for any $x\in [0,\pi]$. In particular, a map $c:[0,1]\to X$ is continuous for $\delta$ if and only if it is for $\dist$. Moreover, since $\delta$ is a length distance, $\dist \le \delta$ implies $\overline{\dist} \le \delta$, so we are left with proving the converse inequality. Let $c$ be a path in $X$. Being continuous, $c$ is also uniformly continuous: for any $\eps\in (0,1)$, there exists $\alpha>0$ such that for any $t,s \in [0,1]$,
\[
|t-s| < \alpha \implies \delta(c(t),c(s)) < \eps.
\]
Since $x- 2\sin(x/2)\le x^2$ for any $x\ge0$, then
\[
\delta - \dist \le \delta^2,
\]
so that for any $t,s \in [0,1]$:
\[
|t-s| < \alpha \implies \delta(c(t),c(s)) - \dist(c(t),c(s)) \le \eps \delta(c(t),c(s)).
\]
This implies $L_{\delta,\alpha}(c) - L_{\dist,\alpha}(c) \le \eps L_{\delta,\alpha}(c)$ and thus $(1-\eps)L_\delta(c) \le L_\dist(c)$ by letting $\alpha$ tend to $0$. Letting $\eps$ tend to $0$ provides $L_\delta(c) \le L_\dist(c)$. This implies $\delta \le \overline{\dist}$.
\end{proof}

\subsection*{Busemann functions}

Let $(X,\dist)$ be a metric space. A geodesic ray in $X$ is a continuous function $\gamma:[0,+\infty)\to X$ such that $\dist(\gamma(t),\gamma(s))=|t-s|$ for any $s,t\ge 0$. The Busemann function associated to a geodesic ray $\gamma$ is defined by
\[
b_\gamma(x)=\lim_{t \to +\infty} t - \dist(x,\gamma(t)).
\]
Note that this limit is well-defined for any $x \in X$ since the function $t \mto t - \dist(x,\gamma(t))$ is non-decreasing and bounded from above by $\dist(o,x)$. Note also that $b_\gamma$ is $1$-Lipschitz, since for any $x, y \in X$ and any $t>0$, one has $t - \dist(x,\gamma(t)) - (t - \dist(y,\gamma(t)) \le \dist(x,y)$, henceforth $b_\gamma(x) - b_\gamma(y) \le \dist(x,y)$ by letting $t \to +\infty$. Moreover, for any $s>0$, one can easily check that
\begin{equation}%\label{eq:Busemann}
b_\gamma(\gamma(s))=s.
\end{equation}

We shall need the following lemma.

\begin{lemma}\label{lem:preparatory}
Let $(X,\dist,\mu)$ be a metric measure space equipped with a strongly local and regular Dirichlet form $\cE$ with associated operator $L$. Assume that:
\begin{enumerate}
\item\label{lem:preparatory1} $\Lip(X)$ embeds continuously in $\cD_{\loc}(\cE)$, \ie there exists $C>0$ such that $\di\Gamma(f)\le C \Lip^2(f)\di \mu$ for any $f \in \Lip(X)$,
\item\label{lem:preparatory2} there exists $\alpha>-1$ such that $L\dist(x,\cdot)=\alpha/\dist(x,\cdot)$ on $X\backslash \{x\}$ for any $x \in X$.
\end{enumerate}
Then any Busemann function on $(X,\dist)$ is locally $L$-harmonic.
\end{lemma}

\begin{proof}
Let $b_\gamma$ be a Busemann function on $(X,\dist)$. Since $b_\gamma$ is $1$-Lipschitz, we have $b_\gamma \in \cD_{\loc}(\cE)$. To conclude we must prove that $\cE(b_\gamma,\phi)=0$ for any $\phi \in \cD_\rmc(\cE)$.

Recall that strong locality implies that constant functions are locally $L$-harmonic. Set $f_s:=s-\dist(\gamma(s),\cdot)$ for any $s>0$ and note that \eqref{lem:preparatory2} implies that $f_s$ is a local solution of $Lu = -\alpha/\dist(\gamma(s),\cdot)$ on $X \backslash \{\gamma(s)\}$. Then for any $\phi \in \cD_\rmc(\cE)$, since for $s$ large enough $\gamma(s) \notin \supp \phi$ and thus $\dist(\gamma(s),\cdot)>0$ on $\supp \phi$, we get
\[
|\cE(f_s,\phi)|= \left|\int_X (Lf_s) \phi \di \mu\right| \le \frac{|\alpha|}{\dist(\gamma(s),\supp \phi)} \int_X |\phi| \di \mu \to 0 \quad \text{when $s \to +\infty$}.
\]
Thus to conclude, it suffices to show that for any $\phi \in \cD_\rmc(\cE)$,
\[
\lim_{s \to +\infty} \cE(f_s,\phi)=\cE(b_\gamma,\phi).
\]
Let $\phi \in \cD_\rmc(\cE)$ and $K\subset X$ be a compact set containing the support of $\varphi$.
We introduce
\[
\cD_{K}(\cE):=\left\{ f \in L^2(K,\mu): \int_{K}\di \Gamma(f)<+\infty \right\},
\]
which is a Hilbert space when equipped with the norm
\[
\|f\|_{\cD_{K}(\cE)}:=\left(\|f\|^2_{L^2(K,\mu)} + \int_{K}\di \Gamma(f)\right)^{1/2}.
\]
As $(f_s)$ is increasing and converges pointwise to $b_\gamma$, we have $f_s \to b_\gamma$ in $L^2_{\loc}(X,\mu)$. As a consequence, the sequence $(\left.f_s\right|_K)$ converges to $\left.b_\gamma\right|_K$ weakly in $L^2(K,\mu)$. Since the functions $f_s$ are all $1$-Lipschitz, by \eqref{lem:preparatory1} we get that the sequence $(\left.f_s\right|_K)$ is bounded in $\cD_{K}(\cE)$, so it admits limit points for the weak convergence in $\cD_{K}(\cE)$. This weak convergence implies the weak convergence in $L^2(K,\mu)$. Therefore, by uniqueness of the weak limit in $L^2(K,\mu)$, we get that $\left.b_\gamma\right|_K$ is the unique limit point of $(\left.f_s\right|_K)$ for the weak convergence in $\cD_{K}(\cE)$. The strong locality implies that
\[
\cE(f_s,\phi)=\cE(\left.f_s\right|_K,\phi)\quad\text{and}\quad \cE(b_\gamma,\phi)=\cE(\left.b_\gamma\right|_K,\phi).
\]
Moreover $v\in \cD_{K}(\cE)\mto \cE(v,\phi)$ is continuous hence
\[
\cE(f_s,\phi)\to \cE(b_\gamma,\phi).\qedhere
\]
\end{proof}

\subsection*{Laplace transform}

Let $F:[0,+\infty)\to \setR$ be a locally integrable function such that $F(t)=O(e^{\gamma t})$ when $t \to +\infty$ for some $\gamma \in \setR$. The Laplace transform of $F$ is the complex-valued function $\mathcal{L}\{F\}$ defined by:
\[
\mathcal{L}\{F\}(z) = \int_0^{+\infty} F(\xi)e^{-z\xi}\di \xi, \qquad \forall z \in \{\mathrm{Re}(\cdot) > \gamma\}.
\]
Lerch's theorem asserts that if $F_1, F_2: [0,+\infty)\to\setR$ are two continuous, locally integrable functions satisfying $F_1(t), F_2(t)=O(e^{\gamma t})$ when $t \to +\infty$ and $\mathcal{L}\{F_1\} = \mathcal{L}\{F_2\}$ on $\{\mathrm{Re}>\gamma\}$, then $F_1=F_2$ (see \eg \cite[Th.\,2.1]{Cohen}). This provides the following lemma.

\begin{lemma}\label{lem:Laplace}
Let $F:(0,+\infty)\to \setR$ be a continuous and locally integrable function such that $F(t)=O(e^{\eps t})$ when $t \to +\infty$ for any $\eps>0$ and
\begin{equation}\label{eq:Laplace}
\mathcal{L}\{F\}(\lambda) = \lambda^{-\alpha-1} \qquad \forall \lambda >\lambda_o
\end{equation}
for some $\alpha>0$ and $\lambda_o \ge 0$. Then $F(\xi)=\xi^\alpha/\Gamma(\alpha+1)$ for any $\xi\ge0$.
\end{lemma}

\begin{proof}
Since $F$ is locally integrable, one can apply the classical theorem on holomorphy under the integral sign to get that $\mathcal{L}\{F\}$ is holomorphic on any compact subset of $\{\mathrm{Re}>0\}$. Therefore, by analytic continuation, \eqref{eq:Laplace} implies $\mathcal{L}\{F\}(z) = z^{-\alpha-1}$ for any $z \in \{ \mathrm{Re}(\cdot)>0\}$. Since the Laplace transform of $\xi \mto \xi^{\alpha}$ is $z\mto \Gamma(\alpha +1) z^{-\alpha-1}$, Lerch's theorem gives $F(\xi)=\xi^\alpha/\Gamma(\alpha+1)$ for any $\xi\ge0$.
\end{proof}

\section{First rigidity results for spaces with an Euclidean heat kernel}\label{sec:3}

In this section, we establish several properties of metric measure spaces equipped with a Dirichlet form admitting an $\alpha$-dimensional Euclidean heat kernel. We shall use most of these results in the next section to prove Theorem \ref{th:main}.

\subsection{Stochastic completeness and consequences}

We begin with stochastic completeness.

\begin{lemma}\label{lem:stochastic}
Let $(X,\dist,\mu,\cE)$ be with an $\alpha$-dimensional Euclidean heat kernel.
Then $(X,\dist,\mu,\cE)$ is stochastically complete.
\end{lemma}

\begin{proof}
Take $t,s>0$ and $x \in X$. By \eqref{eq:ChapmanKolmogorov}, for any $y \in X$ we have
\begin{equation}\label{eq:Lem3.1}
\int_X p(x,z,t)e^{-\sfrac{\dist^2(z,y)}{4s}} \di \mu(z) = \Bigl(\frac{s}{t+s}\Bigr)^{\alpha/2} e^{-\sfrac{\dist^2(x,y)}{4(t+s)}}.
\end{equation}
Letting $s \to +\infty$ and applying the monotone convergence theorem, we get the result.
\end{proof}

As a consequence of Lemma \ref{lem:stochastic}, we can show that spaces with an $\alpha$-dimensional Euclidean heat kernel have an $\alpha$-dimensional volume.

\begin{lemma}\label{lem:euclideanvolume}
Let $(X,\dist,\mu,\cE)$ be with an $\alpha$-dimensional Euclidean heat kernel. Then $(X,\dist,\mu)$ has an $\alpha$-dimensional volume.
\end{lemma}

\begin{proof}
Take $x \in X$. By Lemma \ref{lem:stochastic}, we have:
\[
\int_X p(x,y,t) \di \mu(y) = 1 \qquad \forall \, t>0,
\]
so that the hypothesis on the heat kernel implies:
\begin{equation}\label{eq:001}
\int_X e^{-\sfrac{\dist^2(x,y)}{4t}} \di \mu(y) = (4\pi t)^{\alpha/2} \qquad \forall \, t>0.
\end{equation}
By Cavalieri's principle (see for instance \cite[Lem.\,5.2.1]{AmbrosioTilli}), we have
\[
\int_X e^{-\sfrac{\dist^2(x,y)}{4t}}\di \mu(y) = \int_0^{+\infty} \mu(\{e^{-\sfrac{\dist^2(x,\cdot)}{4t}}>s\}) \di s.
\]
Since for any $y \in X$, one has $e^{-\sfrac{\dist^2(x,y)}{4t}}>s$ if and only if $\dist^2(x,y) < - 4t \log(s)$, then
\[
\Bigl\{e^{-\sfrac{\dist^2(x,\cdot)}{4t}}>s\Bigr\}
= \begin{cases}
\emptyset & \text{if $s\ge 1$},\\
B_{\sqrt{-4t\log(s)}}(x) & \text{if $s<1$}.
\end{cases}
\]
Therefore, the change of variable $\xi=-4t\log(s)$ yields to
\[
\int_X e^{-\sfrac{\dist^2(x,y)}{4t}}\di \mu(y) = \frac{1}{4t} \int_0^{+\infty} e^{-\sfrac{\xi}{4t}} \mu(B_{\sqrt{\xi}}(x)) \di \xi.
\]
Coupled with \eqref{eq:001}, and setting $\lambda=1/(4t)$, this leads to:
\[
\int_0^{+\infty} e^{-\lambda \xi} \mu(B_{\sqrt{\xi}}(x)) \di \xi = \pi^{\alpha/2} \lambda^{-\alpha/2-1} \qquad \forall \lambda>0.
\]
Applying Lemma \ref{lem:Laplace} and \eqref{eq:omega_n} provides the result.
\end{proof}

A second consequence is that complete spaces with an $\alpha$-dimensional Euclidean heat kernel are proper; in particular, they are locally compact. Note that the space $(\setR^n \backslash \{0\}, \dist_e,\leb^n)$ shows that completeness is a non-removable assumption.

\begin{lemma}
Let $(X,\dist,\mu,\cE)$ be with an $\alpha$-dimensional Euclidean heat kernel and such that $(X,\dist)$ is complete. Then any closed ball in $X$ is compact.
\end{lemma}

\begin{proof}
From Lemma \ref{lem:euclideanvolume}, we know that $(X,\dist,\mu)$ has an $\alpha$-dimensional volume, hence it is measure doubling. Thus $(X,\dist)$ is metric doubling. Since $(X,\dist)$ is also complete, by \cite[Exer.\,10.17]{Heinonen}, any closed ball in $X$ is totally bounded and hence compact.
\end{proof}

Let us conclude with an important lemma.

\begin{lemma}
Let $(X,\dist,\mu,\cE)$ be with an $\alpha$-dimensional Euclidean heat kernel such that $(X,\dist)$ is complete. Then $(X,\dist)$ is a geodesic space.
\end{lemma}

\begin{proof}
Let us begin with showing that any two points $x,y \in X$ admit a midpoint, \ie a point $m \in X$ such that
\begin{equation}\label{eq:middle}
\dist(x,m)=\dist(y,m)=\frac{\dist(x,y)}{2}.
\end{equation}
Set $F(z)=\dist^2(x,z)+\dist^2(z,y)$ for any $z \in X$. Since $F(z) \to +\infty$ when $\dist(x,z), \dist(z,y) \to +\infty$, then there exists a ball $B \subset X$ such that $\inf_X F = \inf_B F$. From the previous lemma, we know that balls in $X$ are compact, so $\inf_B F$ is attained in some $m \in B$. Therefore, setting
\[
\lambda:=F(m)=\dist^2(x,m)+\dist^2(m,y),
\]
we have
\[
\|e^{-\sfrac{F}{4}}\|_{L^\infty(X,\mu)} = e^{-\sfrac{\lambda}{4}}.
\]
By the Chapman-Kolmogorov identity \eqref{eq:ChapmanKolmogorov}, we have for any $t>0$
\[
\int_X e^{-\psfrac{\dist^2(x,z)+\dist^2(z,y)}{4t}}\di \mu(z) = e^{-\sfrac{\dist^2(x,y)}{8t}}(2\pi t)^{\alpha/2}
\]
which can be raised to the power $t$ to provide
\[
\|e^{-\sfrac{F}{4}}\|_{L^{1/t}(X,\mu)} = e^{\sfrac{\dist^2(x,y)}{8}}(2\pi t)^{\alpha t/2}.
\]
Letting $t$ tend to $0$, this yields to $e^{-\sfrac{\lambda}{4}}=e^{-\sfrac{\dist^2(x,y)}{8}}$ hence $\lambda=\sfrac{\dist^2(x,y)}{2}$, thus \begin{equation}\label{eq:lambda2} \dist^2(x,m)+\dist^2(m,y) = \frac{\dist^2(x,y)}{2}
\end{equation}
by definition of $\lambda$. Since for any $z \in X$,
\begin{align*}
\dist^2(x,z)+\dist^2(z,y) & = \frac{1}{2}(\dist(x,z)+\dist(z,y))^2 + \frac{1}{2}(\dist(x,z)-\dist(z,y))^2\\
&\ge \frac{1}{2}(\dist(x,y))^2 + \frac{1}{2}(\dist(x,z)-\dist(z,y))^2,
\end{align*}
taking $z=m$ and using \eqref{eq:lambda2} implies $\dist(x,m)=\dist(m,y)$.

The existence of midpoints implies that $(X,\dist)$ is a length space, see \cite[Th.\,2.4.16\,(1)]{BuragoBuragoIvanov}. Then the result follows from \cite[Th.\,2.5.23]{BuragoBuragoIvanov} and \cite[Th.\,2.5.9]{BuragoBuragoIvanov}.
\end{proof}
\begin{remark}\label{rem:geodesicplus}
The previous proof can be adapted to show that if a complete proper metric measure space $(X,\dist,\mu)$ can be endowed with a symmetric Dirichlet form $\cE$ admitting a heat kernel $p$ such that for any $x,y\in X$,
\[
\lim_{t\to 0+}-4t\log p(x,y,t)=\dist^2(x,y),
\]
where the convergence holds locally uniformly, then $(X,\dist)$ is geodesic.
\end{remark}

\subsection{Strong locality and regularity of the Dirichlet form}

Let us show now that having an $\alpha$-dimensional Euclidean heat kernel forces a Dirichlet form to satisfy several properties. We begin with the following.

\begin{lemma}\label{lem:lipinD}
Let $(X,\dist,\mu,\cE)$ be with an $\alpha$-dimensional Euclidean heat kernel. Then $\Lip_\rmc(X) \subset \cD(\cE)$.
\end{lemma}

\begin{proof}
Let $f \in \Lip_\rmc(X)$ be with support $K$. Thanks to \eqref{eq:Grigoryan}, we only need to show that
\[
F \sep (0,+\infty) \ni t \mto \frac{1}{2t} \iint_{X\times X} (f(x)-f(y))^2 \frac{1}{(4 \pi t)^{\alpha/2}}\, e^{-\sfrac{\dist^2(x,y)}{4t}} \di \mu(x) \di \mu(y)
\]
is a bounded function. Set $D(x,y)=f(x)-f(y)$ for any $(x,y) \in X \times X$. It is easily checked that $\supp(D) \subset (K \times X) \cup (X \times K)$, so for any $t>0$, using the symmetry in~$x$ and~$y$ of the integrand we get
\begin{align*}
F(t) & \le \frac{\Lip(f)^2}{t} \int_K\int_X \dist^2(x,y) \frac{1}{(4 \pi t)^{\alpha/2}} e^{-\sfrac{\dist^2(x,y)}{4t}} \di \mu(x) \di \mu(y)\\
& = \frac{4 \Lip(f)^2}{(4 \pi t)^{\alpha/2}} \int_K \int_X \frac{\dist^2(x,y)}{4t} e^{-\sfrac{\dist^2(x,y)}{4t}} \di \mu(x) \di \mu(y).
\end{align*}
For any measurable function $g:X\to [0,\infty)$ and any $C^1$ function $\phi: [0,\infty) \to[0,\infty)$ satisfying $\lim_{\lambda\to+\infty}\varphi(\lambda)=0$ and $\int_0^{+\infty} |\varphi'|(\lambda)\ \mu\left(\left\{g<\lambda\right\}\right)d\lambda<\infty$, writing $\varphi(g(x))=\int_{g(x)}^{+\infty} \varphi'(\lambda)d\lambda$ and applying Fubini's theorem leads to
\begin{equation}\label{eq:Fub}
\int_X \varphi(g(x))\di\mu(x)=-\int_0^{+\infty} \varphi'(\lambda)\ \mu\left(\left\{g<\lambda\right\}\right)\di\lambda.
\end{equation}
For any $y \in K$, using this fact with $g(x)=\dist^2(x,y)/(4t)$ and $\phi(\xi)=\xi e^{-\xi}$, we get
\[
F(t) \le \frac{4 \Lip(f)^2}{(4 \pi t)^{\alpha/2}} \int_K \int_0^{+\infty} (\lambda - 1) e^{-\lambda} \mu(B_{\sqrt{4 t \lambda}}(y)) \di \lambda \di \mu(y).
\]
Setting $C_o=C_o(\alpha):=4 \int_0^{+\infty}(\lambda-1)e^{-\lambda}\lambda^{\alpha/2} \di \lambda$ and recalling that $\mu(B_{\sqrt{4 t \lambda}}(y))= \omega_{\alpha} (4 t \lambda)^{\alpha/2}$, we obtain $F(t) \le \Lip(f)^2 \mu(K) C_o \omega_{\alpha} \pi^{-\alpha/2}$, thus $F$ is bounded.
\end{proof}

Next we prove the following crucial result.

\begin{proposition}\label{prop:slandreg}
Let $(X,\dist,\mu,\cE)$ be with an $\alpha$-dimensional Euclidean heat kernel. Then $\cE$ is strongly local and regular.
\end{proposition}

\begin{proof}
In the proof of \cite[Th.\,4.1]{Grigor'yan}, it is shown that if $(X,\dist,\mu,\cE)$ admits a stochastically complete heat kernel $p$ satisfying
\[
t^{-\gamma/\beta} \Phi_1(\dist(x,y)t^{-1/\beta})\le p(x,y,t)\le t^{-\gamma/\beta} \Phi_2(\dist(x,y)t^{-1/\beta})
\]
for any $x,y \in X$ and $t>0$, where $\beta$ and $\gamma$ are positive constants and $\Phi_1, \Phi_2$ are monotone decreasing functions from $[0,+\infty)$ to itself such that $\Phi_1>0$ and $\int^{+\infty}s^{\beta + \gamma-1}\Phi_2(s)\di s<+\infty$, then
\begin{equation}\label{eq:Grigor'yan}
\cE(f) \simeq \limsup_{t\to 0} t^{-\sfrac{(\beta + \gamma)}{2}} \iint_{\{\dist(x,y)<t^{1/2}\}}(f(x)-f(y))^2\di \mu(x)\di \mu(y)
\end{equation}
holds for all $f \in \cD(\cE)$, what in turn implies strong locality of $\cE$. Here we have used $A\simeq B$ to denote the existence of a constant $c>1$ such that $c^{-1}A\le B \le cA$. Choosing $\Phi_1(s)=\Phi_2(s)=e^{-s^2/4}$, $\beta=2$ and $\gamma = \alpha$, we can apply this result in our context to get strong locality of $\cE$.

To prove regularity, let us show that $\Lip_\rmc(X)$ is a core for $\cE$.

\subsubsection*{Density in $C_\rmc(X)$} Let $f \in C_\rmc(X)$ be with support $K$. For any $R>0$ and $x \in X$, set $f_R(x)=\inf_y \{f(y)+R\dist(x,y)\}$. Note that $f_R(x)\le f(x)$ for any $x \in X$, and since $f(y)+R\dist(x,y)\to +\infty$ when $\dist(x,y) \to +\infty$ and $(X,\dist)$ is proper, the infimum in the definition of $f_R$ is always attained at some $x' \in X$. Then $f_R(x)=f(x')+R\dist(x,x')$ implies
\[
\dist(x,x') \le \frac{2\|f\|_{\infty}}{R}.
\]
Being continuous with compact support, $f$ is uniformly continuous, so it admits a modulus of continuity $\omega$ which we can assume non-decreasing with no loss of generality.
Then for any $x \in X$,
\[
|f_R(x)-f(x)|=f(x)-f_R(x)=f(x)-f(x')+\underbrace{f(x')-f_R(x)}_{=-R\dist(x',x)\le 0} \le \omega(\dist(x,x')),
\]
so that
\[
\|f_R-f\|_{\infty} \le \omega\left(2 \|f\|_{\infty}/R\right) \to 0
\]
when $R \to +\infty$. Therefore, setting $\phi_K:=\max(1-\dist(\cdot,K),0)$ and $g_R:=\phi_K f_R$ for any $R>0$, we get a sequence of compactly supported Lipschitz functions $(g_R)_R$ converging uniformly to $f$.

\subsubsection*{Density in $\cD(\cE)$} Let $\Lip_o(X)$ be the set of Lipschitz functions $f$ on $X$ vanishing at infinity, \ie such that for some $o \in X$ one has $f(x)\to 0$ when $\dist(o,x)\to+\infty$. We are going to show that $\Lip_\rmc(X)$ is dense in $\Lip_o(X)\cap \cD(\cE)$ for the norm $|\cdot|_{\cE}$. By \eqref{eq:Grigor'yan}, we know that that there exists a constant $C_\alpha>0$ such that if $f\in \cD(\cE)$, then
\[
\frac{1}{C_\alpha}\limsup_{t\to 0}\int_X E(f,x,t)\di\mu(x)\le \cE(f) \le C_\alpha\limsup_{t\to 0}\int_X E(f,x,t)\di\mu(x),
\]
where
\[
E(f,x,t)= t^{-\sfrac{(\alpha+2)}{2}} \int_{B_{\sqrt{t}}(x)}(f(x)-f(y))^2\di \mu(y).
\]

Let $f\in \Lip_o(X)\cap \cD(\cE)$. For any $R>0$, we set
\[
\varphi_R(x):=\Bigl(1-\frac{\dist(x,B_R(o)}{R}\Bigr)_+
\]for any $x \in X$, and $f_R:= f \varphi_R$.
By monotone convergence, we have
\[
\lim_{R\to +\infty} \|f-f_R\|_2=0.
\]
We look now at $E(f_R,x,t)$ and we distinguish 3 cases:
\begin{itemize}
\item if $x\in B_{R-\sqrt{t}}(o)$, then $E(f_R,x,t)=E(f,x,t)$;
\item if $x\not\in B_{2R+\sqrt{t}}(o)$, then $E(f_R,x,t)=0$;
\item if $x\in B_{2R+\sqrt{t}}(o)\setminus B_{R-\sqrt{t}}(o)$,
\end{itemize}
then using $f_R(x)-f_R(y)=\varphi_R(x)(f(x)-f(y))+f(y)(\varphi_R(x)-\varphi_R(y))$ we get
\begin{align*}
E(f_R,x,t)&\le 2 t^{-\sfrac{(\alpha+2)}{2}} \biggl(\int_{B_{\sqrt{t}}(x)}(f(x)-f(y))^2\di \mu(y)\\[-5pt]
&\hspace*{4cm}+\int_{B_{\sqrt{t}}(x)}f^2(y)(\varphi_R(x)-\varphi_R(y))^2\di\mu(y)\biggr)\\
&\le 2E(f,x,t)+ 2 t^{-\sfrac{(\alpha+2)}{2}} \int_{B_{\sqrt{t}}(x)}f^2(y)\frac{t}{R^2}\di\mu(y),
\end{align*}
where we have used the fact that $\varphi_R$ is $1/R$-Lipschitz. By Fubini's theorem,
\[
\int_X \int_{B_{\sqrt{t}}(x)}f^2(y)\di \mu(y) \di \mu(x) = \int_X f^2(y)\mu(B_{\sqrt{t}}(y))\di \mu(y) = \omega_\alpha t^{\sfrac{\alpha}{2}} \int_X f^2(y)\di \mu(y),
\]
thus
\begin{align*}\cE(f_R)&\le C_\alpha \limsup_{t\to 0}\int_{X} E(f,x,t)\di\mu(x)+C_\alpha \frac{2}{R^2}\, \omega_\alpha \int_X f^2(y)\di\mu(y)\\
&\le C_\alpha^2 \cE(f)+C_\alpha \frac{2}{R^2}\, \omega_\alpha \int_X f^2(y)\di\mu(y).
\end{align*}
Hence $\{f_R\}_{R\ge 1}$ is bounded in $\cD(\cE)$. The fact that $\lim_{R\to +\infty} \|f-f_R\|_2=0$ implies that $f_R$ converges weakly to $f$ in $\cD(\cE)$ when $R\to+\infty$. By Mazur's lemma, there exists a sequence $\{u_\ell\}_\ell \subset \Lip_\rmc(X)$ made of convex combination of $\{f_R\}_{R\ge 1}$ such that
\[
\lim_{\ell\to +\infty} \|f-u_\ell\|_\cE=0.
\]
Therefore, it is enough to show that $\Lip_o(X)$ contains a subset that is dense in $\cD(\cE)$ for $|\cdot|_\cE$. Let $L^2_\rmc$ be the set of compactly supported functions $f$ in $L^2(X,\mu)$. Then $P_t(L^2_\rmc) \subset \Lip_o(X)$ for any fixed $t>0$. Indeed, for any $f \in L^2_\rmc$ and $x, y \in X$,
\[
|P_tf(x) - P_tf(y)|\le \frac{1}{(4\pi t)^{\alpha/2}} \int_X \Big|e^{-\sfrac{\dist^2(x,z)}{4t}}-e^{-\sfrac{\dist^2(y,z)}{4t}}\Big| |f(z)| \di \mu(z).
\]
Setting $\phi(s)=e^{-\sfrac{s^2}{4t}}$ and noting that $|\phi'(s)|\le |\phi'(\sqrt{2t})|=:c_t$ for all $s>0$, we get from the mean value theorem, the triangle inequality and Hölder's inequality, that $|P_tf(x) - P_tf(y)|\le C(t,f) \dist(x,y)$ with $C(t,f):=c_t(4 \pi t)^{-\alpha/2} \mu(\supp f)^{1/2} \| f \|_{L^2}$. Moreover,
\begin{align*}
|P_tf(x)| &= \left|\frac{1}{(4\pi t)^{\alpha/2}} \int_{\supp f} f(y) e^{-\sfrac{\dist^2(x,y)}{4t}} \di \mu(y)\right|\\
&\le \frac{ e^{-\sfrac{\dist^2(o,x)}{8t}} }{ (4\pi t)^{\alpha/2} } \int_{\supp f} |f(y)| e^{\sfrac{\dist^2(o,y)}{4t}} \di \mu(y) \to 0
\end{align*}
when $\dist(o,x) \to +\infty$.

Let us show now that $P_t(L^2_\rmc)$ is dense in $\cD(\cE)$ by proving that its $\langle \cdot, \cdot \rangle_{\cE}$-orthogonal complement $F$ in $\cD(\cE)$ reduces to $\{0\}$. For any $v\in F$, we have:
\begin{equation}\label{eq:18.07}
\int_X v P_t f \di \mu + \cE(v,P_t f)=0 \qquad \forall f \in L^2_\rmc.
\end{equation}
Since $P_t$ maps $L^2(X,\mu)$ into $\cD(L)$ and is self-adjoint, then
\begin{align*}
\cE(v,P_t f) & = - \int_X v L (P_t f) \di \mu = -\int_X v \frac{\di}{\di t} P_t f \di \mu = -\frac{\di}{\di t} \int_X v P_t f \di \mu\\
& = - \frac{\di}{\di t} \int_X (P_tv) f \di \mu
= -\int_X \frac{\di}{\di t}(P_tv) f \di \mu = - \int_X L(P_tv) f \di \mu = \cE(P_tv,f)
\end{align*}
for any $f \in L^2_\rmc$, so \eqref{eq:18.07} becomes:
\[
\int_X (P_tv) f \di \mu + \cE(P_t v,f)=0 \qquad \forall f \in L^2_\rmc.
\]
This implies $P_t v \in \cD(L)$ with $L (P_t v) = P_t v$. Since $L$ is a non-positive operator, $1$~cannot be an eigenvalue of $L$, so we necessarily have $P_t v = 0$. This implies $v=0$ since the spectral theorem ensures that $0$ cannot be an eigenvalue of $P_t$.
\end{proof}

Proposition \ref{prop:slandreg} has several consequences. The first one is the existence of a $\Gamma$ operator defined on $\cD_{\loc}(\cE)$ for any $(X,\dist,\mu,\cE)$ with an $\alpha$-dimensional Euclidean heat kernel. Coupled with Lemma \ref{lem:lipinD}, this yields the following.

\begin{corollary}\label{cor:1}
Let $(X,\dist,\mu,\cE)$ be with an $\alpha$-dimensional Euclidean heat kernel. Then $\Lip(X) \subset \cD_{\loc}(\cE)\cap C(X)$ and for some constant $C_1$ depending only on $\alpha$, we have:
\begin{equation}\label{eq:Lip}
\Gamma(f) \le C_1 \Lip(f)^2 \mu \qquad \forall f \in \Lip(X).
\end{equation}
\end{corollary}

\begin{proof}
Take $f \in \Lip(X)$. For any compact set $K \subset X$, the function $\phi_K:=\max(1-\dist(\cdot,K),0)$ is a compactly supported Lipschitz function constantly equal to $1$ on $K$. Therefore, $f \phi_K$ coincides with $f$ on $K$, and thanks to Lemma \ref{lem:lipinD}, $f \phi_K$ belongs to $\cD(\cE)$. This shows that $f \in \cD_{\loc}(\cE)$. Moreover, for any non-negative $\phi \in C_\rmc(X)$ and $t>0$, a direct computation like in the proof of the previous lemma implies
\[
\frac{1}{2t} \iint_{X\times X} \phi(x)(f(x)-f(y))^2 p(x,y,t) \di \mu(x) \di \mu(y) \le C_1 \Lip(f)^2 \int_X \phi(x) \di \mu(x)
\]
with $C_1$ depending only on $\alpha$, so that letting $t$ tend to $0$ and applying formula \eqref{eq:Grigor'yan2} yields to \eqref{eq:Lip}.
\end{proof}

Proposition \ref{prop:slandreg} also implies that we can define the pseudo-distance $\dist_\cE$ recalled in Section \ref{sec:2} for any $(X,\dist,\mu,\cE)$ with an $\alpha$-dimensional Euclidean heat kernel. It turns out that in this case, $\dist_\cE$ is equivalent to the initial distance $\dist$, as shown in the next proposition.

\begin{proposition}\label{eq:propA}
Let $(X,\dist,\mu,\cE)$ be with an $\alpha$-dimensional Euclidean heat kernel. Then there exists $C_2>0$ depending only on $\alpha$ such that $C_2 \dist \le \dist_{\cE} \le \dist$. In particular, the assumption \eqref{eq:A} is satisfied.
\end{proposition}

\begin{proof}
Let us first show that $C_2\dist \le \dist_\cE$ for some $C_2>0$. Set
\[
\Lambda:=\{f \in \Lip(X) \sep C_1\Lip(f)^2 \le 1\},
\]
where $C_1$ is as in \eqref{eq:Lip}. It follows from Corollary \ref{cor:1} that $\Lambda$ is included in the set of test functions in \eqref{eq:defdist}. Noting that $f:=C_1^{-1/2}\dist(x,\cdot)$ is in $\Lambda$ for all $x \in X$ and that $|f(x)-f(y)|=C_1^{-1/2}\dist(x,y)$ for all $x,y \in X$, with $C_2:=C_1^{-1/2}$ we get:
\[
\dist_\cE(x,y)\ge C_2\dist(x,y) \qquad \forall x,y \in X.
\]

Let us show now that $\dist_\cE \le \dist$. To this aim, we follow the lines of \cite{Grigor'yanEd}. Let $v \in \cD_{\loc}(\cE)\cap C(X)$ be bounded and such that $\Gamma(v)\le \mu$. For any $a \in \setR$, $t>0$ and $x \in X$, set $\xi_a(x,t):=a v(x) - \sfrac{a^2t}{2}$.

\begin{claim}
For any $f \in L^2(X,\mu)$, the quantity
\[
I(t):=\int_X f_t^2 e^{\xi_a(\cdot,t)} \di \mu,
\]
where $f_t:=P_t f$, does not increase when $t>0$ increases.
\end{claim}

Indeed, for any $t>0$, writing $\xi_a$ for $\xi_a(\cdot,t)$ and $\xi_a'$ for $\frac{\di}{\di t} \xi_a(\cdot,t)$, we have
\begin{align*}
\frac{\di}{\di t} (f_t^2e^{\xi_a}) & = 2 f_t \Bigl(\frac{\di}{\di t} f_t \Bigr)e^{\xi_a} + f_t^2\xi_a' e^{\xi_a} = 2 f_t Lf_t e^{\xi_a}-\frac{a^2}{2}f_t^2 e^{\xi_a}.
\end{align*}
Since $e^{\xi_a} \le e^{a \|v\|_\infty}$, this implies
\[
\left|\frac{\di}{\di t} (f_t^2e^{\xi_a}) \right| \le e^{a \|v\|_\infty}(2|f_t||Lf_t| + a^2|f_t|^2/2) \, \in \, L^1(X,\mu),
\]
so we can differentiate under the integral sign to get
\[
I'(t) = 2 \int_X f_t Lf_t e^{\xi_a} \di \mu - \frac{a^2}{2} \int_X f_t^2 e^{\xi_a} \di \mu.
\]
The Leibniz rule \eqref{eq:Leibniz} implies
\begin{align*}
\int_X f_t Lf_t e^{\xi_a} \di \mu = & \, \, - \cE(f_t,f_te^{\xi_a}) = - \int_X \Gamma(f_t,f_t e^{\xi_a}) \di \mu \\ = & - \int_X f_t \underbrace{\Gamma(f_t,e^{\xi_a})}_{=\Gamma(f_t,e^{\xi_a/2}e^{\xi_a/2})} \di \mu - \int_X e^{\xi_a} \Gamma(f_t) \di \mu\\
= & - 2 \int_X f_t e^{\xi_a/2} \Gamma(f_t,e^{\xi_a/2})\di \mu - \int_X e^{\xi_a} \Gamma(f_t) \di \mu
\end{align*}
and, starting from $\Gamma(f_te^{\xi_a/2})$,
\[
- 2 f_t e^{\xi_a/2} \Gamma(f_t,e^{\xi_a/2}) = - \Gamma(f_te^{\xi_a/2}) + f_t^2 \Gamma(e^{\xi_a/2}) + e^{\xi_a} \Gamma(f_t),
\]
so that
\begin{align*}
I'(t) & = 2 \int_X f_t^2 \Gamma(e^{\xi_a/2}) - 2 \int_X \Gamma(f_t e^{\xi_a/2})\di \mu - \frac{a^2}{2}\int_X f_t^2 e^{\xi_a} \di \mu \\ & \le \int_X f_t^2 \Bigl(2 \Gamma(e^{\xi_a/2})-\frac{a^2}{2} e^{\xi_a} \Bigr) \di \mu.
\end{align*}
Since $v$ is bounded, we can apply the chain rule \eqref{eq:chain} with $\eta(\xi):=e^{\xi/2}$ to get $\Gamma(e^{\xi_a/2})=(1/4)e^{\xi_a}\Gamma(\xi_a)$ and thus
\[
2 \Gamma(e^{\xi_a/2}) - \frac{a^2}{2}e^{\xi_a} = \frac{1}{2}e^{\xi_a}\Gamma(\xi_a) - \frac{a^2}{2}e^{\xi_a} = \frac{1}{2}e^{\xi_a} a^2 \Gamma(v)-\frac{a^2}{2}e^{\xi_a}\le 0,
\]
so $I'(t) \le 0$.

From now on, assume $a \ge 0$. Apply the claim to $f=1_A$, where $A$ is some Borel subset of $X$. Then for any $t>0$ and any Borel subset $B$ of $X$,
\[
\int_B f_t^2 e^{\xi_a(\cdot,t)} \di \mu \le \int_X f_t^2 e^{\xi_a(\cdot,t)} \di \mu =I(t)\le I(0) = \int_{A} e^{a v} \di \mu,
\]
hence
\[
\int_B f_t^2 e^{\xi_a(\cdot,t)} \di \mu \le \mu(A) e^{a \sup_A v}.
\]
Moreover, since the heat kernel is Euclidean, we have
\begin{align*}
\int_B f_t^2 e^{\xi_a(\cdot,t)} \di \mu & = \int_{B} \left(\int_A p(x,y,t) \di \mu(y)\right)^2 e^{\xi_a(x,t)} \di \mu(x)\\ & \ge \frac{e^{-\psfrac{\sup_{A\times B} \dist^2}{2t}}}{(4\pi t)^\alpha}\, \mu(B) \mu(A)^2 e^{a \inf_B v - a^2 \sfrac{t}{2}},
\end{align*}
thus
\[
\frac{e^{-\psfrac{\sup_{A\times B} \dist^2}{2t}}}{(4\pi t)^\alpha}\, \mu(B) \mu(A) e^{a (\inf_B v - \sup_A v) - a^2 \sfrac{t}{2}} \le 1.
\]

Take $x,y \in X$. With no loss of generality we can assume $v(y)-v(x)>0$. Choose~$t$ such that $\sqrt{t}<\dist(x,y)/3$. Apply the previous inequality with $A=B_{\sqrt{t}}(x)$ and $B=B_{\sqrt{t}}(y)$. In this case, $\sup_{A\times B} \dist^2 = \dist^2(x,y) + 2 \sqrt{t}$. Moreover, since $v$ is continuous, we have $\inf_B v - \sup_A v = v(y)-v(x) + \eps(t)$ where $\eps(t)\to 0$ when $t \to 0^+$. Then
\[
\frac{e^{-\psfrac{\dist^2(x,y) + 2\sqrt{t}}{2t}}}{(4\pi)^\alpha}\, \omega_\alpha^2\, e^{a (v(y)-v(x) + \eps(t)) - a^2 \sfrac{t}{2}} \le 1
\]
for any $t \in (0,\dist^2(x,y)/9)$ and any $a \ge 0$. Now for $t \in (0,\dist^2(x,y)/9)$, choose $a=a(t)=(v(y)-v(x) + \eps(t))/t$ to get
\[
\frac{e^{-\psfrac{\dist^2(x,y) + 2\sqrt{t}}{2t}}}{(4\pi)^\alpha}\, \omega_\alpha^2\, e^{\Psfrac{1}{2t}(v(y)-v(x) + \eps(t))^2} \le 1.
\]
Apply the logarithm function, multiply the resulting inequality by $2t$ and then add $\dist^2(x,y)$ to get
\[
-2\sqrt{t} + 2t \ln(\omega_\alpha^2/(4\pi)^\alpha) + (v(y)-v(x) + \eps(t))^2 \le \dist^2(x,y).
\]
Letting $t$ tend to $0$ gives
\begin{equation}\label{16sept2}
(v(x)-v(y))^2 \le \dist^2(x,y).
\end{equation}

Since for any $u \in \cD_{\loc}(\cE)$ and any $R>0$, the function $u_R:=\max(u,R)$ is in $\cD_{\loc}(\cE)$ with $\Gamma(u_R)\le \Gamma(u)$, approximating any possibly unbounded $v \in \cD_{\loc}(\cE)\cap C_\rmc(X)$ with $\Gamma(v) \le \mu$ by $(v_R)_{R>0}$ provides \eqref{16sept2} for any $x,y \in X$ for such a $v$. This implies $\dist_\cE \le \dist$.
\end{proof}

\begin{remark}
Though we will not use it in the sequel, let us point out that Proposition \ref{eq:propA} can be upgraded into $\dist=\dist_\cE$ provided a suitable technical assumption holds: see \cite[Th.\,2.5\,(I)]{terElstRobinsonSikora}.
\end{remark}

\subsection{Evaluation of $L$ on squared distance functions}

Let us show now that the operator associated to the Dirichlet form of a space with an $\alpha$-dimensional Euclidean heat kernel behaves on squared distance functions as the Laplacian does in $\setR^n$.

\begin{lemma}\label{lem:Ld^2}
Let $(X,\dist,\mu,\cE)$ be with an $\alpha$-dimensional Euclidean heat kernel. Then $L\dist^2(x,\cdot)=2\alpha$, $L\dist(x,\cdot)=(\alpha-1)/\dist(x,\cdot)$ on $X\backslash \{x\}$ and $\Gamma(\dist(x,\cdot))=1$ $\mu$-a.e.\ on $X$.
\end{lemma}

\begin{proof}
Take $x\in X$. Note first that Corollary \ref{cor:1} guarantees that $\dist^2(x,\cdot), \dist(x,\cdot) \in \cD_{\loc}(\cE)$. For any $t>0$, a direct computation relying on the chain rule \eqref{eq:chainrule} and starting from the equation
\[
\Bigl(\frac{\di}{\di t} - L \Bigr) \frac{e^{-\sfrac{\dist^2(x,\cdot)}{4t}}}{(4 \pi t)^{\alpha/2}} = 0
\]
(recall the remark after Definition 2.3) provides
\[
\frac{\dist^2(x,\cdot)}{4t^2} + \frac{\alpha}{2t} - \frac{1}{4t}L\dist^2(x,\cdot) - \frac{1}{(4t)^2}\Gamma(\dist^2(x,\cdot)) =0.
\]
Multiplying by $(4t)^2$ and letting $t$ tend to $0$ gives $\Gamma(\dist^2(x,\cdot))=4\dist^2(x,\cdot)$ hence $\Gamma(\dist(x,\cdot))=1$ by \eqref{eq:chain}, while multiplying by $4t$ and letting $t$ tend to $+\infty$ gives $L\dist^2(x,\cdot)=2\alpha$, from which follows $L\dist(x,\cdot)=(\alpha-1)/\dist(x,\cdot)$ by \eqref{eq:chainrulesquare}.
\end{proof}

\begin{remark}\label{rk:applypreparatory}
Note that any space $(X,\dist,\mu,\cE)$ with an $\alpha$-dimensional Euclidean heat kernel satisfies the assumptions of Lemma \ref{lem:preparatory}. Indeed, Proposition \ref{prop:slandreg} states that $(X,\dist,\mu,\cE)$ is strongly local and regular, Corollary \ref{cor:1} yields assumption \eqref{lem:preparatory1} and Lemma \ref{lem:Ld^2} yields assumption \eqref{lem:preparatory2}.
\end{remark}

As a consequence of Lemma \ref{lem:Ld^2}, we can show that locally $L$-harmonic functions on spaces with an $\alpha$-dimensional Euclidean heat kernel are necessarily strongly harmonic.

\begin{lemma}\label{lem:Lstrong} Let $(X,\dist,\mu,\cE)$ be with an $\alpha$-dimensional Euclidean heat kernel. Let $\Omega \subset X$ be an open set and $h$ a locally integrable local solution of $Lh=0$ on $\Omega$. Then for any $x \in \Omega$, the function defined on $(0,\dist(x,^\rmc\!\Omega))$ by
\[
r\mto \fint_{B_r(x)}h \di \mu
\]
is a constant. Therefore, $h$ has a continuous representative strongly harmonic in $\Omega$.
\end{lemma}
\begin{proof}
Take $x \in \Omega$ and set $R:=\dist(x,^\rmc\!\Omega)$. From $Lh=0$ we get that for any $\phi \in \cD(L)$ with compact support in $\Omega$,
\[
\langle h,L\phi\rangle_{L^2}=0.
\]
Take $\phi \in C_\rmc^\infty((0,R))$ and set
$u=\phi \circ \dist(x,\cdot)$ on $X$. Then $u$ belongs to $\cD(L)$ and has compact support included in $\Omega$. The chain rule \eqref{eq:chainrule} and Lemma \ref{lem:Ld^2} yields to
\[
Lu = \chi\circ \dist(x,\cdot),
\]
where we have set
\[
\chi(r):= \phi''(r) + \frac{\alpha - 1}{r} \phi'(r) = r^{1-\alpha} \left(r^{\alpha-1}\, \phi'\right)'
\]
for any $r \in (0,R)$. Since $\chi \circ \dist(x,\cdot)=-\int_{\dist(x,\cdot)}^R \chi'(s)ds$, we have
\begin{align*}
-\int_0^R \chi'(s)\biggl(\int_{B_s(x)}h\di\mu\biggr)\di s & = -\int_0^R \chi'(s)\biggl(\int_X h1_{B_s(x)}\di\mu\biggr)\di s\\
& = \int_X \biggl(- \int_0^R \chi'(s) 1_{(\dist(x,\cdot),+\infty)}(s) \di s \biggr) h \di \mu\\
& = \int_X (\chi \circ \dist(x,\cdot)) h \di \mu = \langle h,Lu\rangle_{L^2} = 0.
\end{align*}
This implies that the function
$s\mto I(s):= \int_{B_s(x)}h\di \mu$ satisfies the equation
\[
\left[r^{\alpha-1}\left(r^{1-\alpha}y'\right)'\right]'=0
\]
in the distributional sense on $(0,R)$. Then there exists real-valued constants $a$, $b$ and~$c$ such that for any $s \in (0,R)$,
\[
I(s)=as^\alpha+bs^2+c.
\]
Since $h$ is locally integrable, $c=\lim_{s\to 0} I(s) = 0$. Then $s^{-\alpha}I(s)= a+b s^{2-\alpha}$ for any $s \in (0,R)$. Using test functions $\phi$ that are constantly equal to $1$ in a neighborhood of~$0$, we can get
$(2-\alpha)b=0$, from which follows that $s \mto s^{-\alpha} I(s)$ is a constant. Since $(X,\dist,\mu)$ has an $\alpha$-dimensional volume, this provides the result.
\end{proof}

\subsection{Spaces of locally $L$-harmonic functions with polynomial growth}

Let us conclude with a result that shall be crucial in the next section. We recall that for any positive integer $m$, a function $h:X \to \setR$ has polynomial growth of rate $m$ if there exists $C>0$ such that $|h|\le C(1+\dist^m(o,\cdot))$ holds for some $o \in X$. The case $m=1$ corresponds to a linear growth. Note that functions with a fixed polynomial growth rate form a vector space.

\begin{proposition}\label{prop:finitedim}
Let $(X,\dist,\mu,\cE)$ be with an $\alpha$-dimensional Euclidean heat kernel. Then for any $m \in \setN\backslash \{0\}$, the space of locally $L$-harmonic functions $h:X\to\setR$ with polynomial growth of rate $m$ is finite dimensional.
\end{proposition}

\begin{proof}
Having an $\alpha$-dimensional Euclidean heat kernel, $(X,\dist,\mu,\cE)$ trivially satisfies the Gaussian estimate \eqref{eq:LiYau}. Moreover, we know from Proposition \ref{eq:propA} that it satisfies the assumption \eqref{eq:A}. Consequently, by Proposition \ref{prop:important}, $(X,\dist,\mu,\cE)$ has the doubling and Poincaré properties. Therefore, the arguments of \cite{ColdingMinicozzi} carry over.
\end{proof}

\section{Construction of the isometry}\label{sec:4}
\Changel
In this section, we construct an isometry between a given metric measure space $(X,\dist,\mu)$ equipped with a Dirichlet form $\cE$ admitting an $\alpha$-dimensional Euclidean heat kernel and an Euclidean space $\setR^l$ equipped with a distance $\dist_Q$ associated to a suitable quadratic form $Q$.

Let us recall that a quadratic form on a $\setR$-vector space $V$ is a map $Q: V \to \setR$ for which there exists a bilinear symmetric form $\beta: V \times V \to\setR$ such that $Q(u)=\beta(u,u)$ for any $u \in V$, in which case one has
$\beta(u,v)=\frac{1}{2}(Q(u+v)-Q(u)-Q(v))$ and then
$Q(u+v)=Q(u) +2\beta(u,v) + Q(v)$ for any $u,v \in V$. Moreover, when $Q$ is positive definite, setting\begin{equation}\label{eq:dist_Q}\dist_Q(u,v):=\sqrt{Q(u-v)}\end{equation} for any $u,v \in V$ defines a distance on $V$ canonically associated to $Q$. When $V=\setR^l$, Sylvester's law of inertia states that $Q$ can be transformed into $(v_1,\ldots,v_l) \mto \sum_i v_i^2$ via a suitable change of basis. This implies that $(\setR^l,\dist_Q)$ and $(\setR^l,\dist_e)$ are isometric, so that the construction made in this section proves Theorem \ref{th:main}.

\subsection{The quadratic form $Q$ and the coordinate function $H$}\label{subsec:4.1}

Let us explain how to define $Q$ in our context. We first fix a base point $o \in X$ and set
\[
B(x,y):=\frac{1}{2}(\dist^2(o,x)+\dist^2(o,y)-\dist^2(x,y))
\]for any $x,y \in X$. Note that
\begin{equation}\label{eq:D6.1}
\dist^2(x,y)= B(x,x)+B(y,y)-2B(x,y).
\end{equation}

\begin{remark}
In case $(X,\dist)=(\setR^l,\dist_{e})$ and $o$ is the origin in $\setR^l$, the law of cosines gives $B(x,y)=\scal{x}{y}$ for any $x,y \in \setR^l$, where $\scal{\cdot}{\cdot}$ is the usual Euclidean scalar product in $\setR^l$.
\end{remark}

For any $x \in X$, it follows from Lemma \ref{lem:Ld^2} and the fact that constant functions are locally $L$-harmonic that $B(x,\cdot)$ is locally $L$-harmonic. Moreover, for any $x,y \in X$, since $\dist^2(o,y)-\dist^2(x,y) = (\dist(o,y)-\dist(x,y))(\dist(o,y)+\dist(x,y))$, $\dist(o,y)-\dist(x,y)\le \dist(o,x)$ and $\dist(x,y)\le \dist(x,o)+\dist(o,y)$, we have
\begin{align*}
B(x,y) & \le \frac{1}{2}(\dist^2(o,x)+\dist(o,x)(\dist(o,x)+2\dist(o,y)))\\
& = \dist^2(o,x)+\dist(o,x)\dist(o,y) \le C_x(1+\dist(o,y))
\end{align*}
with $C_x:=\max(\dist^2(o,x),\dist(o,x))$. This shows that $B(x,\cdot)$ has linear growth for any $x \in X$. Then $\cV:=\Span\{B(x,\cdot): x \in X\}$ is a subspace of the space of locally $L$\nobreakdash-harmonic functions with linear growth. Using Proposition \ref{prop:finitedim}, we know that this space has a finite dimension, so $\cV$ has a finite dimension which we denote by $l$.

Let us then consider the subspace $\mathcal{D}:=\Span\{\delta_x: x \in X\}$ of the algebraic dual~$\cV^*$ of $\cV$. If $f \in \cV$ is such that $\theta(f)=0$ for any $\theta \in \cD$, then $f=0$; since the duality pairing $\cV\times\cV^* \to\setR$ is non-degenerate, this implies $\cD = \cV^*$. Therefore, there exist $x_1,\dots, x_l \in X$ such that $\{\delta_{x_1},\dots,\delta_{x_l}\}$ is a basis of $\mathcal{V}^*$. Let $\{h_1,\dots,h_l\}$ be the associated basis of $\mathcal{V}$. Then for any $x \in X$,
\begin{equation}\label{eq:basis}
B(x,\cdot) = \sum_{i=1}^l \delta_{x_i}(B(x,\cdot)) h_i = \sum_{i=1}^l B(x,x_i) h_i
\end{equation}
and for any $i \in \{1,\dots,l\}$,
\[
B(x,x_i) = B(x_i,x) = \sum_{j=1}^l \delta_{x_j}(B(x_i,\cdot)) h_j(x)= \sum_{j=1}^l B(x_i,x_j)h_j(x).
\]
Therefore, we have
\begin{equation}\label{eq:D6.2}
B(x,y) = \sum_{i,j=1}^l B(x_i,x_j)h_j(x)h_i(y)
\end{equation}
for any $x,y \in X$. We now define $Q$ on $\setR^l$ by setting:
\[
Q(\xi):= \sum_{i,j=1}^lB(x_i,x_j)\xi_i \xi_j \qquad \forall \, \xi=(\xi_1,\dots,\xi_l) \in \setR^l.
\]
Then $Q$ is a quadratic form whose associated symmetric form $\beta$ is given by
\[
\beta(\xi,\xi') = \sum_{i,j=1}^l B(x_i,x_j)\xi_i \xi_j'
\]
for any $\xi, \xi' \in \setR^l$. Note that $\beta$ is non-degenerate. Indeed, if $\xi \in \setR^l $ is such that $\beta(\xi,\cdot)=0$, then for any $y \in X$ we have $\sum_{i=1}^l \xi_i B(x_i,y)=0$ because
\[
\sum_{i=1}^l \xi_i B(x_i,y)=\sum_{i,j=1}^l \xi_i h_j(y)B(x_i,x_j)=\beta(\xi,(h_1(y),\ldots,h_l(y)))=0.
\]
But $\{B(x_i,\cdot)\}_i$ is a basis of $\cV$ since
\[
B(x,y)=\sum_{i=1}^l \delta_{x_i}(B(\cdot,y))h_i(x)= \sum_{i=1}^l B(x_i,y)h_i(x)
\]
for any $x,y \in X$, thus $\xi=0_l$, where $0_l$ denotes the origin in $\setR^l$.

We are now in a position to introduce our ``coordinate'' function $H:X\to\setR^l$ which we define as
\[
H:=(h_1,\dots,h_l).
\]
This function $H$ is continuous because $h_1,\dots,h_l$ are so. Moreover, for any $x,y \in X$, we have
\begin{equation}\label{eq:D6.3}
\beta(H(x),H(y)) = B(x,y)
\end{equation}
thanks to \eqref{eq:D6.2} and
\begin{equation}\label{eq:D6.4}
\dist^2(x,y) = Q(H(x)-H(y))
\end{equation}
thanks to \eqref{eq:D6.1}. Note that $H(o)=0_l$ because $B(x,o)=0$ for any $x \in X$. Moreover, \eqref{eq:D6.4}, the continuity of $H$ and the completeness of $(X,\dist)$ imply that $H(X)$ is a closed set of $\setR^l$.

\begin{claim}\label{claim}
$H$ is an injective map. Moreover, $\Span(H(X)) = \setR^l$ -- in fact, the closed convex hull of $H(X)$ is $\setR^l$.
\end{claim}

\begin{proof}
If $H(x)=H(y)$ then \eqref{eq:D6.4} gives $\dist^2(x,y)=Q(0)=0$, so $x=y$. Then $H$ is injective. For the second statement, let us recall that the closed convex hull $\overline{\mathrm{conv}}(E)$ of a closed set $E \subset \setR^l$ is defined as the smallest convex subset of $\setR^l$ containing $E$; moreover, $\overline{\mathrm{conv}}(E)$ coincides with
\[
\bigcap_{\lambda \in \mathcal{A}(E)} \{\lambda \ge 0\},
\]
where $\mathcal{A}(E)$ is the set of affine functions on $\setR^l$ that are non-negative on $E$. Note that being closed and convex, $\Span(E)$ contains $\overline{\mathrm{conv}}(E)$. Take $\lambda \in \mathcal{A}(H(X))$. Then $\lambda \circ H: X \to \setR$ is an affine combination of locally $L$-harmonic functions, hence it is a locally $L$-harmonic function too. Since $\lambda \circ H$ is non-negative on $X$, Lemma \ref{lem:ellipticHarnack} implies that it is a constant. Therefore, $\overline{\mathrm{conv}}(H(X))=\setR^l$, what brings the result.
\end{proof}

Note that the right-hand side in \eqref{eq:D6.4} does not define any squared distance on $\setR^l$ unless $Q$ is shown to be positive definite, see Section \ref{subsec:4.3}.

\subsection{Conical structure of tangent cones at infinity}

Let $(\uX,\udist,\uo)$ be a tangent cone at infinity of $(X,\dist)$ at $o$. We denote by $\{(X_i,\dist_i:=r_i^{-1}\dist,o)\}_i$, where $\{r_i\}_i \subset (0,+\infty)$ converges to $+\infty$, the sequence of rescalings of $(X,\dist,o)$ converging in the pointed Gromov-Hausdorff topology to $(\uX,\udist,\uo)$. Note that whenever $\ux \in \uX$, there exists a sequence $\{x_i\}_i\subset X$ such that $x_i \To{\sGH} \ux$; in particular, $\dist_i(o,x_i) \to \udist(\uo,\ux)$ and $\dist(o,x_i)\to +\infty$.

\Subsubsection*{Step 1. Construction of a Busemann function $h_\infty$ associated with a divergent sequence}

Let $\{x_i\}_i \subset X$ be a sequence such that $\dist(o,x_i) \to +\infty$. For any $i$, setting
\[
D_i:=\dist(o,x_i),
\]
we define
\[
h_i(y):=D_i-\dist(x_i,y)
\]
for any $y \in X$ and call $c_i:[0,D_i] \to X$ a minimizing geodesic joining $o$ to $x_i$.

On one hand, the triangle inequality implies that the functions $h_i$ are all $1$-Lipschitz, so by the Ascoli--Arzelà theorem, up to extracting a subsequence, we can assume that the sequence $\{h_i\}_i$ converges uniformly on compact subsets of $X$ to a $1$-Lipschitz function $h_\infty$. On the other hand, the minimizing geodesics $c_i$ being $1$-Lipschitz too, we can use again the Ascoli-Arzelà theorem to assume, up to extraction, that they converge uniformly on compact sets of $[0,+\infty)$ to a geodesic ray $\gamma$.

\begin{claim}\label{claim2}
The function $h_\infty$ constructed as above coincides with the Busemann function $b_\gamma$ associated with $\gamma$.
\end{claim}

\begin{proof}
Thanks to Lemma \ref{lem:Ld^2} and the fact that constant functions are locally $L$-harmonic, we know that for any $i$ we have $h_i \in \cD_{\loc}(\cE)$ and
\[
Lh_i = \frac{\alpha-1}{\dist(x_i,\cdot)} \qquad \text{on $X \backslash \{x_i\}$}.
\]
Therefore, for any $R\in(0,\dist(o,x_i))$, since $\dist(x_i,y)\ge D_i-\dist(y,o)\ge D_i-R$ holds for any $y \in B_R(o)$, we get
\begin{equation}\label{eq:D5.1}
|Lh_i| \le \frac{\alpha-1}{D_i-R} \qquad \text{on $B_R(o)$.}
\end{equation}
Then
\begin{equation}\label{eq:wesh}
|\cE(h_i,\phi)| = |\langle \phi, L h_i\rangle_{L^2}| \le \frac{\alpha-1}{D_i-R}\,\|\phi\|_1 \to 0 \qquad \text{when $i \to +\infty$}
\end{equation}
for any $\phi \in \cD_\rmc(\cE)$. Since $h_i \to h_\infty$ uniformly on compact sets, we can act as in the proof of Lemma \ref{lem:preparatory} to show that $h_\infty$ is locally $L$-harmonic. Lemma \ref{lem:preparatory} also shows that $b_\gamma$ is locally $L$-harmonic, hence $h_\infty - b_\gamma$ is so too.

Let us show now that $h_\infty - b_\gamma$ is non-negative. For any $i$ and $s \in [0,D_i]$, set:
\[
h_{i,s}(y)=s-\dist(y,c_i(s)) \qquad \forall y \in X.
\]
Since $\dist(x_i,y) \le \dist(x_i,c_i(s))+\dist(y,c_i(s))=D_i -s + \dist(y,c_i(s))$ for all $y \in X$, we have
\[
h_{i,s} \le h_i.
\]
As the curves $c_i$ pointwise converge to $\gamma$, the functions $h_{i,s}$ pointwise converge to $g_s:y \mto s - \dist(y,\gamma(s))$, so that letting $i$ tend to $+\infty$ provides
\[
g_s \le h_\infty
\]
and then letting $s$ tend to $+\infty$ gives
\[
b_\gamma \le h_{\infty}.
\]

By Lemma \ref{lem:ellipticHarnack}, since $h_\infty - b_\gamma$ is non-negative and locally $L$-harmonic, it is a constant function. But $b_\gamma(o)=0=h_i(o)$ for any $i$, so $h_\infty - b_\gamma$ is constantly equal to $0$.
\end{proof}

\subsubsection*{Step 2. Behavior of $H$ in the convergence $(X,\dist_i) \to (\uX,\udist)$ and link with $h_\infty$}

Recall that for any $1 \le j \le \ell$, the function $h_j$ has linear growth: $|h_j(x)| \le C_j (1+\dist(o,x))$ for any $x \in X$, where $C_j>0$ is some constant. Then the rescalings $h_j^{i}:=r_i^{-1}h_j$ are such that $|h_j^{i}(x)| \le C_j(r_i^{-1}+ \dist_i(o,x))$ for any $x \in X$, hence
\[
\|h_j^{i}\|_{L^\infty(B_r^{\dist_i}(o))} \le C_j(r_i^{-1}+r)\le C_j(1+r)
\]
holds for any $r>0$ and any $i$ such that $r_i>1$. Moreover, since $h_j$ is locally $L$-harmonic, it is strongly harmonic by Lemma \ref{lem:Lstrong} and then Lipschitz by Lemma \ref{lem:0109}: there is some constant $C_j'>0$ such that
\[
|h_j(x)-h_j(y)|\le C_j' \dist(x,y)
\]
for any $x,y \in X$. This implies that the sequence $\{h_j^{i}\}_i$ is asymptotically uniformly continuous on $\overline{B}_r(x)$. It is immediate to check that the rescalings $h_j^{i}$ are also strongly harmonic in $(X,\dist_i,\mu_i)$ where $\mu_i:=r_i^{-\alpha}\mu$. Then Propositions \ref{prop:AA1} and \ref{prop:stabstrongharmonic} imply that up to extracting a subsequence, we can assume that for any $j=1,\dots, l$, the functions~$h_j^{i}$ converge uniformly on all compact sets to a strongly harmonic function $\uh_j: \uX \to \setR$. We set
\[
\uH:= (\uh_1,\ldots,\uh_l): \uX \to \setR^l.
\]

\begin{claim}
For any given $\ux \in \uX$, the function $X \ni y \mto \beta(\uH(\ux),H(y))$ is a multiple of a Busemann function.
\end{claim}

\begin{proof}
Let $\{x_i\}_i \subset X$ be such that $x_i \To{\sGH} \ux$. Denote by $h_\infty$ the Busemann function associated to $\{x_i\}_i$ as in the previous step. Then for any $y \in X$,
\begin{align*}
\beta(\uH(\ux),H(y)) & = \lim_{i \to +\infty} \beta(H_i(x_i),H(y))\\ & = \lim_{i \to +\infty} \frac{1}{r_i} \beta(H(x_i),H(y))\\[-1pt]
& = \lim_{i \to +\infty} \frac{1}{2r_i} (\dist^2(o,x_i)+\dist^2(o,y)-\dist^2(x_i,y)) \qquad \text{by \eqref{eq:D6.3}}\\[-1pt]
& = \lim_{i \to +\infty} \frac{(\dist(o,x_i)-\dist(x_i,y))(\dist(o,x_i)+\dist(x_i,y))}{2r_i}\\[-1pt]
& = h_\infty(y) \Bigl(\frac{\udist(\uo,\ux)}{2} + \lim_{i \to +\infty} \frac{\dist(x_i,y)}{2r_i} \Bigr)
\end{align*}
since $\dist(o,x_i)-\dist(x_i,y) \to h_\infty(y)$. Now
\[
\underbrace{\frac{\dist(x_i,o)-\dist(o,y)}{r_i}}_{\to \udist(\ux,\uo)} \le \frac{\dist(x_i,y)}{r_i} \le \underbrace{\frac{\dist(x_i,o)+\dist(o,y)}{r_i}}_{\to \udist(\ux,\uo)}\,,
\]
so\vspace*{-4pt}
\begin{equation}
\beta(\uH(\ux),H(y)) = \udist(\uo,\ux) h_\infty(y).\qedhere
\end{equation}
\end{proof}

Note that we also have the following.

\begin{claim}\label{uclaim}
$\uH$ is an injective map. Moreover, $\Span(\uH(\uX)) = \setR^l$.
\end{claim}

\begin{proof}
Dividing $\eqref{eq:D6.4}$ by $r_i^2$ and taking $i\to+\infty$ implies $\udist^2(x,y) = Q(\uH(x)-\uH(y))$ for any $x,y \in X$, hence the injectivity of $\uH$. To prove the second part of the statement, let us show that $\uH(\uX)$ is contained in no hyperplane of $\setR^l$. Take a linear form $\lambda:\setR^l\to\setR$ vanishing on $\uH(\uX)$. Considering the convergent sequence $(X_i,\dist_i,o)\to(\uX,\udist,\uo)$, Proposition \ref{prop:AA1} implies that up to extraction the equi-Lipschitz functions $r_i^{-1} \lambda \circ H: X_i \to \setR$ converge to $0=\lambda \circ \uH: \uX \to \setR$ over $\overline{B}^{\dist_i}_r(o) \to\overline{B}_r(o)$ for any $r>0$. Therefore, we have\vspace*{-4pt}
\[
\sup_{\partial B_{r_i}(o)} |\lambda \circ H| = \sup_{\partial B_1^{\dist_i}(o)} |\lambda \circ H|= o(r_i).
\]
Being a linear combination of locally $L$-harmonic functions, $\lambda \circ H$ is locally $L$-harmonic hence strongly harmonic by Lemma \ref{lem:Lstrong}. Thus Lemma \ref{lem:0109} implies that $\lambda \circ H$ is constantly equal to $0$. Since $\Span(H(X))=\setR^l$, this implies $\lambda = 0$.
\end{proof}

\subsubsection*{Step 3. Construction of the bijection}

Our goal now is to construct a natural bijection between $\uX \backslash \{\uo\}$ and $\uS \times (0,+\infty)$, where $\uS:=\{\ux \in \uX: \udist(\uo,\ux)=1\}$.

Let us start with some heuristics. For $\ux \in \uX \backslash \{\uo\}$ given, we look for $\usigma \in \uS$ and $t\in(0,+\infty)$ uniquely determined by $\ux$. Here is how we are going to proceed:

\begin{enumerate}
\item
prove that there exists only one minimizing geodesic $\uc$ joining $\uo$ to $\ux$;

\item
show that $\uc$ extends in an unique way into a geodesic ray $\ugamma$.
\end{enumerate}
\noindent Indeed, the unique geodesic ray $\ugamma$ that we are going to construct necessarily crosses~$\uS$ at a single point $\usigma$ (otherwise $\ugamma$ would fail to be a minimizing geodesic) and it is such that $\ugamma(t)=\ux$ for a unique time $t>0$. Conversely, a pair $(\usigma,t)$ would uniquely determine a point $\ugamma(t) \in \uX$.

Let us proceed now with the construction. Take $\ux \in \uX \backslash \{\uo\}$. Let $\{x_i\}_i \subset X$ be such that $x_i \To{\sGH} \ux$. For any $i$, let $c_i$ be the minimizing $\dist_i$-geodesic joining $o$ to $x_i$. As~done previously, up to extracting a subsequence we can assume that $\{c_i\}_i$ converges uniformly on compact subsets of $[0,+\infty)$ to a geodesic ray $\gamma: [0,+\infty) \to X$. We know from Claim \ref{claim2} and the previous step that
\begin{equation}\label{eq:asabove}
\beta(\underline{H}(\ux),H(y)) = \udist(\uo,\ux) b_\gamma(y)
\end{equation}
holds for any $y \in X$, where $b_\gamma$ is the Busemann function associated with $\gamma$. Consider now a minimizing geodesic $\uc$ in $(\uX,\udist)$ joining $\uo$ to $\ux$ and set
\[
\uD:=\udist(\uo,\ux).
\]
For any $s \in [0,\uD]$, acting as we did to establish \eqref{eq:asabove}, we can prove that
\[
\beta(\underline{H}(\uc(s)),H(y)) = s b_\gamma(y)
\]
holds for any $y \in X$. Subtracting \eqref{eq:asabove} to this latter equality yields to
\begin{equation}%\label{eq:draft23}
\beta\Bigl(\underline{H}(\uc(s)) - \frac{s}{\uD}\,\uH(\ux), H(y)\Bigr) = 0
\end{equation}
for any $y \in X$. By Claim \ref{claim}, this implies
\begin{equation}%\label{eq:draft23}
\beta\Bigl(\underline{H}(\uc(s)) - \frac{s}{\uD}\,\uH(\ux), \xi \Bigr) = 0
\end{equation}
for any $\xi \in \setR^n$ and then
\begin{equation}\label{eq:H}
\underline{H}(\uc(s)) = \frac{s}{\uD}\,\uH(\ux)
\end{equation}
since $\beta$ is non-degenerate. Uniqueness of $\uc$ follows: if $\uc_1$ and $\uc_2$ are two minimizing geodesics joining $\uo$ to $\ux$, for any $s \in [0,\uD]$ one has
\[
\uH(\uc_1(s)) = \frac{s}{\uD}\,\uH(\ux) =\uH(\uc_2(s))
\]
and thus $\uc_1(s)=\uc_2(s)$ since $\uH$ is injective.

Let us show now that $\uc$ extends in an unique way into a geodesic ray. Our argument is inspired by the analysis done by Cheeger about generalized linear functions \cite[\S 8]{CheegerRademacher}. For any $i$, set $D_i:=\dist(o,x_i)$ and write $\gamma_i:[0,+\infty) \to X$ for the geodesic ray in $(X,\dist_i)$ defined by:
\[
\gamma_i(s) = \gamma(sr_i + D_i) \qquad \forall s>0.
\]
On the one hand, by Proposition \ref{prop:AA2}, we know that up to extracting a subsequence we can assume that $\{\gamma_i\}_i$ converges uniformly on compact subsets of $[0,+\infty)$ to a geodesic ray $\tilde{\gamma}:[0,+\infty] \to \uX$ whose associated Busemann function we denote by~$b_{\tilde{\gamma}}$. On the other hand, if we write $b_{\gamma_i}$ for the Busemann function associated with $\gamma_i$, we~can proceed as in Step 2 with $\dist_{i}, H_{i},\gamma_{i}$ in place of $\dist, H, \gamma$ respectively to get
\begin{equation}
b_{\gamma_{i}}(y) = \frac{\beta(\uH(\ux),H_{i}(y))}{\uD}
\end{equation}
for any $y \in X$. Then the sequence $(b_{\gamma_i})$ converges pointwise to the function
\[
F:\uX \ni \uy \mto \frac{\beta(\uH(\ux),\uH(\uy))}{\uD}
\]
and we have the following:

\skpt
\begin{claim}\mbox{}\vspace*{-1.2\baselineskip}
\begin{equation}\label{eq:uniqueness}
F=\underline{b}_{\tilde{\gamma}}+\uD.
\end{equation}
\end{claim}

\begin{proof}
Observe first that $F$ is strongly harmonic since it is a linear combination of the strongly harmonic functions $\uh_1,\ldots,\uh_l$. Let us show that $b_{\tilde{\gamma}}$ is strongly harmonic too. For any $i$, set
\[
p_i(x,y,t):=\frac{1}{(4\pi t)^{\alpha/2}}\, e^{-\sfrac{\dist_i^2(x,y)}{4t}}= r_i^\alpha p(x,y,r_i^2 t)
\]
for any $x,y \in X$ and $t>0$, and
\[
\underline{p}(\ux,\uy,t):= \frac{1}{(4 \pi t)^{\alpha/2}}\, e^{-\sfrac{\udist^2(\ux,\uy)}{4t}}
\]
for any $\ux,\uy \in X$ and $t>0$. Then for any $x,y \in X$ and $s,t>0$,
\begin{align*}
p_i(x,y,t+s) & = r_i^\alpha p(x,y,r_i^2 t + r_i^2 s) = r_i^\alpha \int_X p(x,z,r_i^2 t)p(z,y,r_i^2 s) \di \mu(z)\\
& = \int_X p_i(x,z,t)p(z,y,r_i^2s) \,r_i^\alpha\, \frac{\di \mu(z)}{r_i^\alpha} = \int_X p_i(x,z,t)p_i(z,y,s) \di \mu_i(z).
\end{align*}
For any $\ux, \uy \in \uX$ and $\{x_i\}_i, \{y_i\}_i \in X$ such that $x_i \To{\sGH} \ux$ and $y_i \To{\sGH} \uy$, the convergence $\dist_i(x_i,y_i) \to \udist(\ux,\uy)$ implies $p_i(x_i,y_i,t) \to \upp(\ux,\uy,t)$ for any $t>0$, hence:
\[
\underline{p}(\ux,\uy,t+s) = \int_X \underline{p}(\ux,\uz,t) \underline{p}(\uz,\uy,s) \di \umu(\uz) \qquad \forall \ux, \uy \in \uX, \, \forall \, t,s>0.
\]
By a standard procedure described for instance in \cite[\S 2]{Grigor'yan}, we can construct a Dirichlet form $\underline{\cE}$ on $(\uX,\udist,\umu)$ admitting a heat kernel given by $\underline{p}$. In particular, $(\uX,\udist,\umu,\underline{\cE})$ has an $\alpha$-dimensional Euclidean heat kernel. Writing $\underline{L}$ for the associated self-adjoint operator, we deduce from Lemma \ref{lem:preparatory} that $b_{\tilde{\gamma}}$ is locally $\underline{L}$-harmonic, then Lemma \ref{lem:Lstrong} implies that $b_{\tilde{\gamma}}$ is strongly harmonic.

Let us show now that $F - \underline{b}_{\tilde{\gamma}} \ge 0$. Take $i\in \setN$, $s>0$ and $y \in X$. Then $b_\gamma(y) \ge r_i s + D_i - \dist(\gamma(r_i s+D_i),y)$ by definition of a Busemann function, hence \hbox{$r_i^{-1}b_\gamma(y) \ge s+ r_i^{-1} D_i - \dist_i(\gamma_i(s),y)$}. Since
\begin{align*}
r_i^{-1}b_\gamma(y) & = \lim_{s \to +\infty} r_i^{-1} s - r_i^{-1} \dist(\gamma(s),y)\\
& = D_i r_i^{-1} + \lim_{s' \to +\infty} s' - \dist_i(\gamma(D_i+r_i s'),y),\qquad \text{where $s'=(s-D_i)r_i^{-1}$}\\ & = D_i r_i^{-1} + b_{\gamma_i}(y),
\end{align*}
we get $b_{\gamma_i}(y) \ge s - \dist_i(\gamma_i(s),y)$. Letting $i$ tend to $+\infty$ provides $F(y)\ge s - \udist(\tilde{\gamma}(s),y)$, after what letting $s$ tend to $+\infty$ gives $F\ge b_{\tilde{\gamma}}$.

By Lemma \ref{lem:ellipticHarnack}, we get that $F-b_{\tilde{\gamma}}$ is a constant function. Since $F(\ux)=\uD=\udist(\uo,\ux) $ and $b_{\tilde{\gamma}}(\ux)=0$, the claim is proved.
\end{proof}

Let $\ugamma:[0,+\infty)\to \uX$ be the concatenation of $\uc$ and $\tilde{\gamma}$, \ie
\[
\ugamma(t):=
\begin{cases}
\uc(t) & \text{if $0<t\le \uD$,}\\
\tilde{\gamma}(t-\uD) & \text{if $t\ge\uD$.}
\end{cases}
\]
By construction, $\ugamma$ is $1$-Lipschitz: $\udist(\ugamma(t),\ugamma(s))\le |t-s|$ for any $s,t>0$. Moreover, \eqref{eq:H} implies $F(\ugamma(t))=t$ when $0<t\le \uD$ while \eqref{eq:uniqueness} implies $F(\ugamma(t))=t$ when $t\ge \uD$. Since the function $F$ is $1$-Lipschitz we get $|t-s| \le \udist(\ugamma(t),\ugamma(s))$ for any $s,t>0$, thus~$\ugamma$ is a geodesic ray that extends $\uc$.

Let us show that this extension $\ugamma$ is unique. By \eqref{eq:uniqueness}, we have
\begin{equation}\label{eq:un}
\beta(\uH(\ux),\uH(\uy)) = \uD(\ub_{\tilde{\gamma}}(\uy)-\uD)
\end{equation}
for any $\uy \in \uX$ and we can obtain
\begin{equation}
\beta(\uH(\tilde{\gamma}(t)),\uH(\uy)) = t(\ub_{\tilde{\gamma}}(\uy)-\uD)
\end{equation}
for any $\uy \in \uX$ and $t>\uD$ by a similar reasoning. Then if $\ugamma'$ is another extension of $c$ obtained from a geodesic ray $\tilde{\gamma}'$ emanating from $\ux$, we get
\[
\beta(\uH(\tilde{\gamma}(t))-\uH(\tilde{\gamma}'(t)),\uH(\uy)) =0
\]
for any $\uy \in \uX$ and $t>\uD$, from which Claim $\ref{uclaim}$ yields $\tilde{\gamma}(t)=\tilde{\gamma}'(t)$.

\begin{remark}
Note that \eqref{eq:un} implies $\beta(\uH(\tilde{\gamma}(t)),\uH(\uy)) = t\beta(\uH(\tilde{\gamma}(1)),\uH(\uy))$ for all $\uy \in \uX$, hence
\begin{equation}\label{eq:conicalstructure}
\uH(\tilde{\gamma}(t)))= t\uH(\tilde{\gamma}(1)).
\end{equation}
\end{remark}

\subsubsection*{Step 4. Construction of the isometry}

Let $\Phi$ be the inverse of the bijection constructed in the previous step, \ie
\[
\Phi: \begin{aligned}
(0,+\infty) \times \uS & \to \uX \backslash\{\uo\} \\
(t,\usigma) & \mto \ugamma_{\usigma}(t),
\end{aligned}
\]
where $\ugamma_{\usigma}$ is the geodesic ray obtained by extending the minimizing geodesic joining $\uo$ to $\usigma$. Note that \eqref{eq:conicalstructure} implies
\begin{equation}\label{eq:conicalstructure2}
\uH(\Phi(t,\usigma))= t\uH(\Phi(1,\usigma))
\end{equation}
for any $(t,\usigma)\in (0,+\infty) \times \uS$. Let $\dist_C$ be the cone distance on $(0,+\infty)\times \uS$ defined by
\[
\dist_C^2((t,\usigma),(t',\usigma')):=(t-t')^2 + 4 t t' \sin^2\bigl(\sfrac{\udist_{\uS}(\usigma,\usigma')}{2}\bigr)
\]
for any $(t,\usigma),(t',\usigma') \in (0,+\infty) \times \uS$, where $\udist_{\uS}$ is the length distance associated with the distance on $\uS$ obtained by restricting $\udist$ to $\uS \times \uS$. We are going to establish
\begin{equation}\label{eq:isometry}
\udist(\ux,\ux') = \udist_C((t,\usigma),(t',\usigma'))
\end{equation}
for any $\ux = \Phi(t,\usigma), \, \ux' = \Phi(t',\usigma') \in \uX \backslash \{\uo\}$.

\begin{claim}\label{claim:isometrywithdelta}
There exists $\delta(\usigma,\usigma') \in [0,\pi]$ such that
\begin{equation}\label{eq:coeur2}
\udist^2(\ux,\ux') = (t-t')^2 + 4 tt' \sin^2(\sfrac{\delta(\usigma,\usigma')}{2}).
\end{equation}
\end{claim}

\begin{proof}
Choose $\{x_i\}_i, \{x_i'\}_i \subset X$ such that $x_i \To{\sGH} \ux$ and $x_i' \To{\sGH} \ux'$. For any $i$, divide \eqref{eq:D6.4} by $r_i^2$ to get $
\udist_i^2(x_i,x_i') = Q(H_i(x_i)-H_i(x_i'))$.
Letting $i$ tend to $+\infty$ implies $
\udist^2(\ux,\ux') = Q(\uH(\ux)-\uH(\ux'))$, hence
\[
\udist^2(\ux,\ux')
= Q(t \uH(\Phi(1,\usigma)) - t' \uH(\Phi(1,\usigma'))),
\]
thanks to \eqref{eq:conicalstructure2}. To compute $Q(t \uH(\Phi(1,\usigma)) - t' \uH(\Phi(1,\usigma')))$, let us use $\underline{h}_i(\usigma)$ as a shorthand for $\underline{h}_i(\Phi(1,\usigma))$. Then:
\begin{align*}
\udist^2(\ux,\ux') & = Q(t\uh_1(\usigma) - t' \uh_1(\usigma'), \ldots, t\uh_l(\usigma) - t' \uh_l(\usigma')) \nonumber \\
& = \sum_{i,j=1}^l B(x_i,x_j) (t\underline{h}_i(\usigma)-t'\underline{h}_j(\usigma'))(t \underline{h}_j(\usigma)-t'\underline{h}_j(\usigma')) \\
& = \begin{multlined}[t]
\biggl(\sum_{i,j=1}^l B(x_i,x_j) \uh_i(\usigma)\uh_j(\usigma)\biggr) t^2 + \biggl(\sum_{i,j=1}^l B(x_i,x_j) \uh_i(\usigma')\uh_j(\usigma')\biggr) (t')^2\\
 - 2 t t' \biggl(\sum_{i,j=1}^l B(x_i,x_j) \uh_i(\usigma')\uh_j(\usigma)\biggr)
\end{multlined}\\
& = Q(\uH(\usigma)) t^2 + Q(\uH(\usigma')) (t')^2 - 2 t t' \beta(\uH(\usigma),\uH(\usigma'))\\
& = \udist^2(\usigma,\uo)t^2 + \udist^2(\usigma',\uo) (t')^2 - 2 t t' \beta(\uH(\usigma),\uH(\usigma'))\\
& = t^2 + (t')^2 - 2 t t' \beta(\uH(\usigma),\uH(\usigma').
\end{align*}
Set $\uB(\ux,\ux'):=\frac{1}{2}(\udist^2(\uo,\ux)+\udist^2(\uo,\ux')-\udist^2(\ux,\ux'))$. Write \eqref{eq:D6.3} with $x=x_i$, $x'=x_i'$, divide by $r_i^2$ and let $r_i$ tend to $+\infty$ to get $\ubeta(\uH(\ux),\uH(\ux'))=\uB(\ux,\ux')$ and then
\begin{equation}\label{eq:coeur}
\udist^2(\ux,\ux') = t^2 + (t')^2 - 2 t t' \uB(\usigma,\usigma').
\end{equation}
Assuming $t=t'=1$ provides $\uB(\usigma,\usigma')=1-\frac{1}{2}\udist^2(\usigma,\usigma')$. The triangle inequality implies $\udist^2(\usigma,\usigma')\le 4$, thus $\uB(\usigma,\usigma') \in [-1,1]$, so we can set
\[
\cos(\delta(\usigma,\usigma')):=\uB(\usigma,\usigma')
\]
for some $\delta(\usigma,\usigma') \in [0,\pi]$.
\end{proof}

In particular, \eqref{eq:coeur2} implies
$\udist(\usigma_0,\usigma_1)=2\sin(\delta(\usigma_0,\usigma_1)/2)$
for any $\usigma_0,\usigma_1 \in \uS$.

\begin{claim}\label{claim:geodesicdelta}
The function $\delta$ defines a geodesic distance on $\uS$.
\end{claim}

\begin{proof}
Let us first show that $\delta$ defines a distance on $\uS$. We only prove the triangle inequality since the two other properties are immediate. For given $\usigma_0, \usigma_1,\usigma_2 \in \uS$, let us set $\alpha:= \delta(\usigma_0,\usigma_1)$, $\beta:=\delta(\usigma_1,\usigma_2)$ and $\gamma:=\delta(\usigma_0,\usigma_2)$. We can assume $\alpha + \beta \le \pi$ because otherwise we would have $\alpha + \beta > \pi \ge \gamma$, thus nothing to prove. For any $t,s,r>0$, the triangle inequality for $\udist$ written with \eqref{eq:coeur2} gives
\[
\sqrt{(t-s)^2 + 4 ts\sin^2(\alpha/2)} + \sqrt{(s-r)^2 + 4 sr\sin^2(\beta/2)} \ge \sqrt{(t-r)^2 + 4tr\sin^2(\gamma/2)}.
\]
Considering the three complex numbers $z_0=t$, $z_1 = se^{i\alpha}$ and $z_2 = re^{i(\alpha + \beta)}$, this can be rewritten as
\[
|z_0-z_1| + |z_1-z_2| \ge \sqrt{(t-r)^2 + 4tr\sin^2(\gamma/2)}.
\]
Choosing $s$ so that $z_0, z_1, z_2$ are aligned implies $|z_0-z_2| = |z_0-z_1| + |z_1 - z_2|$ thus
\[
\sqrt{(t-r)^2 + 4tr\sin^2((\alpha+\beta)/2)} = |z_0-z_2| \ge \sqrt{(t-r)^2+ 4tr\sin^2(\gamma/2)},
\]
which yields to $\alpha + \beta \ge \gamma$.

Let us show now that $\delta$ is geodesic. For given $\usigma_0, \usigma_1 \in \uS$ with $\usigma_0 \neq \usigma_1$, we aim at finding $\usigma_m \in \uS$ such that
\[
\delta(\usigma_0,\usigma_m)=\delta(\usigma_m,\usigma_1)=\frac12\delta(\usigma_0,\usigma_1).
\]
Let $c:[0,\udist(\usigma_0,\usigma_1)]\to X$ be the minimizing $\udist$-geodesic between $\usigma_0$ and $\usigma_1$. Assume first $\delta(\usigma_0,\usigma_1)<\pi$ so that $c(\udist(\usigma_0,\usigma_1)/2)\neq 0$. Then $c(\udist(\usigma_0,\usigma_1)/2)$ writes as $\Phi(s,\usigma_m)$ for some $(s,\sigma_m) \in (0,1) \times \uS$. We have
\[
\udist(\usigma_0,\usigma_m) = \udist(\usigma_m,\usigma_1) = \frac{1}{2}\udist(\usigma_0,\usigma_1),
\]
from which follows
\begin{equation}\label{eq:star}
(1-s)^2 + 4 s \sin^2\left(\sfrac{\alpha_0}{2}\right)=(1-s)^2 + 4 s \sin^2\left(\sfrac{\alpha_1}{2}\right)= \sin^2\left(\sfrac{\beta}{2}\right)
\end{equation}
thanks to \eqref{eq:coeur2}, where we have set $\alpha_0:=\delta(\usigma_0,\usigma_m)$, $\alpha_1:=\delta(\usigma_m,\usigma_1)$ and $\beta:=\delta(\usigma_0,\usigma_1)$.
Note first that \eqref{eq:star} immediately implies $\alpha_0=\alpha_1$. Moreover, for any $t>0$,
\[
\udist\left(\usigma_0,\Phi(t,\sigma_m)\right)+\udist\left(\usigma_1,\Phi(t,\sigma_m)\right)\ge \udist\left(\usigma_0,\usigma_1\right)
\]
implies
\[
\sqrt{(1-t)^2 + 4t \sin^2(\alpha_0/2)} + \sqrt{(1-t)^2 + 4t \sin^2(\alpha_1/2)} \ge 2 \sin(\sfrac{\beta}{2}),
\]
thus
\[
(1-t)^2 + 4t \sin^2(\alpha_0/2) \ge \sin^2(\sfrac{\beta}{2}).
\]
Therefore, the polynomial function $F: t \mto (1-t)^2 + 4t\sin^2(\alpha_o/2) - \sin^2(\beta/2)$
is non-negative and vanishes only at $t=s$, so $F'(s)=0$ hence
\[
2(1-s) = 4 \sin^2(\sfrac{\alpha_0}{2}).
\]
Plugging this in \eqref{eq:star} leads to $\sin^2(\beta/2) = \sin^2(\alpha_0/2)$
hence $\alpha_0=\beta_2$.
\end{proof}

Claim \ref{claim:geodesicdelta} and Lemma \ref{lem:length} implies $\delta=\udist_{\uS}$ which yields \eqref{eq:isometry} by Claim \ref{claim:isometrywithdelta}.

\subsection{Equality $l=\alpha$ and positive definiteness of $Q$}\label{subsec:4.3}

Since $\beta$ is non-degenerate, we can write
\begin{equation}\label{eq:decomposition}
\setR^l = E_+\oplus E_-,
\end{equation}
where $E_+$ is a subspace of $\setR^l$ with maximal dimension where $\beta$ is positive definite and $E_-$ is its $\beta$-orthogonal complement; $\beta$ is negative definite on $E_-$. We call $p_+$ the dimension of $E_+$ and $p_-$ the dimension of $E_-$. Note that $l = p_+ + p_-$ so in particular, $l \ge p_+$. Let us prove $p_+ \ge \alpha$, then $l = \alpha$, in order to reach our conclusion that is $l=p_+=\alpha$.

\subsubsection*{Step 1. $p_+ \ge \alpha$}

Let us write $\uH=(\uH_+,\uH_-)$, where $\uH_+:=\mathrm{proj}_{E_+} \circ \uH$ and $\uH_-:=\mathrm{proj}_{E_-} \circ \uH$, and $\mathrm{proj}_{E_+}, \mathrm{proj}_{E_-}$ are the projections associated to the decomposition \eqref{eq:decomposition}. Moreover, we set $q_+(v_+):=\beta(v_+,v_+)$ for any $v_+ \in E_+$ and $q_-(v_-):=\beta(v_-,v_-)$ for any $v_- \in E_-$. Then for any $\ux, \uy \in \uX$,
\[
Q(\uH(\ux)-\uH(\uy)) = q_+(\uH_+(\ux)-\uH_+(\uy)) + q_-(\uH_-(\ux)-\uH_-(\uy)),
\]
thus
\begin{equation}\label{eq:D9.1}
\udist^2(\ux,\uy) - q_-(\uH_-(\ux)-\uH_-(\uy)) = q_+(\uH_+(\ux)-\uH_+(\uy)).
\end{equation}
Since $q_- \le 0$, it follows from \eqref{eq:D9.1} that $q_+(\uH_+(\ux)-\uH_+(\uy)) \ge \udist^2(\ux,\uy)$. Moreover, $- q_-(\uH_-(\ux)-\uH_-(\uy))$ is bounded from above by $\lambda \udist^2(\ux,\uy)$, where $\lambda$ is the largest modulus an eigenvalue of $Q$ can have, so $q_+(\uH_+(\ux)-\uH_+(\uy)) \le (1+\lambda)\udist^2(\ux,\uy)$. Finally, since $\uH$ is injective, then $\uH_+$ is injective too. Therefore, the map $\uH_+$ is a bi-Lipschitz embedding of $(\uX,\udist)$ into $(E_+,\dist_{q_+})$ where $\dist_{q_+}(v_+,v_+'):=\sqrt{q_+(v_+-v_+')}$ for any $v_+, v_+' \in E_+$. This implies that $p_+$ is greater than or equal to the local Hausdorff dimension of $\uX$ which is equal to $\alpha$.

\subsubsection*{Step 2. $l=\alpha$}

Set $\tilde{\mu}:= (\Phi^{-1})_{\#}(\umu \measrestr \uX \backslash \{\uo\})$. Then $\tilde{\mu}$ is a Borel measure on $(0,+\infty)\times \uS$ equipped with $\dist_C$. We complete $(0,+\infty)\times \uS$ by adding the point $\uo$ corresponding to the tip of this metric cone. Let $\di t \otimes \nu_t$ be the disintegration of $\tilde{\mu}$ with respect to the first variable $t$ (we refer to \cite[2.5]{AmbrosioFuscoPallara} for the definition of disintegration of a measure). Since for any $\lambda,r>0$, we have $\umu(h_{\lambda}(B_r(\uo)))=\lambda^{\alpha} \umu(B_r(\uo))$, where $h_\lambda: (t,\usigma)\mto (\lambda t, \usigma)$, then $\di \nu_t=t^{\alpha-1}\nu_1$ for any $t>0$ and $\unu_1(\uS)=\alpha \omega_\alpha$. Let us write~$\unu$ instead of $\unu_1$.

\begin{claim}
For any $\usigma \in \uS$ and $\uh \in \ucV:=\Span(\uh_1,\dots,\uh_l)$,
\begin{equation}\label{proK}
\uh(\usigma) = \frac{\alpha}{\unu(\uS)} \int_{\uS} \cos(\udist_{\uS}(\usigma,\uphi)) \uh(\uphi) \di \unu(\uphi).
\end{equation}
\end{claim}

\begin{proof}
Take $\uh \in \ucV$ and $t>0$. Since $\uh_1,\dots,\uh_l$ are locally $\uL$-harmonic, then for any $\ux \in \uX$,
\begin{equation}\label{eq:13.56}
\uh(\ux) = \int_{\uX} \upp(\ux,\uy,t)\uh(\uy)\di\umu(\uy).
\end{equation}
Use the notation $\ux = \Phi(r,\usigma)$ and $\uy=\Phi(s,\uphi)$ and note that $\uh(\Phi(t,\usigma)) = r \uh(\usigma)$ and $\uh(\Phi(s,\uphi)) = s \uh(\uphi)$ thanks to \eqref{eq:conicalstructure}. Then \eqref{eq:13.56} writes
\begin{equation}\label{eq:draft21}
r \uh(\usigma)=\frac{1}{(4\pi t)^{\alpha/2}}\int_0^{+\infty} \int_{\uS} e^{\psfrac{-r^2-s^2+2rs\cos(\udist_{\uS}(\usigma,\phi))}{4t}}\uh(\uphi)s^\alpha\di s \di \unu(\uphi).
\end{equation}
Since
\[
\frac{\di}{\di r} \left(e^{\psfrac{-r^2+2rs\cos(\udist_{\uS}(\usigma,\uphi))}{4t}}\right) = \Bigl(- \frac{r}{2t} + \frac{s\cos(\udist_{\uS}(\usigma,\uphi))}{2t}\Bigr)e^{\psfrac{-r^2+2rs\cos(\udist_{\uS}(\usigma,\uphi))}{4t}},
\]
differentiating \eqref{eq:draft21} with respect to $r$ and evaluating at $r=0$ gives
\[
\uh(\usigma) = \frac{1}{(4\pi t)^{\alpha/2}} \int_0^{+\infty} e^{-\sfrac{s^2}{4t}}\,\frac{s^{\alpha+1}}{2t} \di s \int_{\uS} \cos(\udist_{\uS}(\usigma,\uphi)) \uh(\uphi) \di \unu(\uphi).
\]
A direct computation using the change of variable $\xi=\sfrac{s^2}{4t}$ shows that
\[
\frac{1}{(4\pi t)^{\alpha/2}} \int_0^{+\infty} e^{-\sfrac{s^2}{4t}}\,\frac{s^{\alpha+1}}{2t} \di s = \frac{1}{\omega_\alpha} = \frac{\alpha}{\unu(\uS)}.\qedhere
\]
\end{proof}
Set $\cW:=\{\uh_{|S} \sep \uh \in \ucV\}$. Note that \eqref{eq:conicalstructure} implies that the restriction map $\ucV\to \cW$ is a bijection, hence $\dim \cW=l$. We introduce the operator
$\mathcal{K}\colon L^2(\uS,\di\unu)\to L^2(\uS,\di\unu)$ defined by
\[
\mathcal{K}(f)(\usigma):= \frac{\alpha}{\unu(\uS)} \int_{\uS} \cos(\udist_{\uS}(\usigma,\uphi)) f(\uphi) \di \unu(\uphi)
\]
for any $f \in L^2(\uS,\di\unu)$ and $\umu$-a.e.\ $\usigma \in \uS$. Since for any $\usigma,\usigma' \in \uS$,
\begin{equation}\label{expcos}\cos(\dist_S(\usigma,\usigma'))=\uB(\usigma,\usigma')=\beta\left(\uH(\Phi(1,\usigma)),\uH(\Phi(1,\usigma'))\right)=\sum_{i,j} B(x_i,x_j) \uh_i(\usigma)\uh_j(\usigma'),
\end{equation}
then the image of $\mathcal{K}$ is contained in $\cW$ and according to \eqref{proK}, we have
\[
\mathcal{K}f=f \qquad \text{for every $f \in \cW$.}
\]
Hence $\mathcal{K}$ is the orthogonal projection onto $\cW$ and
if $\underline{k}_1,\ldots, \underline{k}_l$ form an orthonormal basis of $\cW$ for the $L^2(\uS,\unu)$ scalar product, then for any $f\in L^2(\uS,\di\nu)$:
\[
\mathcal{K}(f)(\usigma)= \sum_{i=1}^l \underline{k}_i(\usigma) \int_{\uS} \underline{k}_i(\uphi) f(\uphi) \di \unu(\uphi).
\]
This implies
\begin{equation}\label{eq:0409}
\frac{\alpha}{\unu(\uS)}\, \cos(\udist_{\uS}(\usigma,\uphi)) = \sum_{i=1}^l \underline{k}_i(\usigma) \underline{k}_i(\uphi)
\end{equation}
for $\unu\otimes \unu$-a.e.\ $(\usigma,\uphi) \in \uS \times \uS$. Since for any $i$, the function $\underline{k}_i$ admits a continuous representative -- still denoted by $\underline{k}_i$ -- defined by
\[
\underline{k}_i (\usigma) = \int_{\uS} \cos(\udist_{\uS}(\usigma,\uphi)) \uh(\uphi) \di \unu(\uphi)
\]
for any $\usigma \in \uS$, then \eqref{eq:0409} holds for all $(\usigma,\uphi) \in \uS \times \uS$. In particular, we can take $\usigma = \uphi$ in \eqref{eq:0409} to get
\[
\frac{\alpha}{\unu(\uS)} = \sum_{i=1}^l \underline{k}_i(\usigma)^2.
\]
Integrating over $\uS$ with respect to $\unu$ gives $\alpha = l$.

\subsection{Conclusion}

From the previous subsections, we get that $H$ is an isometric embedding of $(X,\dist)$ into $(\setR^l,\dist_Q)$ or, as explained at the beginning of this section, into $(\setR^l,\dist_e)$. Therefore, $H(X)$ equipped with the restriction of $\dist_e$ is geodesic. Minimizing geodesics in $(\setR^l,\dist_e)$ being straight lines, this implies that $H(X)$ is convex. Being also closed, $H(X)$ is equal to its closed convex hull that is equal to $\setR^l$ by the proof of Claim \ref{claim}, hence Theorem \ref{th:main} is proved.

\Changelback
\section{Almost rigidity result for the heat kernel}\label{sec:5}

In this section, we show how our rigidity result (Theorem \ref{th:main}) provides an almost rigidity result (Theorem \ref{th:almostrigidity}). We fix a positive constant $T>0$, a positive integer $n$, and we recall that $\setB^n_r$ stands for an Euclidean ball in $\setR^n$ with radius $r>0$ (where this ball is centered as no importance), and $\dist_{\sGH}$ for the Gromov-Hausdorff distance.

We begin with the following lemma:
\begin{lemma}\label{lem:voleu} If $(X,\dist,\mu)$ is a complete metric measure space endowed with a symmetric Dirichlet form $\cE$
admitting a heat kernel $p$ such that for some $\gamma >1$,
\begin{equation}\label{eq:estheatkernel}
\frac{\upgamma ^{-1}}{(4 \pi t)^{n/2}}\, e^{-\upgamma\sfrac{\dist^2(x,y)}{4t}}\le p(x,y,t) \le \frac{\upgamma }{(4 \pi t)^{n/2}}\, e^{-\sfrac{\dist^2(x,y)}{4\upgamma t}}
\end{equation}
for all $x, y \in X$ and $t \in (0,T]$, then there exists positive constants $c(n,\upgamma), C(n,\upgamma)$ such that for any $x\in X$ and $r\le \sqrt{T}$,
\[
c(n,\upgamma)\, r^n\le \mu(B_r(x)))\le C(n,\upgamma)\, r^n.
\]
\end{lemma}
\begin{remark} The upper bound is quite classical, the novelty is the lower bound which was nonetheless known for stochastically complete spaces (see \cite[Th.\,2.11]{Grigor'yan}).
\end{remark}
\proof For any $x \in X$ and $r>0$, integrating the lower bound in \eqref{eq:estheatkernel} gives
\[
e^{-\upgamma\sfrac{r^2}{4t}} \mu\left(B_r(x)) \right)\le \int_{B_r(x)} e^{-\upgamma\sfrac{\dist^2(x,y)}{4t}}\di\mu(y)\le \gamma (4 \pi t)^{n/2},
\]
hence $\mu\left(B_r(x)) \right)\le \gamma (4 \pi t)^{n/2} e^{\upgamma\sfrac{r^2}{4t}}$ for any $t \in (0,T]$. Consequently, when $r\le \sqrt{T}$, choosing $t=r^2$ provides
\begin{equation}\label{eq:uppervolume} \mu\left(B_r(x) \right)\le e^{\sfrac \upgamma 4} \upgamma(4 \pi)^{n/2}\, r^n,
\end{equation}
while when $r\ge \sqrt{T}$, choosing $t=T$ gives
\begin{equation}\label{eq:uppervolumeT}
\mu\left(B_r(x) \right)\le \upgamma\,(4 \pi T)^{n/2}\, e^{\upgamma\sfrac{r^2}{4T}}.
\end{equation}
Note that \eqref{eq:uppervolume} is the desired upper bound. Take $t \in (0,T/2]$. Combining \eqref{eq:estheatkernel} with the Chapman-Kolmogorov formula, we get
\[
\frac{\upgamma^{-1}}{(8\pi t)^{n/2}} \le p(x,x,2t) = \int_X p(x,y,t)^2 \di \mu(y) \le \frac{\upgamma^2}{(4 \pi t)^n} \int_X e^{-\sfrac{\dist^2(x,y)}{2 \gamma t}} \di \mu(y),
\]
hence
\begin{equation}\label{eq:est4}
\upgamma^{-3}\,(2\pi)^{n/2}\, t^{n/2} \le \int_X \hspace*{-2mm}e^{-\sfrac{\dist^2(x,y)}{2\upgamma t}}\di\mu(y) \le \mu(B_r(x))+\int_{X\setminus B_r(x)}\hspace*{-6mm} e^{-\sfrac{\dist^2(x,y)}{2\upgamma t}}\di\mu(y).
\end{equation}
From now on, assume $r \le \sqrt{T}$. By Cavalieri's principle and the estimates \eqref{eq:uppervolume} and~\eqref{eq:uppervolumeT}, we get
\begin{multline}\label{eq:est1}
\int_{X\setminus B_r(x)} e^{-\sfrac{\dist^2(x,y)}{2\upgamma t}}\di\mu(y)=\int_r^{+\infty} e^{-\sfrac{\rho^2}{2\upgamma t}}\,\frac{\rho}{\upgamma t}\,\mu(B_\rho(x))\di\rho\\
\le (4 \pi)^{n/2}e^{\sfrac{\upgamma}{4}} \int_r^{\sqrt{T}} e^{-\sfrac{\rho^2}{2\upgamma t}}\,\frac{\rho}{ t} \, \rho^n \di \rho+(4 \pi T)^{n/2}\int_{\sqrt{T}}^{+\infty} e^{-\sfrac{\rho^2}{2\upgamma t}}\,\frac{\rho}{ t}\, e^{\upgamma\sfrac{\rho^2}{4T}}\di\rho.
\end{multline}
A direct computation shows that for any $n \in \setN$, there exists $C_0>0$ depending only on $n$ such that for any $A\ge 1$,
\[
\int_A^{+\infty} e^{-\sfrac{\xi^2}{2}} \xi^{n+1} \di \xi \le C_0 A^n e^{-\sfrac{A^2}{2}}.
\]
Therefore, using the change of variable $\xi= \rho/\sqrt{\upgamma t}$ to get
\[
\int_r^{\sqrt{T}} e^{-\sfrac{\rho^2}{2\upgamma t}}\,\frac{\rho}{ t}\, \rho^n \di \rho\le \int_r^{+\infty} e^{-\sfrac{\rho^2}{2\upgamma t}}\frac{\rho}{ t}\, \rho^n \di\rho = \upgamma^{\Psfrac n2+1} t ^{n/2} \int_{r/\sqrt{\upgamma t}}^{+\infty} e^{-\sfrac{\xi^2}{2}}\xi^{n+1} \di \xi,
\]
we obtain that $r\ge \sqrt{\upgamma t}$ implies
\begin{equation}\label{eq:est2} \int_r^{\sqrt{T}} e^{-\sfrac{\rho^2}{2\upgamma t}}\frac{\rho}{\ t}\, \rho^n \di \rho\le C_0 \upgamma r^n e^{-\sfrac{r^2}{2\upgamma t}}.
\end{equation}
To bound the second term in \eqref{eq:est1}, assume $t\le T/\upgamma^2$. Then a straightforward computation shows that $-\sfrac{\rho^2}{2\gamma t} + \sfrac{\gamma \rho^2}{4T} \le - \sfrac{\rho^2}{4 \gamma t}$ holds, thus
\begin{equation}\label{eq:est3}
\int_{\sqrt{T}}^{+\infty} e^{-\sfrac{\rho^2}{2\upgamma t}}\,\frac{\rho}{ t}\, e^{\upgamma\sfrac{\rho^2}{4T}}\di \rho \le \int_{\sqrt{T}}^{+\infty} e^{-\sfrac{\rho^2}{4\upgamma t}}\,\frac{\rho}{ t}\, \di \rho = 2\upgamma e^{-\sfrac{T}{4\upgamma t}}.
\end{equation}
Combining \eqref{eq:est1}, \eqref{eq:est2} and \eqref{eq:est3}, we get existence of a constant $C>0$ depending only on $n$ such that if
$r^2 \ge \gamma^2 t$ (this implies both $r\ge \sqrt{\gamma t}$ and $t \le T/\gamma^2$), then
\[
\int_{X\setminus B_r(x)} e^{-\sfrac{\dist^2(x,y)}{2\upgamma t}}\di\mu(y)\le C\bigl(\upgamma e^{\sfrac{\upgamma}{4}}r^n e^{-\sfrac{r^2}{2\upgamma t}}+\upgamma T^{n/2} e^{-\sfrac{T}{4\upgamma t}}\bigr).
\]
Then \eqref{eq:est4} implies
\begin{equation}\label{eq:lem}
\upgamma^{-3}\,(2\pi)^{n/2}\, t^{n/2} \le \mu(B_r(x))+C\bigl(\upgamma e^{\sfrac{\upgamma}{4}}r^n e^{-\sfrac{r^2}{2\upgamma t}}+\upgamma T^{n/2} e^{-\sfrac{T}{4\upgamma t}}\bigr)
\end{equation}
for any $t \in (0,r^2/\gamma^2)$, what can be rewritten as
\[
c'(n,\gamma)t^{n/2}\le \mu(B_r(x))+Ct^{n/2}\bigl(\upgamma e^{\sfrac{\upgamma}{4}} F(r^2/t)+\upgamma G(T/t)\bigr),
\]
where $c'(n,\gamma):=\upgamma^{-3}\,(2\pi)^{n/2}$ and $F(s):=s^{n/2}e^{-\sfrac{s}{2\gamma}}$, $G(s):=s^{n/2}e^{-\sfrac{s}{4\gamma}}$ for any $s\ge 0$. The function $G$ is decreasing on $(2 n\gamma,+\infty)$ so if $r^2/t \ge 2 n\gamma$, since $r^2 \le T$, we~get $G(T/t) \le G(r^2/t)$. As $\lim_{s \to +\infty}F(s)=\lim_{s \to +\infty} G(s)=0$, then there exists $s(n,\gamma)>0$ such that if $s\ge s(n,\gamma)$,
\[
C\bigl(\upgamma e^{\sfrac{\upgamma}{4}} F(s)+\upgamma G(s)\bigr) \le \frac{c'(n,\gamma)}{2}.
\]
Then for any $t >0$ such that $r^2/t \ge \max(\gamma^2,2 n\gamma,s(n,\gamma))=:\theta(n,\gamma)$, we get
\[
\frac{c'(n,\gamma)}{2}\,t^{n/2} \le \mu(B_r(x)).
\]
Choosing $t=t(r)$ such that $\theta(n,\gamma) t \le r^2 \le 2 \theta(n,\gamma)t$, we get
\[
\frac{c'(n,\gamma)}{2^{n/2+1}\theta(n,\gamma)^{n/2}}\, r^{n} \le \mu(B_r(x)).\qedhere
\]
\endproof

We shall also need the next proposition.
\begin{proposition}\label{prop:5.3} Let $(X,\dist,\mu)$ be a measure metric space satisfying the local doubling condition, namely there exists $r_o>0$ and $C_D>0$ such that $\mu(B_{2r}) \le C_D \mu(B_r)$ for any $r \in (0,r_o)$, and such that for some $\alpha >0$, we have
\[
\int_X \frac{1}{(4 \pi t)^{\alpha/2}}\, e^{-\sfrac{\dist^2(x,z)}{4t}}\frac{1}{(4 \pi s)^{\alpha/2}}\, e^{-\sfrac{\dist^2(z,y)}{4s}} \di \mu(z)=\frac{1}{(4 \pi (t+s))^{\alpha/2}}\, e^{-\sfrac{\dist^2(x,y)}{4(t+s)}}
\]
for all $x,y \in X$ and $t,s \in (0,T)$. Then there exists a symmetric Dirichlet form $\cE$ on $(X,\dist,\mu)$ admitting an $\alpha$-Euclidean heat kernel.
\end{proposition}
\proof By \cite[Lem.\,3.9]{Carron}, the space $(X,\dist,\mu)$ satisfies $
\mu(B_R(x))/\mu(B_r(x)) \le c_o e^{c_1 R/r}$ for any $x \in X$, $r\in(0,r_o)$ and $R\ge r$, where $c_o$ and $c_1$ depend only on $C_D$. For any $C>0$ and $z \in X$, applying \eqref{eq:Fub} with $\phi(\lambda)=\lambda^2 e^{-C\lambda^2}$ and $g(y)=\dist(z,y)$ yields to
\begin{equation}
\begin{aligned}\label{eq:exp}
\int_X e^{-C\dist^2(z,y)}&\di \mu(y) = \int_0^{+\infty} 2 C \lambda e^{-C\lambda^2} \mu(B_\lambda(y))\di \lambda \\
& \le \left(\int_0^r 2 C \lambda e^{-C\lambda^2} \di \lambda + \int_r^{+\infty} 2 C \lambda c_o e^{-C\lambda^2 + c_1 \lambda/r} \di \lambda\right)\mu(B_r(y)) \\
& \le c_2 \mu(B_r(y))
\end{aligned}
\end{equation}
for any $r \in(0,r_o)$, where $c_2$ depends only on $r$, $C_D$ and $C$. For any $x,y\in X$ and $t\in \setC\setminus \{0\}$,
we set
\[
\bP_\alpha(x,y,t):=\frac{1}{(4 \pi t)^{\alpha/2}}\, e^{-\sfrac{\dist^2(x,y)}{4t}}.
\]
Take $x,y \in X$ and $t\in(0,T)$. By assumption, the identity
\[
\int_X \bP_\alpha(x,z,t)e^{-\sfrac{\dist^2(z,y)}{4s}}\di \mu(z)= \Bigl(\frac{s}{t+s}\Bigr)^{\sfrac \alpha2}e^{-\sfrac{\dist^2(z,y)}{4(t+s)}}
\]
is valid for any $s \in (0,T-t)$. However both expressions are holomorphic in $s\in \setC_+:=\{z\in \setC\sep \mathrm{Re} z>0\}$ (the left-hand side can be proved holomorphic by a suitable application of the dominated convergence theorem using \eqref{eq:exp}), thus the identity holds for any $s\in \setC_+$. Freezing $s$ and letting $t$ be variable, we can apply the same reasoning to get the identity valid for any $s, t \in \setC_+$. In particular, we obtain:
\[
\int_X \bP_\alpha(x,z,t) \bP_\alpha(z,y,s)\di\mu(z)= \bP_\alpha(x,y,t+s) \qquad \forall s,t >0.
\]
Thus for any $x \in X$ and $t>0$,
\begin{align*}
\int_X \bP_\alpha(x,z,t) \di \mu(z) & = (4 \pi t)^{\alpha/2} \int_X \Bigl(\frac{1}{(4 \pi t)^{\alpha/2}} e^{-\sfrac{\dist^2(x,z)}{8t}} \Bigr)^2 \di \mu(z)\\
& = (16 \pi t)^{\alpha/2} \int_X \bP_\alpha(x,z,2t)^2 \di \mu(z)\\ & =(16 \pi t)^{\alpha/2} \bP_\alpha(x,x,4t) = 1.
\end{align*}
This easily implies that for any $f\in L^2(X,\mu)$, if
$f_t(x)=\int_X \bP_\alpha(x,z,t)f(z) \di\mu(z)$, then
\[
\lim_{t\to 0+} \|f_t-f\|_{L^2}=0.
\]
Then by a standard procedure described for instance in \cite[\S 2]{Grigor'yan}, we can build a symmetric Dirichlet form whose heat kernel is $\bP_\alpha$.
\endproof

We can now prove Theorem \ref{th:almostrigidity}.
\proof The metric spaces considered in this proof are all complete. Assume that the result is not true. Then there exists some $\upepsilon>0$ such that for any $\updelta>0$ we can find:\begin{itemize}
\item $T_\updelta>0$,
\item a metric measure space $(X_\updelta,\dist_\updelta, \mu_\updelta)$ endowed with a symmetric Dirichlet form $\cE_\updelta$
admitting a heat kernel $p_\updelta$ satisfying
\[
(1-\updelta)\, \frac{1}{(4 \pi t)^{n/2}}\, e^{-\sfrac{\dist_\updelta^2(x,y)}{4(1-\updelta)t}}\le p_\updelta(x,y,t) \le (1+\updelta)\, \frac{1}{(4 \pi t)^{n/2}}\, e^{-\sfrac{\dist_\updelta^2(x,y)}{4(1+\updelta) t}}
\]
for any $x,y \in X_\updelta$ and $t\in (0,T_\updelta]$,
\item $x_\updelta\in X_\updelta$ and $r_\updelta\in (0,\sqrt{T_\updelta}]$ such that $\dist_{\sGH}\left(B_{r_\updelta}(x_\updelta), \setB^n_{r_\updelta}\right) \ge \upepsilon r_\updelta$.
\end{itemize}
By a rescaling of the distance and of the measure, we can assume that $r_\updelta=1$ and $T_\updelta=1$.
It follows from Lemma \ref{lem:voleu} that the set of pointed metric measure space
\[
\left\{(X_\updelta,\dist_\updelta,\mu_\updelta,x_\updelta)\right\}_{\updelta\in (0,1/2)}
\]
satisfies a uniform local doubling condition, thus it is precompact for the pointed measure Gromov-Hausdorff topology. Therefore, we can consider an infinitesimal sequence $\{\updelta_\ell\}_\ell \subset (0,1/2)$ and a sequence of pointed metric measure spaces
\[
\left\{(X_\ell,\dist_\ell,\mu_\ell,x_\ell)\right\}_{\ell}
\]
converging to some pointed metric measure space $(X_\infty,\dist_\infty,\mu_\infty,x_\infty)$ such that for any~$\ell$:
\begin{itemize}
\item the space $(X_\ell,\dist_\ell,\mu_\ell)$ is endowed with a symmetric Dirichlet form $\cE_\ell$
admitting a heat kernel $p_\ell$ satisfying
\begin{multline}\label{eq:14.17}
(1-\updelta_\ell)\, \frac{1}{(4 \pi t)^{n/2}}\, e^{-\sfrac{\dist_\ell^2(x,y)}{4(1-\updelta_\ell)t}}\le p_\ell(x,y,t)\\[-5pt]
\le (1+\updelta_\ell)\, \frac{1}{(4 \pi t)^{n/2}}\, e^{-\sfrac{\dist_\ell^2(x,y)}{4(1+\updelta_\ell)t}}
\end{multline}
for any $x,y \in X_\ell$ and $t\in (0,1]$,
\item $\dist_{\sGH}\left(B_1(x_\ell), \setB^n_1\right) \ge \upepsilon. $
\end{itemize}
In particular, letting $\ell$ tend to $+\infty$ gives:
\begin{equation}\label{disHG}\dist_{\sGH}\left(B_1(x_\infty), \setB^n_1\right) \ge \upepsilon.\end{equation}
Since for any $\ell$ we have
\[
\int_{X_\ell}p_\ell(x,z,t)p_\ell(z,y,s)d\mu_\ell(z)=p_\ell(t+s,x,y)
\]
for all $x,y \in X_\ell$ and $t,s>0$, we deduce from \eqref{eq:14.17} that when $t+s<1$,
\[
\frac{(1-\updelta_\ell)^{\Psfrac n2+1}}{(1+\updelta_\ell)^{n+1}}\ \bP_n\Bigl(x,y,\frac{1-\updelta_\ell}{1+\updelta_\ell}(t+s)\Bigr)\le \int_{X_\ell}\bP_n(x,z,t)\bP_n(z,y,s)d\mu_\ell(z)
\]
and
\[
\int_{X_\ell}\bP_n(x,z,t)\bP_n(z,y,s)d\mu_\ell(z)\le \frac{(1+\updelta_\ell)^{\Psfrac n2+1}}{(1-\updelta_\ell)^{n+1}}\ \bP_n\Bigl(x,y,\frac{1+\updelta_\ell}{1-\updelta_\ell}(t+s)\Bigr).
\]
From this, we obtain for any $x,y\in X_\infty$ and any $t,s>0$ with $t+s<1$,
\[
\int_{X_\infty}\bP_n(x,z,t)\bP_n(z,y,s)d\mu_\infty(z)= \bP_n(x,y,t+s).
\]
Then Proposition \ref{prop:5.3} and Theorem \ref{th:main} imply that $(X_\infty,d_\infty)$ is isometric to $(\setR^n,\dist_e)$. But this is in contradiction with \eqref{disHG}.
\endproof

\section{A new proof of Colding's almost rigidity theorem}\label{sec:6}

In this section, we show how our almost rigidity result (Theorem \ref{th:almostrigidity}) can be used to give an alternative proof of the almost rigidity theorem for the volume of Riemannian manifolds with non-negative Ricci curvature (Theorem \ref{th:Colding}). Here again $n$ is a fixed positive integer and $\mathbb{B}_r^n$ is an Euclidean ball in $\setR^n$ with radius $r>0$.

We recall that whenever $(M^n,g)$ has non-negative Ricci curvature, the Bishop-Gromov comparison theorem states that the function $r\mto \omega_n^{-1} r^{-n}\vol (B_r(x))$ is non-increasing for any $x \in M$
and the quantity
\begin{equation}\label{eq:theta}
\uptheta=\lim_{r\to +\infty} \frac{\vol (B_r(x))}{\omega_n\, r^n}
\end{equation}
does not depend on $x$. When $\uptheta>0$, we say that $(M^n,g)$ has Euclidean volume growth, in which case one has
\begin{equation}\label{eq:volumelowerbound}
\vol (B_r(x)) \ge \uptheta \omega_n r^n
\end{equation}
for any $x \in X$ and $r>0$. Note that a manifold satisfying \eqref{eq:vol} has Euclidean volume growth with $\uptheta \ge 1-\delta$. Our proof of Theorem \ref{th:Colding} is a direct application of Theorem~\ref{th:almostrigidity} together with the following heat kernel estimate.
\begin{theorem}\label{th:estimate} There exists a function $\upgamma\colon [0,1]\to [1,\infty)$ satisfying $\lim_{\theta\to 1^-} \upgamma(\theta)=1$ such that
whenever $(M^n,g)$ is a complete Riemannian manifold with non-negative Ricci curvature and Euclidean volume growth, then the heat kernel $p$ of $(M^n,g)$ satisfies
\[
\frac{1}{(4\pi t)^{\sfrac n2}}\,e^{-\sfrac{d^2(x,y)}{4t}}\le p(x,y,t)\le \upgamma(\uptheta) \frac{1}{(4\pi t)^{\sfrac n2}}\,e^{-\sfrac{d^2(x,y)}{\upgamma(\uptheta)4t}}
\]
for all $x,y \in M$ and $t>0$, where $\uptheta$ is given by \eqref{eq:theta}.
\end{theorem}

\begin{remark} Our proof of the above heat kernel upper bound follows the arguments of P. Li, L-F. Tam and J. Wang \cite{LiTamWang}.
\end{remark}

\begin{proof}
The lower bound is the comparison theorem of J.~Cheeger and S-T.~Yau \cite{CheegerYau}: for any $t>0$ and $x,y\in M$, we have
\begin{equation}\label{eq:CheegerYau}\bP_n(x,y,t)\le p(x,y,t),\end{equation}
where
$\bP_n(x,y,t)=(4\pi t)^{-\sfrac n2}e^{-\sfrac{d^2(x,y)}{4t}}$. Consequently we only need to prove the upper bound.

Take $x,y \in X$ and $t>0$. We shall need the following estimates from P. Li and S-T. Yau (see \cite[Form.\,(2.1)]{LiTamWang}): for any $r,\tau>0$,
\begin{equation}\label{eq:LY1}\int_{B_r(x)}p(x,z,\tau)\di\vol(z)\ge \int_{\mathbb{B}^n_r}\frac{1}{(4\pi \tau)^{\sfrac n2}}\,e^{-\sfrac{\|\xi\|^2}{4\tau}}\di \xi
\end{equation}
and
\begin{equation}\label{eq:LY2}
\int_{M\setminus B_r(x)}p(x,z,\tau)\di \vol(z)\le \int_{\setR^n\setminus \mathbb{B}_r^n}\frac{1}{(4\pi \tau)^{\sfrac n2}}\,e^{-\sfrac{\|\xi\|^2}{4\tau}}\mathrm{\di}\xi.
\end{equation}

For $\delta>0$ to be precisely chosen later, set $r:=(1+\delta)^{-1}\dist(x,y)$ and $\tau:= (1+\delta)t$. Note that $B_{r}(x) \cap B_{\delta r}(y) = \emptyset$. By the Harnack inequality of P.~Li and S-T.~Yau \cite{LiYau}, we have
\begin{equation}\label{LY}
p(x,y,t)\le \left(\frac{\tau}{t}\right)^{\sfrac n2}\, e^{\sfrac{\dist(z,y)^2}{4(\tau-t)}} p(x,z,\tau)\end{equation}
for every $z\in M$, so that averaging over the ball $B_{\delta r}(y)$ gives
\begin{equation}
\begin{aligned}\label{eq:Colding111}p(x,y,t)&\le e^{\sfrac{\delta^2r^2}{4(\tau-t)} } \left(\frac{\tau}{t}\right)^{\sfrac n2}\fint_{B_{\delta r}(y)}p(x,z,\tau)\di\vol(z) \\
&\le e^{\sfrac{\delta d^2(x,y)}{4(1+\delta)^2t}} \left(\frac{\tau}{t}\right)^{\sfrac n2}\fint_{B_{\delta r}(y)}p(x,z,\tau)\di\vol(z).
\end{aligned}
\end{equation}
Now
\begin{align*}
\int_{B_{\delta r}(y)}&p(x,z,\tau)\di \vol(z)\\
&=\int_{M\setminus B_r(x)}p(x,z,\tau)\di \vol(z)- \int_{M\setminus (B_r(x)\cup B_{\delta r}(y))} p(x,z,\tau)\di \vol(z)\\
&\le \int_{\setR^n\setminus \mathbb{B}_r^n}\frac{1}{(4\pi \tau)^{\sfrac n2}}e^{-\sfrac{\|\xi\|^2}{4\tau}}\di \xi - \int_{M\setminus (B_r(x)\cup B_{\delta r}(y)}\mathbb{P}_n(x,z,\tau)\di \vol(z)
\end{align*}
thanks to \eqref{eq:LY2} and \eqref{eq:CheegerYau}. Continuing,
\begin{align*}
\int_{B_{\delta r}(y)}&p(x,z,\tau)\di \vol(z) \le \int_{\setR^n\setminus \setB_r^n}\frac{1}{(4\pi \tau)^{\sfrac n2}}\,e^{-\sfrac{\|\xi\|^2}{4\tau}}\di \xi\\
&\hspace*{.8cm} - \int_{M\setminus B_r(x)}\bP_n(x,z,\tau)\di \vol(z)
+ \int_{B_{\delta r}(y)}\bP_n(x,z,\tau)\di \vol(z)\\
&\le \int_{\setR^n\setminus \setB_r^n}\frac{1}{(4\pi \tau)^{\sfrac n2}}\,e^{-\sfrac{\|\xi\|^2}{4\tau}}\di \xi \\
&\hspace*{.8cm}- \int_{M\setminus B_r(x)}\bP_n(x,z,\tau)\di \vol(z)+\vol(B_{\delta r}(y))\, \frac{1}{(4\pi \tau)^{\sfrac n2}}\,e^{-\sfrac{(\dist(x,y)-\delta r)^2}{4\tau}}.
\end{align*}
By Cavalieri's principle and \eqref{eq:volumelowerbound}, we have
\begin{multline*}
\int_{M\setminus B_r(x)}\bP_n(x,z,\tau)\di \vol(z) =\int_{r}^{+\infty} \frac{1}{(4\pi \tau)^{\sfrac n2}}\,e^{-\sfrac{s^2}{4\tau}}\frac{s}{2\tau}\vol(B_s(x))\di s\\
\ge \int_{r}^{+\infty} \frac{1}{(4\pi \tau)^{\sfrac n2}}\,e^{-\sfrac{s^2}{4\tau}}\frac{s}{2\tau} \theta \omega_n s^n \di s = \theta \int_{\setR^n\setminus \setB_r^n}\frac{1}{(4\pi \tau)^{\sfrac n2}}e^{-\sfrac{\|\xi\|^2}{4\tau}}\di \xi,
\end{multline*}
hence
\begin{multline}\label{eq:Colding222}
\int_{B_{\delta r}(y)}p(x,z,\tau)\di \vol(z)\\
\le (1-\theta)\int_{\setR^n\setminus \setB_r^n}\frac{1}{(4\pi \tau)^{\sfrac n2}}\,e^{-\sfrac{\|\xi\|^2}{4\tau}}\di \xi +\vol(B_{\delta r}(y)) \frac{1}{(4\pi \tau)^{\sfrac n2}}\,e^{-\sfrac{(\dist(x,y)-\delta r)^2}{4\tau}}.
\end{multline}
As pointed out in \cite[Form.\,(2.6)]{LiTamWang}, direct computations show that there exists a constant $C=C(n)>0$ such that
\[
\int_{\setR^n\setminus \setB_r^n}\frac{1}{(4\pi \tau)^{\sfrac n2}}\,e^{-\sfrac{\|\xi\|^2}{4\tau}}\di \xi \le C\left(1+\bigl(\sfrac{r}{\sqrt{4\pi\tau}}\bigr)^n\right)e^{-\sfrac{r^2}{4\tau}}.
\]
This together with \eqref{eq:Colding222} and \eqref{eq:Colding111} yields to
\begin{multline*}
p(x,y,t)\le (1-\uptheta) C\left(1+\bigl(\sfrac{r}{\sqrt{4\pi\tau}}\bigr)^n\right)e^{-\sfrac{r^2}{4\tau}}e^{\sfrac{\delta \dist^2(x,y)}{4(1+\delta)^2t}} \left(\frac{\tau}{t}\right)^{\sfrac n2}\frac{1}{\vol (B_{\delta r}(y))}\\
+\frac{1}{(4\pi t)^{\sfrac n2}}\,e^{-\sfrac{(\dist(x,y)-\delta r)^2}{4\tau}}e^{\sfrac{\delta \dist^2(x,y)}{4(1+\delta)^2t}}.
\end{multline*}
It is easily checked that $
-\sfrac{r^2}{4\tau} =-\sfrac{(\dist(x,y)-\delta r)^2}{4\tau} = - \sfrac{\dist^2(x,y)}{4(1+\delta)^3 t}$
hence
\begin{align*}
p(&x,y,t)\\
&\le\biggl[(1-\uptheta) C\left(1+\bigl(\sfrac{r}{\sqrt{4\pi\tau}} \bigr)^n \right)\left(1+\delta\right)^{\sfrac n2}\underbrace{\frac{1}{\vol (B_{\delta r}(y))}}_{\le 1/(\uptheta \omega_n \delta^n r^n)}+\frac{1}{(4\pi t)^{\sfrac n2}}\biggr]\\[-10pt]
&\hspace*{8.5cm}\cdot e^{-\sfrac{(1-\delta-\delta^2)\dist^2(x,y)}{4(1+\delta)^3t}}\\
& \le\left[(\uptheta^{-1}-1) C\left(1+\bigl(\sfrac{r}{\sqrt{4\pi\tau}} \bigr)^n\right)\frac{(4\pi \tau)^{\sfrac n2}}{(4\pi t)^{\sfrac n2}}\frac{1}{\omega_n \delta^n r^n}+\frac{1}{(4\pi t)^{\sfrac n2}}\right]\\
&\hspace*{8.5cm}\cdot e^{-\sfrac{(1-\delta-\delta^2)\dist^2(x,y)}{4(1+\delta)^3t}}\\
& = \left[(\uptheta^{-1}-1)\left(1+\bigl(\sfrac{r}{\sqrt{4\pi\tau}} \bigr)^n\right) \frac{C}{\omega_n}\frac{(4\pi \tau)^{\sfrac n2}}{\delta^n r^n} +1\right]\frac{1}{(4\pi t)^{\sfrac n2}}\\
&\hspace*{8.5cm}\cdot e^{-\sfrac{(1-\delta-\delta^2)\dist^2(x,y)}{4(1+\delta)^3t}}.
\end{align*}
Now we distinguish two cases. According to \cite[Form.\,(2.4)]{LiTamWang}, if $\dist(x,y)\le \delta \sqrt{t}$, then
\begin{equation}\label{esti1}p(x,y,t)\le \frac{1}{\uptheta}\, \frac{1}{(4\pi t)^{\sfrac n2}}\, e^{-\sfrac{d^2(x,y)}{4t}} e^{\sfrac{\delta^2}{4}}.\end{equation}
If $\dist(x,y)\ge \delta \sqrt{t}$, then $\sfrac{r}{\sqrt{\tau}}\ge \sfrac{\delta}{(1+\delta)^{2}}$, thus
\[
\left(1+\bigl(\sfrac{r}{\sqrt{4\pi\tau}}\bigr)^n\right) \frac{(4\pi \tau)^{\sfrac n2}}{\delta^n r^n} = \Bigl(\frac{(4\pi \tau)^{\sfrac n2}}{r^n} + 1 \Bigr)\delta^{-n} \le (4\pi)^{\sfrac n2}\Bigl(\frac{\delta+ 1}{\delta}\Bigr)^{2n}+\Bigl(\frac{1}{\delta}\Bigr)^n.
\]
Therefore, if $\delta<1/2$, we get
\[
\left(1+\bigl(\sfrac{r}{\sqrt{4\pi\tau}}\bigr)^n\right) \frac{(4\pi \tau)^{\sfrac n2}}{\delta^n r^n} \le C'\delta^{-2n},
\]
where $C'$ depends only on $n$, which yields to
\[
p(x,y,t) \le \left[(\uptheta^{-1}-1)\Lambda \delta^{-2n}+1\right]\frac{1}{(4\pi t)^{\sfrac n2}}\,e^{-\sfrac{(1-\delta-\delta^2)\dist^2(x,y)}{4(1+\delta)^3t}},
\]
where $\Lambda:=CC'/\omega_n$ depends only on $n$. Now we choose \begin{equation}\label{eq:delta}
\delta=\delta(\uptheta):=\min\left\{ \frac 12\,, \left(\left(\uptheta^{-1}-1\right)\Lambda\right)^{\spfrac{1}{2n+1}}\right\},\end{equation}
so that when $\left(\left(\uptheta^{-1}-1\right)\Lambda\right)^{\spfrac{1}{2n+1}} < 1/2$ then
\[
\delta(\uptheta)= \left(\uptheta^{-1}-1\right)\Lambda\delta(\uptheta)^{-2n},
\]
hence
\begin{equation}\label{eq:end}
p(x,y,t)\le \frac{\delta(\uptheta)+1}{(4\pi t)^{\sfrac n2}}\,e^{-\sfrac{(1-\delta(\uptheta)-\delta(\uptheta)^2)\dist^2(x,y)}{4(1+\delta(\uptheta))^3t}},
\end{equation}
and when $\left(\left(\uptheta^{-1}-1\right)\Lambda\right)^{\spfrac{1}{2n+1}} \ge 1/2$ -- which corresponds to the case $\uptheta \le 1-\eps_n$ with $\eps_n:=(1+2^{2n+1}\Lambda)^{-1}$ depending only on $n$ -- then $\delta(\theta)=1/2$ implies
\begin{equation}\label{eq:end2}
p(x,y,t)\le \frac{(\uptheta^{-1} - 1)\Lambda 2^{2n} + 1}{(4\pi t)^{\sfrac n2}}\,e^{-\sfrac{(1-\delta(\uptheta)-\delta(\uptheta)^2)\dist^2(x,y)}{4(1+\delta(\uptheta))^3t}}.\end{equation}
Note that $\delta(\uptheta) \to 0$ when $\uptheta \to 1$. Therefore, setting
\begin{align*}
F(\uptheta)&:=(1+\delta(\uptheta))^3/(1-\delta(\uptheta)-\delta(\uptheta)^2)\\
\tag*{and}
\gamma(\uptheta)&:=
\begin{cases}
\max(1+\delta(\uptheta),F(\uptheta)) & \text{if $1-\eps_n<\uptheta<1$,}\\
\max(2^{2n}\Lambda (\uptheta^{-1} - 1)+1,F(\uptheta)) & \text{if $0<\uptheta \le 1-\eps_n$,}
\end{cases}
\end{align*}
we get the result.
\end{proof}

For completeness, let us provide a short proof of Theorem \ref{th:Colding}.

\begin{proof}
Take $\upepsilon>0$. By Theorem \ref{th:almostrigidity}, there exists $\delta'=\delta'(n,\upepsilon)>0$ such that if $(M^n,g)$ is complete and satisfying $\Ric \ge 0$ and \eqref{eq:almostheatkernel} with $\delta$ replaced by $\delta'$, then any ball with radius $r$ in $M$ is $(\upepsilon r)$-GH close from a ball with same radius in $\setR^n$. But Theorem \ref{th:estimate} implies that there exists $\delta=\delta(n,\delta')=\delta(n,\upepsilon)>0$ such that if $1-\delta\le \theta$ holds, then $\gamma(\theta)-1 \le \delta$ and thus \eqref{eq:almostheatkernel} is true. The result follows.
\end{proof}

\section{Case of a spherical heat kernel}\label{sec:7}

In this section, for any Riemannian manifold $(M^n,g)$, we define the operator $L$ acting on $L^2(M)$ as the Friedrich extension of the operator $\tilde{L}$ defined by the formula:\vspace*{-3pt}
\[
- \int_M (\tilde{L}u)v = \int_M \langle \nabla u, \nabla v \rangle \qquad \forall u,v \in C^\infty_\rmc(M).
\]
The spectral theorem implies that $L$ generates a semi-group $(e^{tL})_{t>0}$ which admits a smooth heat kernel.

The heat kernel of $(e^{tL})_{t>0}$ on the sphere $\mathbb{S}^n$ equipped with the canonical spherical metric $g_{\mathbb{S}^n}$ admits a well-known expression, namely\vspace*{-3pt}
\[
K_t^{(n)}(\dist_{\mathbb{S}^n}(x,y))
\]
for any $x,y \in \mathbb{S}^n$ and $t>0$, where $\dist_{\mathbb{S}^n}$ is the Riemannian distance canonically associated with $g_{\mathbb{S}^n}$ and\vspace*{-5pt}
\begin{equation}\label{eq:K_t}
K_t^{(n)}(r):=\sum_{i=0}^{+\infty} e^{\lambda_i t} C_i^{(n)}(r)\end{equation}
for any $r>0$, with
\[
\lambda_i:= -i(i+n-1)\quad\text{and}\quad
C_i^{(n)}(\cdot):=(2i+n-1)(n-1)^{-1}\sigma_n^{-1} G_i^{\psfrac{n-1}{2}}(\cos(\cdot))
\]
for any $i \in \setN$. Here the functions $G_i^\alpha$ are the Gegenbauer polynomials (see \eg \cite{AtkinsonHan}). For our purposes, it is worth mentioning that
\[
C_0^{(n)}(r)=\frac{1}{\sigma_n} \quad \text{and} \quad C_1^{(n)}(r)=\frac{n+1}{\sigma_n}\cos(r)
\]
for any $r>0$. Moreover, the sum in \eqref{eq:K_t} converges uniformly in $C([0,+\infty))$.

\begin{theorem}
Let $(X,\dist,\mu)$ be a complete metric measure space equipped with a Dirichlet form $\cE$ admitting a spherical heat kernel $p$, that is
\begin{equation}\label{eq:spher}
p(x,y,t) = K_t^{(n)}(\dist(x,y))
\end{equation}
for any $x, y \in X$ and $t>0$. Then $(X,\dist)$ is isometric to $(\mathbb{S}^n,\dist_{\mathbb{S}^n})$.
\end{theorem}

\begin{proof}
Let $L$ be the self-adjoint operator canonically associated with $\cE$ and $(P_t)_{t>0}$ the associated semi-group. Assumption \eqref{eq:spher} implies that for any $t>0$ and $f \in L^2(X,\mu)$,
\[
\int_{-\infty}^{+\infty} e^{t \lambda} \di (f,E_\lambda f) = \sum_{i=0}^{+\infty} e^{\lambda_i t} \iint_{X\times X} C_i(\dist(x,y)) f(x) f(y) \di \mu(x) \di \mu(y)
\]
holds, where $(f,E_\lambda f)$ is the projection-valued measure of $L$ associated with $f$, see \eg \cite[p.\,262--263]{ReedSimon}. Uniqueness of the map $f \mto (f,E_\lambda f)$ implies that the spectrum of $L$ is given by $\lambda_0, \lambda_1, \lambda_2, \ldots$ and that the projection operators $P_i: L^2(X,\mu) \to E_i:=\Ker(L-\lambda_i \mathrm{Id})$, for any $i \in \setN$, have a kernel $p_i$ such that for any $x,y \in X$,
\[
p_i(x,y)=C_i(\dist(x,y)).
\]
Since $P_i$ commutes with $L$ for any $i$, we have $P_i Lg = \lambda_i P_i g$ for any $g \in \cD(L)$, thus $
\langle p_i(x,\cdot), Lg \rangle_{L^2} = \lambda_i \langle p_i(x,\cdot), g \rangle_{L^2}$
for any $x \in X$. This implies $p_i(x,\cdot) \in \cD(L)$ with
\[
Lp_i(x,\cdot)=\lambda_i p_i(x,\cdot)
\]
for any $x \in X$. In case $i=0$, as $\lambda_0=0$ and $p_0(x,y)=C_0(\dist(x,y))=1/\sigma_n$ for any $x, y \in X$, we get $L\textbf{1}=0$ thus $P_0 \textbf{1}=\textbf{1}$. This implies $\int_Xp_0(x,y) \di \mu(y)=1$ for any $x \in X$, hence
\begin{equation}\label{eq:volsphere}
\mu(X) = \sigma_n.
\end{equation}
In case $i=1$, we have $\lambda_1=-n$ and
\[
p_1(x,y)=C_1(\dist(x,y))=\frac{n+1}{\sigma_n} \cos(\dist(x,y))
\]
for any $x,y \in X$, hence
\begin{equation}\label{eq:cosdist}
L_x\cos(\dist(x,y)) = - n \cos(\dist(x,y)).
\end{equation}
Let $\phi_1,\ldots, \phi_l$ be continuous functions forming a $L^2(X,\mu)$-orthogonal basis of $E_1$. Observe that
\begin{equation}\label{eq:p1}
P_1f(x) = \int_X p_1(x,y) f(y) \di \mu(y) = \int_X\frac{n+1}{\sigma_n} \cos(\dist(x,y)) f(y) \di \mu(y)
\end{equation}
and\Changel
\[
P_1f(x)=\sum_{j=1}^l \left(\int_X \phi_i(y) f(y) \di \mu(y)\right)\phi_i(x)= \int_X \biggl[\sum_{j=1}^l \phi(y)\phi(x)\biggr] f(y)\di \mu(y)
\]
holds for any $f \in L^2(X,\mu)$ and $x \in X$. This implies
\[
\sum_{i=1}^l\phi(x)^2 = \frac{n+1}{\sigma_n}
\]
for any $x \in X$, hence integration over $X$ and \eqref{eq:volsphere} provides
\[
l=n+1.
\]
\Changelback
Setting
\[
\cV:=\Span\{\cos(\dist(x,\cdot)): x \in X\},
\]
we get $\cV \subset E_1$ thanks to \eqref{eq:cosdist}. Since $E_1$ is the image of $L^2(X,\mu)$ by $P_1$, the reverse inclusion follows from \eqref{eq:p1}, hence
\[
\cV=E_1.
\]
Acting as in Section \ref{subsec:4.1}, we can show that there exist $x_1, \ldots, x_{n+1} \in X$ such that $\{\delta_{x_1},\ldots,\delta_{x_{n+1}}\}$ is a basis of $\cV^*$ whose associated basis $\{h_1,\ldots,h_{n+1}\}$ of $\cV$ permits to write
\begin{equation}\label{eq:cos1}
\cos(\dist(x,y)) = \sum_{i,j=1}^{n+1} c_{ij} h_i(x) h_j(y)
\end{equation}
for any $x,y \in X$, where $c_{ij}:=\cos(\dist(x_i,x_j))$ for any $i,j$. Let $\beta$ be the bilinear form defined by
\[
\beta(\xi,\xi') = \sum_{i,j=1}^{n+1} c_{ij}\xi_i \xi_j'
\]
for any $\xi=(\xi_1,\dots,\xi_{n+1}), \xi'=(\xi_1',\dots,\xi_{n+1}') \in \setR^{n+1}$ and $Q$ the associated quadratic form. Set
\[
H:
\begin{aligned}
X & \to \setR^{n+1}\\
x & \mto (h_1(x),\ldots,h_{n+1}(x)).
\end{aligned}
\]
Then \eqref{eq:cos1} writes as
\begin{equation}\label{eq:cos2}
\cos(\dist(x,y)) = \beta(H(x),H(y)).
\end{equation}
Choosing $y=x$ implies $H(x) \in \Sigma:=\{\xi \in \setR^{n+1} \sep \beta(\xi,\xi) =1\}$, so $H(X)$ is a subset of~$\Sigma$. A direct computation provides:
\begin{equation}\label{eq:cos3}
Q(H(x)-H(y))=4\sin^2\left(\sfrac{\dist(x,y)}{2}\right) \qquad \forall x,y \in X,
\end{equation}
from which follows that $H$ is an injective map. Writing $
\setR^{n+1} = E_+ \oplus E_- \oplus \Ker \beta$ where $E_+, E_-$ are subspaces of $\setR^{n+1}$ where $\beta$ is positive definite and negative definite respectively, we can proceed as in Section \ref{subsec:4.3}, Step 1 (using the same notations) to get that $H_+$ is a bi-Lipschitz embedding of $(X,\dist)$ onto its image in $(E_+,q_+)$. Therefore, $\dim(E_+)$ is greater than the Hausdorff dimension of $X$.

\begin{claim}
The Hausdorff dimension of $X$ is $n$.
\end{claim}

\begin{proof}
The short-time expansion of the heat kernel on Riemannian manifolds \cite{MinakshisundaramPleijel} and the Cheeger-Yau estimate \cite{CheegerYau} implies that for some $C>0$ and $t_o>0$,
\[
\frac{1}{(4\pi t)^{n/2}}\, e^{-\sfrac{r^2}{4t}} \le K_t^{(n)}(r) \le \frac{C}{(4\pi t)^{n/2}}\, e^{-\sfrac{r^2}{5t}}
\]
holds for any $r \in (0,\pi)$ and $t\in (0,t_0)$. Therefore, proceeding as in the proof of Lemma \ref{lem:voleu}, we get existence of a positive constant $C$ such that for any $x \in X$ and any $r\in (0,\sqrt{t_0})$,
\[
C^{-1}r^n\le \mu(B_r(x))\le C r^n.
\]
Hence the claim is proved.
\end{proof}

Thus $n+1 \ge \dim(E_+)> n$, so $\dim(E_+)=n+1$. This shows that $\beta$ is positive definite, thus the distance $\dist_Q$ is well-defined. The associated length distance $\delta$ on $\Sigma$ is then given by:
\[
\dist_Q(\xi,\xi') = 2 \sin\left(\sfrac{\delta(\xi,\xi')}{2} \right) \qquad \forall \xi, \xi' \in \Sigma,
\]
so that one eventually has:
\[
\delta(H(x),H(y)) = \dist(x,y) \qquad \forall x,y \in X,
\]
\ie $H$ is an isometric embedding of $(X,\dist)$ into $\Sigma$ equipped with $\delta$.
Since
\[
\lim_{t\to 0^+} -4t \log K_t^{(n)}(r)=r^2,
\]
we get from Remark \ref{rem:geodesicplus} that
$(X,\dist$) is a geodesic space. Then $H(X)$ is a closed totally geodesic subset of $\Sigma$, meaning that minimizing geodesics joining two points in $H(X)$ are all contained in $H(X)$. We assume that there exists $p \in \Sigma \backslash H(X)$ and set $r:=\delta(p,H(X))$.

\begin{claim}
We have $r< \pi/2$.
\end{claim}

\begin{proof}
Assume $r\ge \pi/2$. Then $H(X)$ is contained in the hemisphere
\[
\{\sigma \in \Sigma: \beta(\sigma,p)\le 0\}.
\]
Set $\lambda(\xi)=\beta(\xi,p)$ for any $\xi \in \setR^{n+1}$. Then $\lambda \circ H:X\to\setR$ is non-positive, and $\lambda \circ H(x)=0$ if and only if $H(x)=0$, which is impossible, so $\lambda \circ H$ is actually negative. But $\lambda \circ H$ is a linear combination of $h_1,\ldots,h_n$ thus it is an element of $\cV$. Since functions in $\cV = E_1$ are $L^2$-orthogonal to constant functions, we reach a contradiction, namely $\int_X \lambda \circ H \di \mu = 0$.
\end{proof}

In fact, the same reasoning can be used to prove that $H(X)$ is contained in no hemisphere of $\Sigma$.

We are now in a position to conclude. Since $H(X)$ is closed there exists $q\in H(X)$ such that $\delta(p,q)=r$. The convexity of $H(X)$ implies that any minimizing geodesic of length $<\pi$ starting at $q$ and passing through the open ball $B_r(p)$ cannot meet $H(\Sigma)$. But the union of these minimizing geodesics is an open hemisphere, so $H(X)$ is contained in the complementary hemisphere, hence a contradiction.
\end{proof}

\backmatter
\bibliographystyle{jepalpha+eid}
\bibliography{carron-tewodrose}
\end{document}
